% !TeX program = pdflatex
%==============================================================================
% A Chronology of the HELM Reproducibility Internship (Kitware, Summer 2026)
%
% Author:  Edward Wang <edward.wang@kitware.com>
% Purpose: An authoritative, detailed technical chronology of the internship
%          work, written as a series of technical reports / blog posts, to
%          serve as source material for an academic paper on benchmark
%          reproducibility.
%
% This document is reconstructed from: the weekly intern status decks
% (docs/internship-powerpoints.zip), the append-only developer journals
% (dev/journals/{claude,codex}.md and their 2026-H1 archives), the git history
% (452 authored commits, 2026-06-03 .. 2026-07-14; a further 127 commits
% 2026-07-16 .. 2026-07-27 recorded in the 2026-07-27 revision), the standing
% research journal (docs/helm-reproduction-research-journal.md), and the
% technical design docs under docs/ (helm-gotchas.md, vllm-vs-huggingface-
% deployment-match.md, helm-unrecorded-deployment-params.md,
% eee-vs-helm-metadata.md, pipeline.md, architecture.md,
% container-execution.md).
%
% Revision history:
%   2026-07-15   interpretive revision (epistemic status of central claims)
%   2026-07-15b  packaging fix
%   2026-07-15c  consensus revision (canonical ledgers, three reproduction
%                targets, evidence-access attestation)
%   2026-07-15d  consensus accepted in full; store-provenance audit furnished
%   2026-07-27   continuation revision: evidence cutoff moved to 2026-07-27;
%                the fp32/substrate thread closed and attributed; the
%                open-judge experiment executed; the Qwen3.5 extension.
%==============================================================================
\documentclass[11pt]{report}

\usepackage[margin=1in]{geometry}
\usepackage[T1]{fontenc}
\usepackage[utf8]{inputenc}
\usepackage{microtype}
\usepackage{xspace}
\usepackage{amsmath}
\usepackage{amssymb}
\usepackage{amsfonts}
\usepackage{booktabs}
\usepackage{tabularx}
\usepackage{longtable}
\usepackage{array}
\usepackage{float}
\usepackage{xcolor}
\usepackage{listings}
\usepackage{enumitem}
\usepackage{graphicx}
\usepackage[
  colorlinks=true,
  linkcolor=black,
  citecolor=black,
  urlcolor=blue
]{hyperref}
\usepackage{bookmark}
\usepackage{url}
\usepackage{xurl}

\newcolumntype{Y}{>{\raggedright\arraybackslash}X}
\newcolumntype{Z}{>{\raggedright\arraybackslash}p{0.30\textwidth}}

\lstdefinestyle{shell}{
  basicstyle=\ttfamily\scriptsize,
  breaklines=true,
  breakatwhitespace=false,
  columns=fullflexible,
  keepspaces=true,
  showstringspaces=false,
  frame=single,
}
\lstset{style=shell}

% Shorthands
\newcommand{\eee}{\textit{Every Eval Ever}\xspace}
\newcommand{\eeeshort}{\textsc{eee}\xspace}
\newcommand{\helm}{\textsc{helm}\xspace}
\newcommand{\code}[1]{\texttt{#1}}
\newcommand{\file}[1]{\texttt{#1}}
\newcommand{\fromspec}{\emph{from-spec}\xspace}

% A consistent "finding" callout
\newcommand{\finding}[1]{\par\noindent\textbf{Finding.}~#1\par}
\newcommand{\statusEstablished}{\textbf{Established}}
\newcommand{\statusImplemented}{\textbf{Implemented}}
\newcommand{\statusPreliminary}{\textbf{Preliminary}}
\newcommand{\statusLocalOnly}{\textbf{Locally supported; not packaged}}
\newcommand{\statusCollaboratorVerified}{\textbf{Collaborator-verified; artifact not packaged}}
\newcommand{\statusInternalEstimate}{\textbf{Internal estimate; regenerate for publication}}
\newcommand{\statusUnresolved}{\textbf{Unresolved}}
\newcommand{\statusPlanned}{\textbf{Planned}}
\newcommand{\statusResolved}{\textbf{Resolved (2026-07-23)}}
% Marks a ruling or passage that a later revision has superseded. The original
% wording is always retained; the marker records what replaced it and where.
\newcommand{\superseded}[1]{\par\smallskip\noindent\textit{\small Superseded
  2026-07-27: #1}\par}
\newcommand{\evidencebox}[2]{\begin{quote}\noindent\textbf{#1.} #2\end{quote}}


\title{\bfseries Master Collaborative Reference: Reconstructing \textsc{helm} Evaluations\\[4pt]
  \large Internship Chronology, Evidence Audit, and Publication Roadmap\\[10pt]
  \normalsize Kitware Research Internship --- Summer 2026}
\author{Edward Wang\\ \texttt{edward.wang@kitware.com}\\[4pt]
  \small Kitware, Inc.\\[6pt]
  \small\normalfont Chapters~\ref{chap:qwen35}--\ref{chap:paperdirection} record
  work led by Jon Crall (Kitware) from 2026-07-15 onward; see
  \S\ref{sec:rev-authorship}.}
\date{Evidence cutoff: 2026-07-27 (chronology through 2026-07-14 unchanged from
  revision 2026-07-15c); collaborative interpretation revised 2026-07-15;
  consensus reached 2026-07-15d; continuation revision 2026-07-27}

\begin{document}
\maketitle

\begin{abstract}
\noindent
This master reference consolidates two independently reconstructed chronologies,
the weekly internship decks, developer journals, Git history, runbooks, paper
sources, and technical design documents for Edward Wang's Kitware internship from
May 15 through July 27, 2026. The initial goal---large-scale reproduction of public
\textsc{helm} results---became a study of a deeper measurement problem: a public
benchmark recipe does not necessarily identify the complete experiment that
produced a score. The surviving record may omit model and tokenizer revisions,
load precision, chat-template behavior, serving engine and kernel details,
hardware-sensitive numerical settings, harness era, and artifact-migration history.
The internship built exact \file{run\_spec.json} replay, containerized and era-pinned
execution, GPU leasing, layered comparison reports, and deployment-matching tools to
attribute disagreement across these layers. The central OLMo float32 result, reported
in the 2026-07-15 revisions as a 12-instance discovery probe awaiting confirmation,
was carried to closure between July 21 and July 23: for OLMo-2 instruct IFEval the
published number is reproduced \emph{byte-exactly} once three unrecorded substrate
variables---load precision, chat-template rendering, and the inference engine---are
recovered together, and each accounts for a measured slice of the drift
(Chapter~\ref{chap:substrate-closed}). Recovering only the first two leaves a
$+0.07$ residual that is attributable to \textsc{vllm}$\leftrightarrow$Hugging~Face
engine numerics alone. A parallel experiment substituted open-weight judges into
the released scoring pipeline of six model-graded benchmarks, and found that judge
non-determinism is universal while metric fragility is a property of the metric
(Chapter~\ref{chap:openjudge}). The historic
classic-model investigation supplies the complementary result: some provenance is
recoverable, whereas migrated class paths and unreleased producing code can make the
original measurement non-identifiable from the surviving evidence examined. This document distinguishes inherited
pre-internship EEE work from internship contributions, preserves contemporaneous
interpretations without presenting superseded claims as final conclusions, and
provides an evidence-status ledger and a bounded publication roadmap.
\end{abstract}

\tableofcontents
%==============================================================================
\chapter{Evidence Status and Paper Lineage}
\label{chap:collaborative-validation}
%==============================================================================

This edition reconciles two independent reconstructions of the internship record.
They agree on the chronology and engineering substance. The main differences
concerned the strength of empirical claims, the distinction between a
deployment-matching probe and an authoritative benchmark reproduction, and the
evidentiary status of live result stores omitted from the review archive. The
following consensus rulings govern all later chapters; where an older passage
records a stronger contemporaneous belief, this chapter and the inline status
callout take precedence.

\section{Consensus rulings from independent reconstructions}

\begin{longtable}{@{}p{0.31\textwidth}p{0.15\textwidth}p{0.46\textwidth}@{}}
\toprule
\textbf{Point} & \textbf{Ruling} & \textbf{Reason and qualification} \\
\midrule
\endhead
The EEE paper and the internship follow-on are distinct paper lines. &
\statusEstablished & The NeurIPS case study and 2,626-line technical report predate
the internship and focus on normalization and discrepancy detection. A separate
TMLR skeleton existed on May 14. The follow-on should be positioned as attribution,
possible closure, and identifiability rather than a repetition of EEE's
Pythia/Vicuna/Falcon study. \\
The positive thesis ``most residual disagreement is small and attributable'' should
be demoted. & \statusEstablished & It is a motivating hypothesis, not a result,
until the denominator is frozen, report stores are fresh, selected configurations
are confirmed through the normal HELM path, discovery is separated from held-out
evaluation, and repeatability/ablation experiments are complete. \\
The OLMo float32 result is discovery rather than completed reproduction. &
\statusResolved & \emph{Correct when written, and the confirmation it demanded was
subsequently run.} The ranking came from a small oracle-scored IFEval probe whose
winning cells used \code{add\_special\_tokens=False}, a knob the ordinary HELM path
does not send. Between July 21 and 23 the confirm step was executed at full $n$
(1082 per-instance rows per side): float32 alone \emph{does not} recover the
official (a $+0.067$/$+0.082$ residual survives, refuting the registered
prediction), and recovering precision, template rendering, \emph{and} the engine
together reproduces the official completions byte-exactly
(Chapter~\ref{chap:substrate-closed}). The ruling stands as stated; its required
action is discharged for one cell (OLMo-2 7B/13B instruct, IFEval) and open
elsewhere. \\
The dtype argument has a deductive source-level leg. & \statusPreliminary & Public
OLMo instruct deployments are recorded as \code{HuggingFaceClient} with no pinned
\code{torch\_dtype}; the client did not inject one, and Transformers 4.x defaulted
such loads to fp32. The January 2025 OLMoE timing and architecture bound the likely
version to the 4.x regime. Because the exact historical environment is absent, the
correct conclusion is ``most likely fp32 under the bounded loader assumptions,''
not a uniquely proven full deployment. \\
Identifiability is per parameter rather than all-or-nothing. & \statusEstablished &
Client class and configured dtype presence are source-readable; tokenizer behavior
and templates may be repository-readable; exact package builds, kernels, batch
composition, hardware, and hosted-provider internals may remain unidentified. A
per-parameter identifiability map is therefore a useful paper contribution. \\
The dense OLMo-2 Hugging Face divergence is scientifically relevant. &
\statusResolved & The ruling was right to keep a dense-model probe artifact on the
candidate list rather than promoting execution-stack sensitivity to a finding: that
is what it turned out to be. Re-running the probe with every axis pinned to the
configuration the officials actually used (\code{float32}, old-template rendering,
\textsc{helm}'s decode) reproduces the official dense OLMo-2 completions
\emph{exactly} in all four forward-pass cells --- \{eager,\,\textsc{sdpa}\} $\times$
\{\code{device\_map=auto},\,single\} --- on 96~GB cards, where nothing shards. The
earlier near-zero quasi-match is attributed to the probe's own unpinned axes
(template rendering swept to the modern behavior; \code{device\_map=auto} sharding a
fits-on-one-card model across the smaller cards of that host, which changes fp32
reduction order), not to a current-HF-versus-historical-HF divergence
(\S\ref{sec:residual-resolved}). The execution-stack sensitivity claim survives, but
its evidence is now the \textsc{vllm}$\leftrightarrow$HF engine gap, which is
measured rather than inferred. \\
The six-category disagreement taxonomy is preferable to ``true reproducibility
failure.'' & \statusEstablished & Non-runnable cases are an upstream eligibility
gate. Observed disagreements should be classified as recipe, deployment,
execution-instrument, or artifact-interpretation mismatch; residual disagreement
under controlled equivalence; or historically non-identifiable reproduction. \\
Recoverable versus non-identifiable provenance is the strongest conceptual spine. &
\statusEstablished & OLMo supplies mechanisms whose missing parameters can be
partly reconstructed. GPT-J/GPT-NeoX/OPT artifacts point to an unreleased producing
instrument and cannot be uniquely mapped to a released harness. OLMo should be
called \emph{partially recoverable} until ordinary-path confirmation is complete. \\
RedPajama-3B ran end-to-end under both proxy eras. & \statusCollaboratorVerified & July 14
journal entries refer to intact era-RedPajama runs and refreshable report artifacts,
and the slide deck states that the E2E script works for v0.2.4 and v0.3.0. The
provided source archive explicitly excludes \path{/data}, so the reports themselves
are not in this evidence package. Preserve and hash them before paper submission.
Its ``recovery'' is a forensic proxy comparison against a public-side serving
artifact, not proof that either later era reproduces the original producing
instrument. \\
Preserving the dated narrative while appending corrections is acceptable. &
\statusImplemented & It preserves research history, but an authoritative reference
must prevent readers from mistaking superseded wording for the final conclusion.
This master edition therefore uses corrected wording inline and retains historical
interpretations only as explicitly labeled contemporaneous beliefs. \\
The claimed commit \code{86ec84af} contains the revisions. & \statusCollaboratorVerified &
A collaborating live-repository review reports that \code{86ec84af} was the current
\textsc{head} of \code{impl/run-from-run-spec} on 2026-07-15 and contained the stated
revisions. The supplied review archive ends at \code{b8b8586}; the later Git object or
bundle is not included in this package. Treat the fact as collaborator-verified until
a Git bundle or immutable repository reference is preserved with the paper artifact.
\emph{Update 2026-07-15d:} the full commit SHA
\code{86ec84af09c0b8f4442afa1ce87b56e1f8b3dc61} and tree
\code{a3734de5b0178077bf241b6026e7d68d74247853} were furnished under
\file{validation/consensus\_accepted\_2026-07-15d/}. A Git bundle was deliberately
\emph{not} exported (the repository is private); the row is therefore
\emph{reference-furnished, bundle pending release}. \\
\bottomrule
\end{longtable}

\section{Evidence precedence}

For this master reference, evidence is prioritized as follows: immutable public or
local artifacts with hashes; source code and Git history; generated reports whose
inputs remain available; journals and runbooks; slide-deck summaries; and finally
retrospective narrative interpretation. A report-store claim whose raw directory is
excluded from the archive is not silently promoted to a packaged result.

\section{Evidence-access asymmetry and attestation}

The collaborating reviewers did not have identical evidence access. One review used
the supplied \code{git archive}, which excludes \path{/data}, ignored files, and
post-archive commits. A later review reports direct access to the live branch and
those stores. This resolves several factual questions, but it does not replace an
archival artifact. Claims based only on that access are marked
\statusCollaboratorVerified{} until the relevant Git objects, manifests, report
inputs, and hashes are copied into the release bundle.

For result stores, three independent freshness questions must be recorded:
\begin{enumerate}[nosep]
  \item \textbf{run-artifact freshness}: were the model outputs generated by the
        intended code, model revision, and execution path?;
  \item \textbf{comparison freshness}: were those outputs compared with the intended
        public target using the corrected matching and deduplication logic?; and
  \item \textbf{report freshness}: were tables and plots regenerated after the latest
        reporting fixes?
\end{enumerate}
A recent report timestamp establishes only the third item unless the first two are
also demonstrated by a manifest.

\section{Canonical claim ledger}

The machine-readable file
\file{chronicle/data/master\_claim\_evidence\_ledger\_2026-07-27.csv} is the
canonical claim-status source for this reference set; it supersedes the
\code{2026-07-15c} ledger, which is retained beside it for the audit trail. Prose
tables are readable snapshots and should be regenerated from, or checked against,
that ledger before any submission. This prevents the chronology, strategy, and review
appendices from drifting into incompatible statuses. The companion
\file{chronicle/data/store\_status\_ledger\_2026-07-27.csv} carries the three
freshness dimensions per store, with the observed \code{run\_artifact\_freshness}
values merged in from the 2026-07-15d disk-level audit.

\section{Two-paper lineage}

The inherited EEE paper demonstrates that a normalized result representation can
expose reproducibility gaps, but explicitly cannot separate checkpoint,
quantization, precision, tokenizer, and generation-kernel effects. The follow-on
paper should use the positioning sentence:

\begin{quote}\itshape
EEE detects reproducibility gaps; this work studies their attribution, possible
closure, and identifiability.
\end{quote}

\section{Three orthogonal reproduction targets}
\label{sec:targets-3}

The six-category mismatch taxonomy answers \emph{why} two executions disagree. A
separate three-level axis answers \emph{what} is being reproduced:
\begin{enumerate}[nosep]
  \item \textbf{Artifact reconstruction}: recompute the published aggregates from the
        released prompts, completions, and metric records without rerunning a model.
  \item \textbf{Procedural reproduction}: execute the recorded recipe under a pinned
        or explicitly equivalent instrument and compare instance-level outputs and
        metrics.
  \item \textbf{Claim-level robustness}: determine whether the substantive conclusion
        (for example, a ranking or model-generation improvement) survives plausible
        execution substrates even when byte-identical outputs do not.
\end{enumerate}
These levels should not be collapsed. Released artifacts are the historical ground
truth for what HELM reported; reruns assess the stability and identifiability of the
procedure that produced them.

\section{Corrections applied in this master revision (2026-07-15c)}
\label{sec:corrections-applied}

Following a collaborating review of the \emph{live} repository (reported HEAD
\code{86ec84af} on \code{impl/run-from-run-spec}, including \path{/data} stores that
the source archive omits), the following corrections were applied \emph{inline}.
Because the Git objects and raw stores are not themselves included, live-only facts
are recorded as collaborator-verified rather than promoted to immutable packaged
evidence. The reasoning and machine-readable attestations are preserved under
\file{validation/claude\_patch/}.

\begin{itemize}[nosep]
  \item \textbf{Commit \code{86ec84af}} reclassified from unresolved to
        \statusCollaboratorVerified{} --- reported as the live branch head, pending
        inclusion of the Git object or an immutable remote reference.
  \item \textbf{Store status split three ways}, replacing the single ``not
        packaged'' label: the GPT-OSS \emph{reports} are reported fresh but still need
        run/comparison acceptance plus preservation; \code{olmo-models-combined} is
        reported stale (on-disk \code{olmo-7b} halving \code{0.295/0.144}) and must be
        regenerated; and the Qwen store is reported genuinely absent
        (Chapter~\ref{chap:newmodels}, Appendix~\ref{app:revision}).
  \item \textbf{OLMo-2 HF divergence} now keeps a dense-model probe artifact on the
        live candidate-cause list, alongside genuine execution-instrument sensitivity
        (Chapter~\ref{chap:deploymatch}).
  \item \textbf{Wording and sourcing:} ``apparently HF-produced'' $\rightarrow$
        ``recorded as a \code{HuggingFaceClient} deployment''; the GPT-OSS-vs-OLMo
        \code{ifeval} drift ratio ``two orders of magnitude'' $\rightarrow$
        ``$\approx\!30\times$''; the 59\% classic-track figure now carries its source.
\end{itemize}

\noindent The structural concerns from that review are also resolved in this
consensus package: the claim ledger is canonicalized, vendor-specific collaboration
meta is neutralized, and the full directory layout includes the referenced appendices
and an accurate README.

\section{Consensus closure (2026-07-15d)}
\label{sec:consensus-closed}

The 2026-07-15c package was accepted in full by the collaborating reviewer on the
same day, with no rejections and no open disagreements. The acceptance letter,
furnished source hashes, and a disk-level store-provenance audit are preserved under
\file{validation/consensus\_accepted\_2026-07-15d/}. Two of its contents changed
this reference set rather than merely endorsing it:

\begin{itemize}[nosep]
  \item \textbf{Furnished immutable references.} Commit and tree SHAs for
        \code{86ec84af}, plus \code{sha256} values for the GPT-OSS, OLMo, and
        both-era RedPajama provenance and manifest files. A Git bundle was
        deliberately withheld (private repository); the SHAs are the immutable
        reference until a Kitware-controlled artifact release exists.
  \item \textbf{A run-artifact acceptance audit that corrected two next-actions.}
        Direct \path{/data} inspection found that the two flagship modern stores
        retain \emph{no} raw run artifacts --- \code{gpt-oss-20b-from-spec} and
        \code{olmo-models-combined} were pruned to their derived layers only --- so
        ``regenerate'' means \emph{re-run}, not refresh a report.
        \code{era-redpajama} (both eras) and the deployment-match sweeps are the
        only stores preservable by copy-and-hash as they stand. The
        non-obvious inversion worth carrying into the paper's artifact statement:
        \emph{the freshest reports sit on the least-preserved inputs}, which is
        precisely why a report timestamp cannot be an acceptance criterion. These
        observations are merged into
        \file{chronicle/data/store\_status\_ledger\_2026-07-27.csv}.
\end{itemize}

The consensus thesis statement, agreed verbatim by both reviewers, is the sentence
this document is organized around:

\begin{quote}\itshape
A public \textsc{llm} benchmark recipe is not a complete experimental record.
Reproduction is a layered system-identification problem spanning recipe, deployment,
execution instrument, and artifact interpretation; controlled reconstruction can
attribute and sometimes close apparent gaps, while some historical measurements
remain non-identifiable from the surviving evidence examined.
\end{quote}

\section{What changed in the continuation revision (2026-07-27)}
\label{sec:rev-2026-07-27}

The 2026-07-15 revisions closed on a consensus and an operational to-do list. This
revision records what happened to that list over the following twelve days
(127 commits, 2026-07-16 through 2026-07-27) and re-scopes the affected claims. The
chronology through July 14 is unchanged; superseded rulings are marked in place
rather than rewritten.

\paragraph{Three substantive additions.}
\begin{enumerate}[nosep]
  \item \textbf{The fp32/substrate thread is closed and attributed}
        (Chapter~\ref{chap:substrate-closed}). The confirm step the consensus
        demanded was run end-to-end. It \emph{refuted} the registered prediction
        that float32 alone would recover the official, and then --- by reproducing
        the official on the original engine --- resolved the residual to
        \textsc{vllm}$\leftrightarrow$Hugging~Face fp32 engine numerics. For
        OLMo-2 instruct IFEval the published number is now recoverable
        byte-exactly, and each unrecorded substrate layer accounts for a measured
        slice of the drift. This upgrades OLMO-001 from \emph{preliminary} to
        \emph{established for one cell} and closes OLMO-003.
  \item \textbf{The open-judge extension became an executed experiment}
        (Chapter~\ref{chap:openjudge}). What \S\ref{sec:judge} recorded as a design
        (a judge registry and a scoped \code{same\_judge} fact) was built out into
        a rejudging pipeline behind a machine-precision identity-replay gate,
        and run: six benchmarks, six judge arms, three replicates. It contributes
        a finding the reproduction thread cannot supply on its own --- that judge
        \emph{non-determinism} is universal at temperature~0 while metric
        \emph{fragility} depends on the metric's granularity --- and a
        contamination caveat that bounds how the agreement numbers may be read.
  \item \textbf{A net-new-model extension path was proven end-to-end}
        (Chapter~\ref{chap:qwen35}). Registering a model \textsc{helm} has never
        seen now costs one preset block, two sidecar YAMLs, and one catalog
        endpoint, with no \textsc{helm} edit and no image rebuild. The exercise
        also produced a self-critique with teeth (\S\ref{sec:freeze-critique}):
        our own freshly computed runs still store a mutable run-key string as
        their source of truth, which is the practice this project criticizes
        \textsc{helm} for.
\end{enumerate}

\paragraph{Status changes applied to earlier chapters.}
\begin{itemize}[nosep]
  \item The consensus ruling on the float32 result is marked \statusResolved{} ---
        \emph{correct as written, and its required action discharged for one cell}.
  \item The dense OLMo-2 Hugging Face divergence
        (\S\ref{sec:residual-resolved}) is marked \statusResolved{}: the dense-model
        probe artifact the 15c review insisted on keeping as a live candidate is
        what it was. The 15c wording is retained above the resolution.
  \item The unrecorded-parameter taxonomy (Appendix~\ref{app:unrecorded}) gains a
        measured effect size for the inference engine, which was previously
        argued from kernel behavior rather than from a metric.
  \item Two gotchas are added (G14, G15; Appendix~\ref{app:gotchas}).
  \item The open-threads list (\S\ref{sec:open-threads}) is rewritten against the
        state at this revision rather than at July 14.
\end{itemize}

\paragraph{What did \emph{not} change.} The preservation position is unchanged and,
if anything, more urgent: the results in Chapter~\ref{chap:substrate-closed} were
produced after the 15d audit found the flagship stores pruned, so they are
\emph{new} raw artifacts that have not yet been copied, hashed, or bundled. Nothing
in this revision promotes a live-store observation to packaged evidence. The
corpus-wide denominator is still an internal estimate awaiting regeneration from a
checked-in manifest, and the Qwen reproduction store is still absent.

\section{Authorship and division of work from 2026-07-15}
\label{sec:rev-authorship}

Through July 14 this is a record of one intern's work. From July 15 a second
Kitware engineer, Jon Crall, joined and led most of the execution recorded in
Chapters~\ref{chap:qwen35}--\ref{chap:paperdirection}: of the 127 commits in this
revision's window, 107 are his and 20 are Edward Wang's (the latter being the TMLR
manuscript work of \S\ref{sec:paper-drafting}). The substrate-recovery experiment
was designed against Edward's earlier probe results and runbooks and executed on
Kitware GPU hosts; the open-judge and Qwen3.5 threads are new work started in this
window. This matters for the manuscript in two ways: the contribution statement must
reflect it, and the internship-timeline framing of Appendix~\ref{app:deliverables}
no longer covers the whole document. Development throughout continued to use an
AI pair-programming harness, recorded per session in the developer journals.


%==============================================================================
\chapter{How to Read This Document}
%==============================================================================

\section{Scope and provenance}

This is an \emph{authoritative reference} assembled after the fact from primary
sources so that a future paper can cite it with confidence. The reconstruction
draws on:

\begin{itemize}[nosep]
  \item the seven weekly intern status decks (2026-06-02 through 2026-07-14),
        which record the work in the author's own words and contain the headline
        result figures;
  \item the append-only developer journals (\file{dev/journals/claude.md} and its
        \file{2026-H1} archive), which capture, per work session, the intent, the
        design alternatives considered, the chosen approach, and the reusable
        lesson --- roughly one hundred design narratives;
  \item the git history: 452 commits authored between 2026-06-03 and 2026-07-14,
        whose messages give a fine-grained, dependency-ordered record of what
        landed and when;
  \item the standing research journal
        (\file{docs/helm-reproduction-research-journal.md}), which frames the
        scientific question and records the early foundational findings; and
  \item the technical design documents accumulated under \file{docs/}, chiefly
        \file{helm-gotchas.md}, \file{vllm-vs-huggingface-deployment-match.md},
        \file{helm-unrecorded-deployment-params.md}, \file{eee-vs-helm-metadata.md},
        \file{pipeline.md}, \file{architecture.md}, and
        \file{container-execution.md}.
\end{itemize}

The internship ran from 2026-05-15. The first $\sim$2.5 weeks (through the 06-02
status update) were spent on onboarding and a literature review; the first
authored commits appear on 2026-06-03. The work then proceeds essentially without
interruption through 2026-07-14, the last date this chronology covers.

\section{A note on inherited foundations versus internship contributions}

The internship did not start from an empty repository. A substantial foundation
was already in place: the \file{eval\_audit} package itself; the \emph{Every Eval
Ever} (\eeeshort) normalized artifact format used as the common comparison
substrate; the \code{kwdagger} scheduler integration; and an early set of
reproducibility findings on \textsc{boolq}/Pythia and Vicuna/\textsc{NarrativeQA}
(described in \S\ref{sec:inherited}). Where this document describes that
foundation, it does so to make the internship's own contributions legible; the
chronology proper (Chapters~\ref{chap:e2e} onward) is the author's work. The
developer journals are written in the first person by the AI pair-programming
harness (Claude Code) that the author drove; throughout this report, ``we'' denotes
the author directing that toolchain, making the design decisions, and executing and
interpreting the runs on the GPU hosts.

\section{The two threads}

Two research threads ran in parallel for the first month, after which the second
was paused in favor of the first:

\begin{description}[style=nextline,leftmargin=1.5em]
  \item[Thread A --- \textsc{helm} reproduction (Chapters \ref{chap:framing}--\ref{chap:reporting}).]
    The main line of work: building the machinery to faithfully replay public
    \textsc{helm} runs on local open-weight hardware, and using it to determine
    what is and is not reproducible, and why. This is the bulk of the internship
    and of this document.
  \item[Thread B --- efficient benchmarking (Chapter \ref{chap:effbench}).]
    An exploratory statistics thread: can a small, well-chosen \emph{coreset} of
    benchmark questions plus a regression predict a model's full-benchmark score,
    and does a Data-Kernel Perspective-Space (\textsc{dkps}) representation help?
    This lived in two sibling repositories (\file{mrmr\_eval}, \file{dkps}) and
    reached a working prototype with preliminary results before being paused.
\end{description}

%==============================================================================
\chapter{The Research Question and Its Framing}
\label{chap:framing}
%==============================================================================

\section{From ``reproduce \textsc{helm}'' to a falsifiable claim}

The stated initial goal was to perform large-scale reproductions of the public
\textsc{helm} (Holistic Evaluation of Language Models) results. The public
corpus publishes, for many (scenario, model) pairs, the prompts sent, the
completions returned, and the computed metrics. Reproducing a result means
re-running the same evaluation locally and checking that the numbers match.

That goal did not survive contact with the data, and the way it failed is itself
the scientific content. Two categories of models dominate the public corpus:
open-weight models that were originally served through a hosted API (frequently
\emph{Together AI}), and closed-weight models that cannot be run locally at all.
The recipe that produced each public number is only partially recorded: the
\file{run\_spec.json} captures \emph{what} to generate (prompts, decoding
parameters, scenario, metrics) and the \emph{name} of the model and its
deployment, but almost nothing about the \emph{execution substrate} --- the
precision the weights were loaded at, the exact checkpoint revision, the attention
kernel, the tokenizer and chat-template versions, the software-stack versions. As
a result, three distinct difficulties compound:

\begin{enumerate}[nosep]
  \item \textbf{Running the benchmarks at all} --- gated models and datasets,
        packaging incompatibilities, download endpoints that have since revoked
        access, and scenario code that requires optional dependencies or
        credentials.
  \item \textbf{Setting up a \emph{faithful} local reproduction} --- matching the
        original recipe when key execution parameters were never written down, and
        when the original was served by an external provider whose stack we cannot
        inspect.
  \item \textbf{Attributing disagreement} --- when local and public numbers
        differ, deciding whether the benchmark is legitimately non-reproducible,
        whether the difference is an artifact of an uncontrollable external
        deployment, or whether our local reproduction is simply missing a detail.
\end{enumerate}

The goal was therefore reframed. The paper-worthy claim is \emph{not} ``all of
public \textsc{helm} is reproducible.'' It is:

\begin{quote}\itshape
A public benchmark recipe is not a complete experimental record. Reproduction is a
layered system-identification problem: controlled reconstruction can attribute and
sometimes close apparent gaps, while lost or transformed provenance can leave the
original measurement non-identifiable.
\end{quote}

\section{The load-bearing distinction: eligibility and disagreement layers}
\label{sec:distinction}

The first gate is whether a run is executable at all. Gated models or datasets,
revoked download access, missing credentials, unsupported packaging, unavailable
judges, and infeasible resource requirements are \emph{eligibility or environment
filters}; they are not negative reproduction measurements.

For runs that execute and can be compared, the final taxonomy is:

\begin{enumerate}[nosep]
  \item \textbf{Recipe mismatch}: scenario, prompt construction, examples, decoding,
        or metrics differ.
  \item \textbf{Deployment mismatch}: the client, hosted endpoint, serving engine,
        or exposed deployment identity differs.
  \item \textbf{Execution-instrument mismatch}: model/tokenizer revision, dtype,
        quantization, template implementation, software stack, kernels, topology,
        batching, caching, or hardware-sensitive numerical behavior differs.
  \item \textbf{Artifact-interpretation or migration mismatch}: stored artifacts
        were reordered, normalized, carried forward, or migrated under later schemas.
  \item \textbf{Residual disagreement under controlled equivalence}: material
        controllable layers are matched and disagreement remains. This is the only
        category analogous to conventional repeatability failure.
  \item \textbf{Historically non-identifiable reproduction}: surviving evidence does
        not uniquely specify what execution should count as faithful.
\end{enumerate}

The project's earlier binary---environment/recipe failure versus ``true
reproducibility failure''---was a useful operational first cut, but it is preserved
only as a contemporaneous framing. The six-layer taxonomy is the authoritative one.

\section{The inherited baseline: what was already known}
\label{sec:inherited}

Before the internship, the project had established a small but load-bearing set of
results, recorded in \file{docs/helm-reproduction-research-journal.md}. They set
the priors that the internship's work refined.

\begin{itemize}
  \item \textbf{Local reruns are stable; the official--local gap is not noise.} For
        \code{boolq:model=eleutherai/pythia-6.9b}, two local repeats agreed at a
        strict run-level ratio of $0.955$ with a maximum run-level delta of
        $0.0015$, whereas official-vs-local agreement was $0.463$ with run-level
        deltas up to $\sim\!12$. The official--local gap is far larger than
        rerun noise, so ordinary nondeterminism cannot be the main explanation ---
        the residual is structural (deployment and execution-spec drift).
  \item \textbf{A dramatic ``irreproducibility'' turned out to be a local
        misconfiguration.} Local Vicuna/\textsc{NarrativeQA} produced 99/100 empty
        completions; the cause was \textsc{helm} silently auto-applying a
        \emph{chat template} to a non-chat prompt. Rerunning with
        \code{apply\_chat\_template: false} recovered near-official scores
        (\code{exact\_match} $0.27$, \code{f1} $0.64$, \code{rouge\_l} $0.64$).
        The lesson --- that \code{apply\_chat\_template} must be an explicit,
        controlled setting, and that high empty-completion rate / near-zero output
        tokens / pervasive \code{finish\_reason\_unknown} are red flags --- recurs
        throughout the internship.
  \item \textbf{The metric-scale caveat.} A single global tolerance is hard to
        interpret because \emph{core} metrics live in $[0,1]$ while
        \emph{bookkeeping} metrics (byte counts, token counts, log-probs) do not.
        The reporting must separate metric classes; this became the
        \code{metrics\_taxonomy} module (\S\ref{sec:taxonomy}).
\end{itemize}

These priors motivated the internship's dominant methodological choice --- prefer
open-weight models runnable on realistic hardware, control the deployment
explicitly rather than trusting \textsc{helm}'s automatic inference, and treat the
reproducibility question as a \emph{distribution} (same-machine, cross-machine,
official-vs-local) rather than a single point estimate.

%==============================================================================
\chapter{Onboarding and Literature Review (mid-May -- June 2)}
\label{chap:litreview}
%==============================================================================

The first status deck (2026-06-02) reports two workstreams: end-to-end tests for
the audit pipeline (Chapter~\ref{chap:e2e}) and a literature review that seeded
both the reproducibility thread and the efficient-benchmarking thread
(Chapter~\ref{chap:effbench}). Three papers were reviewed in depth.

\begin{description}[style=nextline,leftmargin=1.25em]
  \item[LoRA weight-space learning --- \emph{W2T: LoRA Weights Already Know What
        They Can Do} (Han et al., 2026), building on Chen, Helm, Park, Priebe \&
        Villar, \emph{Extracting information from fine-tuned weights}.]
    A LoRA update is $AB^\top$; the factorization $(A,B)$ is non-identifiable
    because $(AM, BM^{-\top})$ yields the same update for any invertible $M$ (a
    $\mathrm{GL}$-group symmetry). W2T canonicalizes via the SVD and embeds the
    canonicalized weights with a transformer into a ``weight-space,'' whose
    embeddings drive downstream tasks: attribute classification (properties of the
    fine-tuning data), performance prediction, and adapter retrieval. The review
    noted that the ablations attribute most of the method's power to the
    \emph{canonicalization} step, and floated directions: generative modelling over
    LoRA weights, simpler statistical embeddings (e.g.\ MDS) in place of the
    transformer, and multi-/merged-adapter settings.
  \item[Efficient benchmarking --- \emph{Efficient Benchmarking Is Just Feature
        Selection and Multiple Regression} (Bowyer, Locatelli \& Cao, 2026).]
    Given $M$ source models, represent each question $X_i$ as a vector in
    $\mathbb{R}^M$ (each coordinate is a model's score on that question) and the
    target $Y$ as the vector of overall scores; \emph{select} an informative
    subset of questions with mRMR (Minimum-Redundancy-Maximum-Relevance) and
    \emph{predict} full-benchmark scores with ridge / kernel-ridge regression.
    Baselines: Lasso, Anchor Points, IRT. This framing became Thread~B.
  \item[Connecting the two.] The review proposed comparing LoRA and \textsc{dkps}
        (Data-Kernel Perspective-Space) representations, using Bayesian / Gaussian-
        process regression to unify selection and regression, and enriching the
        per-question feature vector with question-text embeddings --- directions
        pursued in Chapter~\ref{chap:effbench}.
\end{description}

%==============================================================================
\chapter{End-to-End Tests: Validating the Audit Instrument}
\label{chap:e2e}
%==============================================================================

\emph{Commits: 2026-06-03 through 2026-06-04, and the E2E deck of 2026-06-02.}

\section{Motivation}

Before trusting the audit pipeline's verdicts, the pipeline itself had to be
validated end to end on a case where the ground truth was controllable. The
audit workflow has six stages: (1) stand up a vLLM server for the target model
(via \code{infer-stack}, over LiteLLM / docker-compose); (2) write a benchmark
bundle; (3) execute runs on the servers; (4) aggregate run artifacts; (5) index
them; (6) build the analysis. The tools involved --- \code{infer-stack},
\code{eval-audit}, \code{magnet}, \code{kwdagger}, \code{cmd\_queue} --- each have
seams where a subtle mismatch produces a plausible-but-wrong comparison.

\section{The controlled experiment}

The subject was \code{microsoft/phi-2} on \code{mmlu:philosophy}, run under three
deliberately chosen deployments:

\begin{itemize}[nosep]
  \item \textbf{HF}: \textsc{helm}'s in-process \code{HuggingFaceClient},
        \code{temperature=0};
  \item \textbf{vLLM}: served through vLLM, \code{temperature=0} --- should match
        HF up to engine numerics;
  \item \textbf{vLLM ``incomparable''}: served through vLLM at \code{temperature=1}
        --- a deliberate negative control that \emph{should} disagree, proving the
        instrument can detect a real difference.
\end{itemize}

\section{Three instrument bugs found and fixed}

Validating the instrument surfaced three ways the comparison machinery could
silently mispair runs --- each fixed so that equivalent runs are actually
recognized as equivalent:

\begin{enumerate}
  \item \textbf{Hard-coded serving ports.} \code{infer-stack} used hard-coded
        ports without exposing them in a form \textsc{helm} could consume. Fix:
        read ports from the configuration source of truth and export them into a
        \code{.env} file queryable by the rest of the pipeline.
  \item \textbf{The \code{--groups} flag.} Public benchmark runs carry a
        \code{--groups} flag that has no effect on execution but prevented
        equivalent local runs from being matched to their public counterparts. Fix:
        ignore \code{groups=} tokens when locating comparable public runs, and emit
        a warning that this normalization occurred.
  \item \textbf{The \code{--temperature} run-name elision.} \textsc{helm} strips
        the temperature from the run \emph{name}, creating a divergence between the
        run spec and the run-directory name, so temperature-varying runs could not
        be located. Fix: locate run folders more robustly by reading the original
        run spec out of the generated artifacts rather than trusting the directory
        name.
\end{enumerate}

These three fixes are small but foundational: they are the first instances of a
theme that dominates the rest of the internship --- \textsc{helm}'s run identity is
carried by a drift-prone \emph{string}, and robust reproduction requires matching
on \emph{canonicalized recipe content}, not names (revisited at scale in
\S\ref{sec:canonkey} and in gotchas G1--G2, G10--G12; see
Appendix~\ref{app:gotchas}).

\section{Outcome and next steps recorded at the time}

The e2e tests established that HF and vLLM at \code{temperature=0} agree on the
core metric while the \code{temperature=1} control diverges as intended --- i.e.\
the instrument both confirms equivalence and detects real differences. The deck's
next steps (automate the e2e suite into CI, improve the warnings the systems emit,
and begin auditing more public models) foreshadow the rest of the summer.

%==============================================================================
\chapter{Corpus Cartography: Mapping the Reproducible Subset}
\label{chap:cartography}
%==============================================================================

\emph{Status deck of 2026-06-09.}

\section{The question: what \emph{can} be run at all?}

Before reproducing anything, one must know the shape of the corpus and which
slice is even a candidate. A statistics-and-filtering script was written to crawl
a directory of \textsc{helm} runs, extract per-run metadata, and emit a machine-
readable filter inventory (\file{filter\_inventory.json}). For each run it records
the model (parameter count, access class \{open, limited, closed\}), the scenario
and its class path, the benchmark family, whether the task requires a closed
(proprietary) judge, and a \emph{selection status} with explicit
\emph{failure reasons}. A representative row (a Thai-exam Typhoon-8B run) carries
fields such as \code{model\_num\_parameters}, \code{model\_access},
\code{model\_has\_hf\_client}, \code{size\_threshold\_params}, \code{eligible\_model},
\code{candidate\_pool}, and a \code{failure\_reason\_summary} --- the vocabulary
that the Stage-1 Sankey diagrams later consume.

\section{The numbers}

The crawl over the public corpus produced a stark funnel
(Table~\ref{tab:corpusstats}). Of \textbf{34{,}512} total runs, only
\textbf{1{,}109} were \emph{selected} as candidates for local reproduction under
the open-weight, runnable-recipe policy. The model-access distribution alone
explains much of the attrition: nearly as many \emph{limited}-access (18{,}273)
as \emph{open} (12{,}823) runs, plus 508 closed and 2{,}894 with no recorded
access class. Parameter counts were available for only 13{,}979 runs, so a size
budget could be applied to fewer than half the corpus with certainty.

\begin{table}[H]
\centering
\begin{tabular}{lr}
\toprule
\textbf{Quantity} & \textbf{Count} \\
\midrule
Total runs discovered            & 34{,}512 \\
Structurally incomplete          & 14 \\
Eligible model, out of scope     & 5 \\
\textbf{Selected (candidates)}   & \textbf{1{,}109} \\
\midrule
Model access: open               & 12{,}823 \\
Model access: limited            & 18{,}273 \\
Model access: closed             & 508 \\
Model access: none recorded      & 2{,}894 \\
\midrule
Runs with parameter-count data   & 13{,}979 \\
\midrule
Full grid points (models $\times$ scenarios) & 78{,}690 \\
Full grid: models / scenarios    & 366 / 215 \\
Filtered grid points             & 2{,}013 \\
Filtered grid: models / scenarios & 33 / 61 \\
\midrule
Runs requiring a closed judge    & 484 \\
Runs excluded \emph{solely} for closed judge & 25 \\
\bottomrule
\end{tabular}
\caption{The public-\textsc{helm} corpus funnel as measured on 2026-06-09. The
runnable subset is roughly 3\% of the corpus, and grid coverage is sparse: 198
of 2{,}013 filtered grid points are actually populated (``likely not all
scenario/model pairs make sense'').}
\label{tab:corpusstats}
\end{table}

\finding{The closed-judge dependency is real but not dominant.} Of 484 runs that
require a closed judge, only 25 are excluded \emph{solely} for that reason --- the
rest fail an independent filter too. This number motivated the later
\emph{open-judge extension} (\S\ref{sec:judge}): admitting closed-judge benchmarks
behind a flag and re-running them with an open judge recovers a
non-trivial-but-bounded slice.

\finding{Grid coverage is sparse and non-rectangular.} The full model$\times$scenario
grid has 78{,}690 points but only 6{,}844 are covered; the filtered grid has 2{,}013
points and 198 covered. This is the empirical basis for treating reproduction as a
\emph{coverage problem} --- the denominator of any ``fraction reproducible'' claim
must be the set of populated, in-scope grid points, not the Cartesian product. It
directly motivates the three-level coverage funnel (logical / recipe-canonical /
recipe-identical) that the pipeline reports (\S\ref{sec:funnel}).

\section{Why this is Stage 1}

This cartography is Stage 1 of the pipeline
(\file{index\_historic\_helm\_runs.py}): it discovers runs, applies the eligibility
filters, and emits the filter report whose Sankey diagram shows what was kept or
dropped and \emph{why}. Its output is the ``denominator'' the architecture
document (ADR-6) insists must remain explainable: every excluded run carries its
typed exclusion reason, so a reviewer can audit the claim ``we reproduced $X$ of
$Y$ runnable runs.'' Crucially, the exclusion reasons here are almost all
\emph{recipe/environment} reasons (\S\ref{sec:distinction}) --- gated model,
closed judge, no local deployment, over the size budget --- not reproducibility
verdicts.

\begin{quote}\small
\textbf{Corpus caveat (recorded later).} The curated candidate set that drives the
reproduction runbooks (\file{configs/run\_details.yaml}) is a hand-selected
Pythia/Vicuna/Qwen/OLMo subset, not the full eligible list; and an on-disk
\file{filter\_inventory.json} captured at one moment can go stale relative to the
live corpus (e.g.\ a modern-only pass that contains zero classic-era rows). The
distinction between ``the corpus'' and ``the curated candidate set'' matters when
reading any coverage number.
\end{quote}

%==============================================================================
\chapter{The Unified Comparison Core}
\label{chap:core}
%==============================================================================

\emph{Commits: 2026-06-11 through 2026-06-12 (roughly 30 commits); the ``Phase 0--3''
refactor and the ``Phase 3 / 4.0--4.9'' comparison-core unification.}

Two days of intense work rebuilt the analysis substrate so that every subsequent
reproducibility verdict flows through a single, framework-free comparison engine.
This is infrastructure, but it is infrastructure with research consequences: it is
what makes the later findings \emph{auditable} rather than asserted.

\section{The refactor (Phases 0--2): making the instrument legible}

The repository had grown to $\sim$86 modules and $\sim$32.8k lines with 18 console
scripts, and several ``god modules'' had accreted --- notably
\file{build\_reports\_summary.py} at 5{,}369 lines / 108 functions. The refactor
proceeded in phases: root hygiene and finishing an in-flight
\code{vllm\_service}$\rightarrow$\code{infer\_stack} rename (Phase 0); unifying the
CLI surface so every console script resolves to a thin \code{eval\_audit.cli.*}
wrapper (Phase 1); and decomposing the god modules into cohesive subpackages
(Phase 2).

The methodological point worth recording for the paper's ``reproducible software''
story is \emph{how} the splits were verified. Rather than hand-moving functions,
the god modules were split programmatically: AST-parse the module, compute the
top-level-symbol reference graph, validate the cluster assignment is a DAG,
generate each submodule with exactly the imports its symbols reference, then assert
that all top-level symbols are \emph{AST-identical} to \code{HEAD}. This gives a
machine-checkable ``pure relocation'' guarantee. A recurring subtlety --- that
monkeypatching a facade module no longer reaches functions that bind their
collaborators in a new home module --- caused several test fixes and is a lesson
about characterization tests surviving refactors.

\section{Phase 3 / \code{NormalizedDiff}: one core, many adapters}
\label{sec:normalizeddiff}

The deeper change was to unify two forked comparison cores. The \textsc{helm} path
computed agreement and diagnosis through \code{HelmRunDiff} (a 1{,}902-line class);
the \eeeshort path computed overlapping numbers through \code{normalized/compare.py}.
Both were invoked from one function, \code{\_build\_pair}, gated by a
\code{skip\_diagnosis} flag. The agreement math was already framework-free on both
paths; only the \emph{diagnosis} block and the \file{run\_spec.json} semantic diff
were \textsc{helm}-specific.

The unification introduced \code{eval\_audit/normalized/diff.py::NormalizedDiff}:
one framework-free core consuming two \code{NormalizedRun} objects and producing
agreement rows, agreement curves, quantiles, the facts-grade diagnosis, and the
judge-dependence metric-class split. The \textsc{helm} and \eeeshort paths became
\emph{thin adapters} over this one core; \code{HelmRunDiff}'s \code{\_diagnose\_repro}
was reduced to delegating to \code{normalized.diagnose.diagnose\_repro}.

The design was executed \emph{additive-first}, which is the load-bearing rigor:
\code{NormalizedDiff} was built next to the incumbent and proven \emph{byte-equal}
to committed baseline snapshots (fixtures F1--F10, gated at numerical tolerance
\code{atol}$=10^{-9}$, with the explicit rule that a looser tolerance is ``a
finding to investigate, not a knob to widen'') \emph{before} any renderer was
flipped to use it. Inverting the usual order --- build the replacement, prove
equivalence, then swap --- turned a ``high-risk hinge'' into two routine commits.

\section{The framework-free taxonomy and recipe-facts resolver}
\label{sec:taxonomy}

Three supporting modules were lifted out of the \textsc{helm}-specific package so
they carry no framework dependency:

\begin{itemize}[nosep]
  \item \file{eval\_audit/metrics\_taxonomy.py} --- classifies every metric as
        \emph{core} vs \emph{bookkeeping} (the metric-scale caveat from
        \S\ref{sec:inherited}) and, new here, as \emph{judge-dependent} vs
        \emph{deterministic} (\code{classify\_judge\_dependence}). The latter
        split lets a report separate the part of a comparison that is a
        \emph{reproduction control} (deterministic metrics) from the part that is
        a \emph{substitution measurement} (judge-dependent metrics).
  \item \file{eval\_audit/normalized/recipe\_facts.py} --- one resolver answering
        ``what recipe produced this run?'' in priority order: a native
        \code{recipe\_facts} block (JSON-encoded inside an \eeeshort aggregate's
        \code{source\_metadata.additional\_details}) $\rightarrow$ a sidecar
        \file{run\_spec.json} $\rightarrow$ \emph{unknown} (facts collapse to
        \code{status='unknown'} and the pipeline emits
        \code{comparability\_unknown:*} warnings). This is the seam through which
        the later ``unrecorded parameters'' story (Appendix~\ref{app:unrecorded})
        could be closed without a schema change.
  \item \file{eval\_audit/normalized/diagnose.py} --- the single input-to-label
        diagnosis implementation, including declared-substitution re-labeling
        (\code{intended\_substitution:*}).
\end{itemize}

\section{The open-judge extension (R2)}
\label{sec:judge}

A one-hour ``spike'' against the real corpus reshaped a planned feature and
recorded a genuinely useful finding
(\file{docs/planning/judge-identity-inventory.md}). The question was: where does
judge-model identity live for the six closed-judge benchmarks? Three findings:

\begin{enumerate}[nosep]
  \item \textbf{Judge identity is not in \file{run\_spec.json}.} Annotators carry
        \emph{empty} args; the judge model is hard-coded inside the \textsc{helm}
        annotator class as a function of \textsc{helm} version (e.g.\ a GPT-4o +
        Llama-3.1-405B ensemble). Officials therefore need a curated
        (annotator class, version)$\rightarrow$judge-models map, implemented as
        \file{eval\_audit/judge\_registry.py}.
  \item \textbf{The official ensemble already contains an open judge.} Because
        \textsc{helm} records per-judge sub-scores (e.g.\ \code{safety\_gpt\_score},
        \code{safety\_llama\_score}) as separate metrics, a \emph{same-judge
        reproduction control already exists inside the official data} --- one can
        compare \code{safety\_llama\_score} official-vs-local without any new run.
        This reframes substitution as per-metric rather than per-run.
  \item \textbf{Honesty scoping.} The \code{same\_judge} comparability fact is
        emitted \emph{only} for declared-substitution comparisons; emitting it
        everywhere would spray \code{comparability\_unknown:same\_judge} noise over
        every legacy pair. A crucial design decision recorded here: facts stay
        honest (\code{same\_judge: no} is never re-labeled to a ``substituted''
        status); only the \emph{diagnosis} re-labels the difference as
        \code{intended\_substitution:judge}. A ``substituted'' fact status was
        rejected as dishonest-by-construction.
\end{enumerate}

Stage 1 gained \code{--allow-closed-judge-benchmarks}, admitting these runs through
a distinct \code{judge-substitution} candidate pool that rides
\code{run\_details}$\rightarrow$manifests$\rightarrow$audit index to the planner.

\evidencebox{Status at the 2026-07-27 revision}{This design became an executed
experiment in July: six benchmarks reconstructed and identity-replayed at machine
precision, then rejudged by six open-weight judge arms across three replicates. See
Chapter~\ref{chap:openjudge}. The registry described here is what makes the
official side of that comparison addressable at all.}

\section{Why the entry points stayed separate}

A tempting simplification --- collapse the \textsc{helm} and \eeeshort CLIs into
one auto-detecting command --- was deliberately \emph{rejected}, and the reasoning
is worth preserving because it recurs. The \eeeshort-only path must remain
physically capable of building its full report tree \emph{with zero \textsc{helm}
artifacts on disk} (the ``EEE-only operability'' gate), because one of the paper's
case-study claims is that the \eeeshort per-instance schema alone suffices for
reproducibility analysis --- a claim a reviewer will want to verify by
\code{grep}, not by trusting a code review. Merging the CLIs would import the
\textsc{helm} adapter into the \eeeshort path and destroy that grep-checkable
guarantee.

%==============================================================================
\chapter{The OLMo Campaign}
\label{chap:olmo}
%==============================================================================

\emph{Commits: 2026-06-12 onward; status decks of 06-16, 06-23, 06-30, 07-08,
07-14. The OLMo reproduction is the internship's central case study, and the
vehicle through which every piece of reproduction machinery was hardened.}

\section{The target grid}

Six open-weight OLMo-family models were chosen as the reproduction subjects
(Table~\ref{tab:olmogrid}). Four are recent instruct models evaluated on the
capabilities benchmarks (\code{bbq}, \code{gpqa}, \code{ifeval}, \code{mmlu\_pro});
two are older base models with broad classic-track coverage. The grid spans both
\emph{no-chain-of-thought / temperature=0} benchmarks (bbq, commonsense, gsm,
ifeval, legalbench, med\_qa, mmlu, narrative\_qa, wmt\_14) and \emph{chain-of-thought
/ temperature=1} benchmarks (gpqa, mmlu\_pro).

\begin{table}[H]
\centering
\small
\begin{tabularx}{\textwidth}{lY}
\toprule
\textbf{Model} & \textbf{Benchmarks reproduced} \\
\midrule
\code{olmoe-1b-7b-0125-instruct}   & bbq, gpqa, ifeval, mmlu\_pro \\
\code{olmo-2-1124-7b-instruct}     & bbq, gpqa, ifeval, mmlu\_pro \\
\code{olmo-2-1124-13b-instruct}    & bbq, gpqa, ifeval, mmlu\_pro \\
\code{olmo-2-0325-32b-instruct}    & bbq, gpqa, ifeval, mmlu\_pro \\
\code{olmo-7b}                     & mmlu (62), math (7), legalbench (5), wmt\_14 (5), \\
                                   & natural\_qa (2), commonsense (1), gsm (1), \\
                                   & med\_qa (1), narrative\_qa (1) \\
\code{olmo-1.7-7b}                 & mmlu (57) \\
\bottomrule
\end{tabularx}
\caption{The six OLMo reproduction subjects and their benchmark coverage. The
instruct models exercise the chat-template and judge-dependent paths; the base
models exercise the classic-track, completions path.}
\label{tab:olmogrid}
\end{table}

Early smoke tests already surfaced two environment issues that recur:
a shared-machine configuration-file conflict, and \code{olmo-7b}'s tokenizer
attempting to run external code (an early symptom of the EOS-append bug diagnosed
in \S\ref{sec:eosappend}). These were fixed via HELM tokenizer-alias corrections,
adding \code{sentencepiece}/\code{tiktoken} tokenizer backends as dependencies, and
loading the OLMo-7B tokenizer through a HELM entry-point plugin.

\section{Chat-template overhead blows the context window}

The first substantive reproducibility bug in the OLMo runs was a
\code{ContextWindowExceededError}: the chain-of-thought runs
(\code{num\_output\_tokens=2048}) on the 4096-context chat models were rejected by
vLLM. The mechanism is subtle and paper-worthy. \textsc{helm}'s
\code{LocalWindowService} truncates the prompt to
$\text{max\_seq\_len} - \text{expected\_completion} = 4096 - 2048 = 2048$ tokens,
but it counts the \emph{raw} prompt, \emph{before} the chat template is applied.
When the run is routed through vLLM's OpenAI-compatible chat endpoint, the OLMo-2 /
OLMoE chat template wraps the prompt and adds tokens \textsc{helm} never counted.
Measured with the actual tokenizers, the single-turn wrapper overhead is
\textbf{12 tokens (OLMo-2)} and \textbf{13 tokens (OLMoE)}, so
$2048 + 13 + 2048 = 4109 > 4096$ and vLLM hard-rejects.

The fix reserved a 32-token margin per preset via the sanctioned
\code{helm\_max\_sequence\_and\_generated\_tokens\_length} knob (already used by the
Vicuna chat path), and repaired a latent bug where the flat-preset code path
silently dropped that override. The general lesson --- \emph{when \textsc{helm}
drives a chat model through an OpenAI-compatible server, its prompt-token budget is
blind to the chat template the server applies} --- is the first of several
chat-template-versus-window interactions, and it recurs when the Qwen turbo
endpoints hit the same wall in July (\S\ref{sec:qwenwindow}).

\section{Data-access failures are filtering reasons, not results}

Several OLMo benchmarks could not be run at all, and classifying them correctly was
the point:

\begin{itemize}[nosep]
  \item \code{competition\_math} was disabled on the HF Hub (script-based dataset,
        401 on the loader).
  \item \code{natural\_qa}: the public GCS bucket revoked \emph{authenticated}
        reads (it dropped \code{allAuthenticatedUsers}, not just \code{allUsers}),
        so no credential unblocks it. An entire staging shim built to work around
        the assumed ``authenticated reads work'' had to be reverted once the
        access assumption was falsified by a single test with real credentials.
        The recorded lesson --- \emph{validate the access assumption before
        building the plumbing that depends on it} --- is a methodological point for
        the paper's failure taxonomy.
  \item \code{gpqa} is a gated HF dataset (requires accepting terms).
\end{itemize}

Each was disabled as a documented \emph{environment} failure, joining a taxonomy
that is deliberately kept separate from the reproducibility verdicts.

\section{The canonical-key matching bug: 114 phantom misses}
\label{sec:canonkey}

The most consequential ``bug'' of the OLMo campaign was not in execution but in
\emph{pairing}. The planner matched logical-key \emph{strings} by generating
separator permutations and (under a flag) stripping \code{groups=} tokens, but it
never canonicalized \emph{token order}. The OLMo MMLU keys are the same token set
in a different order, so the local and official variant sets had an \emph{empty
intersection}: all 114 MMLU packets were reported as
\code{missing\_official\_component} even though the public counterparts existed in
the index.

The fix (\file{docs/planning/core-report-planner-robust-matching-plan.md}) replaced
the order-sensitive matcher with a single \code{canonical\_logical\_key} in
\file{eval\_audit/run\_entries.py}: parse the \code{benchmark:k=v,...} name, drop
bookkeeping-only tokens (\code{groups}, \code{model\_deployment}), canonicalize
values (model \code{/}$\leftrightarrow$\code{\_}, \code{mmlu\_pro} \code{subject}
$\rightarrow$ \code{subset}), and re-serialize with the key--value pairs
\emph{sorted by key}. Once you sort, the separator-permutation and groups-strip
variants all fold into one deterministic string, so keeping both matchers would
have re-created the very ``two competing normalizers'' divergence the fix set out
to end. The real-data validation expectation: \code{n\_skipped} $114\rightarrow\sim0$,
\code{n\_built} $35\rightarrow\sim149$.

This is a symmetric equivalence, correct for \emph{grouping}; the asymmetric subset
test (\code{run\_dir\_matches\_requested}) was deliberately preserved for
\code{compare\_batch}, which answers the different question ``does this candidate
satisfy this request.'' The distinction --- symmetric equivalence for grouping vs
asymmetric subset for satisfaction --- is a reusable design point.

\section{The BBQ prompt drift: a genuine recipe disagreement}

Not every difference is an artifact to be fixed. The BBQ runs surfaced a
\emph{real} recipe drift: the official run entry carried
\code{output\_format\_instructions=mcqa} and the prompt began ``Answer with only a
single letter.'' (followed by the standard preamble) while the locally
reconstructed run entry dropped that token and produced only ``The following are
multiple choice questions (with answers).'' The model sees a different prompt and
produces measurably different output. This is gotcha G5 (\S\ref{app:gotchas}), and
it is precisely the kind of difference that must be \emph{surfaced} (via the
\code{same\_instructions} comparability fact), not silently fixed --- and it is one
of the strongest motivations for the from-spec replay of
Chapter~\ref{chap:fromspec}: replaying the official \file{run\_spec.json} verbatim
makes such reconstruction drift vanish by construction.

%==============================================================================
\chapter{Containerization and GPU Leasing}
\label{chap:containers}
%==============================================================================

\emph{Commits: 2026-06-16 through 2026-06-24. Two intertwined infrastructure
efforts: pinning the execution environment in Docker, and letting many HELM runs
fan out across GPUs under a lease.}

\section{Why containerize: the environment is a variable}

By default \code{eval-audit-run} executes each \textsc{helm} run-entry in the host
venv of whatever GPU machine the scheduler lands on, making torch / CUDA /
transformers / \textsc{helm} versions an \emph{uncontrolled variable} --- exactly
the kind of drift the audit exists to separate from residual disagreements under controlled equivalence.
The containerized path runs each run-entry inside a pinned Docker image and records
\emph{which image, by \code{sha256} digest, produced each run}.

The design (confirmed interactively) is worth recording:

\begin{itemize}[nosep]
  \item A brand-new, independent multi-stage Dockerfile (CUDA devel$\rightarrow$
        runtime, \code{uv}-managed standalone CPython, venv at \code{/opt/venv}),
        built from \emph{pristine committed submodule state} via \code{git archive}
        at each submodule's gitlink SHA --- no remote dependency, no \code{.git}
        leakage, SHAs recorded as OCI labels.
  \item The image tag is resolved to an immutable \code{<repo>@sha256:<digest>}
        reference \emph{once at schedule time} and every node is pinned to it, so a
        rebuild automatically re-pins.
  \item The container runs as root; a thin entrypoint chowns \emph{only} the output
        dir back to \code{HOST\_UID:HOST\_GID} on exit, so \code{kwdagger} on the
        host can own results while in-container HF-cache writes still succeed.
  \item The docker-wrapping \code{kwdagger} node \emph{subclasses} the magnet
        materialize node and overrides only \code{command}, inheriting the
        \code{out\_paths}/\code{DONE}-sentinel contract for free --- the reusable
        lesson being ``wrap an existing scheduler node in a new execution substrate
        by subclass + override, not reimplementation.''
\end{itemize}

Several first-real-build bugs are instructive because each could only surface at
the first genuine build or run and each maps to a general lesson:

\begin{itemize}[nosep]
  \item \textbf{Dangling interpreter symlink} --- \code{uv}'s managed standalone
        Python was not copied into the final stage, so \code{/opt/venv/bin/python}
        was a dangling symlink reporting ``command not found.'' Lesson:
        \emph{uv-managed Python + multi-stage Docker = copy the interpreter, not
        just the venv}, and put the import smoke-test in the stage that ships.
  \item \textbf{HF token not reaching the container} --- \code{kwdagger} runs each
        job in a fresh tmux pane whose environment is deliberately empty (secrets
        hygiene), so \code{-e HF\_TOKEN} forwarded nothing; the host-venv path had
        silently relied on the on-disk \code{huggingface-cli login} token. Fix:
        write the resolved token to \code{<hf\_cache\_dir>/token} (mode 0600) at
        schedule time, in the process that \emph{did} inherit the env. Lesson:
        \emph{a redundant channel that silently carries the load hides the breakage
        of the primary one}.
  \item \textbf{Wrong extras and an unpinned hub} --- the image installed
        \code{crfm-helm[heim]} (image-generation extra) which lacks
        \code{langdetect} (needed by ifeval), and left \code{huggingface\_hub}
        unpinned so it floated to a version whose repo-resolution API broke
        \code{wmt\_14}. Fix: install \code{crfm-helm[all]} and co-pin
        \code{huggingface\_hub==0.36.2}, asserted in \emph{both} the install layer
        and the final stage. Lesson: \emph{pin the secrets-of-resolution
        (transitive floats), not just the top-level package}, and pair a build-time
        guard with a runtime probe of the resolved artifact (a stale pinned digest
        is invisible to the build guard).
\end{itemize}

\section{From profiles to a leasing catalog}

Mid-June the \code{infer\_stack} submodule was rewritten from a
profile/contract model to a \emph{catalog} (models + endpoints) plus \emph{leasing}
(\code{acquire}/\code{release}) model. The two in-scope runbooks (phi-2 e2e and the
six OLMo presets) were migrated: \code{load\_profile\_contract} was replaced with a
pure static \code{resolve\_serving\_facts} over the catalog; the config dirs were
reschema'd from \code{config.yaml}+\code{models.yaml} to
\code{settings.yaml}+\code{catalog.yaml}; and an explicit \code{protocol\_mode}
(completions for base models, chat for instruct) was made a required preset fact
rather than a silent default. A subsequent question --- ``do \code{serve} and
\code{acquire} both need to exist?'' --- collapsed the two verbs into
\code{acquire} once a one-line grep proved the distinguishing \code{owner='manual'}
field drove no behavior.

\section{The leasing pipeline: preconditions as job phases}

The deeper design question was how to drive \emph{high-parallelism} reproduction:
each \textsc{helm} run is a \code{kwdagger} \code{ProcessNode} that must bracket
itself with a GPU lease (acquire before, release after), where release must run
after finish \emph{or} crash. The decisive argument, reached across three rounds of
design discussion, was \emph{co-location}: the lease lives in a node-local ledger
and the gateway is localhost, so release must run \emph{where} acquire ran. A
separate cleanup DAG node cannot promise that (and is actively wrong on multi-node
scheduling), and success-only dependencies (\code{afterok}) skip a release node
when the run fails --- the exact leak to avoid. The right primitive is a job
\emph{phase} (setup/teardown), which is also the more general, reusable scheduler
feature.

The findings that settled the design were verified against source, not assumed:
\code{cmd\_queue}'s \code{preamble} already gates the main command, so only an
always-run \code{teardown} trap was new; \code{kwdagger}'s cache guard
(\code{test -e DONE || cmd}) means a \emph{skipped} job acquires nothing; and a
blocking acquire's retry loop \emph{reclaims TTL-expired leaks while it waits}, so
queue-and-wait and leak-recovery become one mechanism --- \emph{iff every pipeline
lease has a finite TTL}. A subtle bash-precedence trap was caught and closed: a
setup chain \code{acquire \&\& \{ snapshot \} || true} lets a \emph{failed}
acquire fall through to \code{|| true}, defeating the gate; the fix scopes
\code{|| true} inside the snapshot brace so only the snapshot is best-effort.

Two further design moves are worth recording because they generalize:

\begin{itemize}[nosep]
  \item \textbf{Decoupling leasing from containerization.} The lease bracket was
        initially bolted onto the docker node only, which would have dragged
        containerization onto the host-venv vLLM scenarios. It was extracted into a
        transport-agnostic \code{LeaseBracketMixin} + shared command renderer;
        leasing and containerization became \emph{orthogonal} axes (a run can be
        docker+lease, docker+no-lease, or --- briefly --- bare+lease). The lesson:
        when a capability seems ``coupled'' to a transport, check whether it is the
        capability that needs the transport or just where someone bolted it on.
  \item \textbf{Scoped cleanup over global teardown.} Every \code{release --all
        --evict} was removed from the runnable path in favor of per-run release
        (\code{reclaim: stop}) plus a scoped \code{infer-stack gc} that reclaims
        only TTL-expired leaks in this user's ledger --- so the shared docker stack
        is never torn out from under a co-tenant. A bootstrap that exists only to
        read a managed secret pays for it with a scoped \code{--env-file} handle,
        not a global teardown.
\end{itemize}

By the end of this arc, containerization was made \emph{mandatory} and the
non-containerized code paths were deleted, with leasing the orthogonal opt-in ---
so the execution environment is always a pinned, digest-recorded artifact
(\code{container\_provenance.json}).

%==============================================================================
\chapter{Faithful Replay: From Run-Entry to Run-Spec}
\label{chap:fromspec}
%==============================================================================

\emph{Commits: 2026-06-25 through 2026-06-30; the \code{impl/run-from-run-spec}
branch on which most subsequent work sits.}

\section{The core methodological principle}

This chapter contains what is arguably the single most important methodological
decision of the internship, and it should anchor the paper's ``methods'' section.

Public \textsc{helm} runs do not record the CLI run-entry string that produced
them, nor the exact \textsc{helm} version. What they \emph{do} ship is a fully
resolved \file{run\_spec.json}: the complete adapter spec, scenario spec, metric
specs, and annotators, with every default already materialized. The prior approach
--- and the approach every earlier reproduction took --- was to \emph{reconstruct}
a run-entry string from the run and re-parse it. But re-parsing re-derives the
recipe under \emph{today's} library defaults, which silently drifts from the
recipe that actually produced the official numbers (this is the mechanism behind
gotchas G1, G5, and the BBQ drift of \S\ref{sec:canonkey}).

\begin{quote}\itshape
The principle: never reconstruct the recipe. Replay the resolved
\file{run\_spec.json} \emph{verbatim}.
\end{quote}

Concretely, a new \code{materialize\_helm\_run\_from\_spec.py} CLI (a faithful-
replay sibling of the run-entry materializer) reuses the discovery/matching
scaffolding but swaps the compute step from ``reconstruct a run-entry string
$\rightarrow$ \code{helm-run} subprocess'' to ``\code{from\_json(run\_spec.json)}
$\rightarrow$ in-process \code{run\_benchmarking}.'' In discovery mode, the
run-entry string now does nothing but \emph{locate} the official directory; the
authoritative resolved recipe drives execution. This memory --- \emph{all
reproductions must be from-spec} --- became a standing rule, and the older OLMo
run-entry path was flagged for migration.

\section{Two failure modes, one loud and one silent}

The design carefully distinguishes the two ways a from-spec replay can go wrong,
because they demand opposite handling (this is gotcha-adjacent and reviewer-
relevant):

\begin{description}[style=nextline,leftmargin=1.25em]
  \item[Class drift is loud.] If the \file{run\_spec.json} references a scenario or
        metric class that does not resolve under the local \textsc{helm}
        (\code{get\_class\_by\_name} raises \code{ImportError}), the preflight
        catches it before any instance runs. This is the correct behavior: a class
        that cannot be resolved is an environment mismatch, and it is surfaced as
        such. (The preflight's exception net had to be broadened from
        \code{(ImportError, AttributeError)} to \code{Exception}, because importing
        a vision--language scenario raises \textsc{helm}'s
        \code{OptionalDependencyNotInstalled} --- itself a recipe/environment
        filtering reason, not a crash to propagate.)
  \item[Field drift is silent.] The cattrs codec ignores unknown keys and fills
        defaults, so a dropped field produces no error --- only wrong numbers.
        The guard is a round-trip test: because \textsc{helm} \emph{writes}
        \file{run\_spec.json} with \code{asdict\_without\_nones} (dropping only
        \code{None} values), every on-disk key is non-\code{None}, and the
        re-serializer re-emits every non-\code{None} field --- so a key missing
        after \code{from\_json}$\rightarrow$\code{to\_json} is an unambiguous silent
        drop. Ten sampled MMLU specs round-tripped with zero key loss.
\end{description}

The reusable insight recorded at the time: \emph{writer/serializer asymmetry
decides whether a diff test is sound} --- the no-key-dropped assertion holds only
because the writer drops a superset of what the re-serializer drops.

\section{Making the deployment substitution visible}

A from-spec replay serves the model on local hardware, which is by definition a
\emph{different deployment} than the original (e.g.\ a hosted Together API). The
subtle danger is that a pure by-name replay would record the \emph{official}
deployment name in the produced run, so the comparison would report
\code{same\_deployment=yes} and the single most important difference the audit
exists to surface --- that the serving engine changed --- would become invisible.

The fix (\file{from-spec-deployment-rewrite-plan.md}) rewrites
\code{adapter\_spec.model\_deployment} to the \emph{local} name after
deserialization, via an explicit \code{--model-deployment} argument. Explicit,
not auto-derived: the whole feature exists to stop \emph{silently} masking the
substitution, so its own failure mode must be loud --- a wrong name is a
\textsc{helm} ``deployment not found'' crash \emph{before any instances run},
whereas auto-derive's failure mode is the silent no-op that fails \emph{toward} the
very bug it fixes. The produced run then records the served endpoint, and the
existing diff logic correctly reports \code{same\_deployment=no} with no downstream
plumbing.

\section{Addressing runs by (root, relative-path) and content-addressed caching}

A further refinement replaced the drift-prone run \emph{name} as the addressing
mechanism with \emph{host-side} resolution by \code{(precomputed\_root,
rel\_path)}: resolve the official \file{run\_spec.json} at schedule time, apply
\emph{raw-JSON} scalar edits (only \code{adapter\_spec.model\_deployment} and
\code{max\_eval\_instances}, so no cattrs round-trip and therefore no field drift),
and materialize a substituted copy under a staging directory. Each run rides one
\code{kwdagger} \emph{submatrix} entry (the verified zip primitive) so the per-run
tuple travels together without an $N\times N$ cross-product. Two consequences worth
recording:

\begin{itemize}[nosep]
  \item \textbf{No image rebuild, no corpus mount.} The runner needs only the tiny
        staging directory, because substitution happens host-side.
  \item \textbf{Content-addressed caching.} \code{max\_eval\_instances} is baked
        into the materialized copy but is not a matrix parameter, so a naive fixed
        filename (\code{run\_spec.json}) would make \code{kwdagger} \emph{skip} and
        reuse a stale run after a cap change. The fix appends a content hash to the
        filename: \code{run\_spec.<sha256(bytes)[:16]>.json}. Because the
        materialized path \emph{is} the node's algorithm identity, an identical
        recipe yields the same path (skip) and any changed field yields a new path
        (recompute). This is the reusable ``content-address generated input files
        so baked-in changes invalidate the cache'' rule.
\end{itemize}

An exporter ``auto-freeze'' resolves each preset run-entry to its exact rel-path
\emph{once} at export time (against the corpus snapshot) and freezes the sources
into the generated manifests --- the only remaining place token-subset discovery
runs, pinned where drift cannot occur --- and a dry-check existence mode validates
that a frozen manifest's rel-paths still resolve.

%==============================================================================
\chapter{Deployment Matching: The Flagship Investigation}
\label{chap:deploymatch}
%==============================================================================

\emph{Commits: 2026-06-30 through 2026-07-10; the \code{dev/tools/deployment\_match/}
tool and the documents \file{vllm-vs-huggingface-deployment-match.md} and
\file{helm-unrecorded-deployment-params.md}. This is the scientific heart of the
internship: the sustained effort to attribute residual disagreement to a specific,
unrecorded execution parameter.}

\section{The provocation: OLMo-7B produces prompt-independent gibberish}
\label{sec:eosappend}

The investigation began with a dramatic symptom. During from-spec reproduction of
\code{narrative\_qa}, local OLMo-7B repeated the same prompt-independent
boilerplate --- ``The first thing you need to do is to make sure that you have
a\ldots before you start playing.'' --- for \emph{every} instance, while the
official \code{together/olmo-7b} conditioned on the prompt correctly. The diff
showed \code{prompts\_equal=True} (the prompt reached the model) and a hugely
negative local completion log-prob against the official's $0$: the model ran but
ignored its input.

The first hypothesis was a documented one: OLMo-1 7B \code{float16} numerical
instability (activation outliers exceed fp16's range $\rightarrow$ inf/NaN in
attention $\rightarrow$ collapse to the highest-prior pretraining attractor). A
minimal deployment matrix was built to test it --- and the hypothesis was
\emph{refuted}. The true cause was a \emph{tokenizer} issue: \code{OLMo-7B-hf}'s
tokenizer appends an \code{<|endoftext|>} token to the prompt under vLLM's default
\code{add\_special\_tokens=True}; the Together deployment did not do this, and
\textsc{helm} never recorded the difference. The production fix serves the
OLMo-1.7 tokenizer (whose post-processor omits the append), and the completions
path forwards \code{add\_special\_tokens=false}.

\finding{This is the paradigm case of the whole project.} A disagreement that
looked like flat irreproducibility (the model outputs nonsense) was in fact a
\emph{recipe/environment} failure caused by a single unrecorded request-time
tokenizer flag. And the first, plausible theory (fp16 instability) was wrong ---
which is exactly why the response was to build a \emph{sweep} rather than trust a
single-knob theory. The lesson --- that single-knob theories about deployment
divergence had by now been wrong twice --- became the design charter for the
deployment-match tool.

\section{The tool: sweep, don't guess}

\code{dev/tools/deployment\_match/} takes \emph{any} public \textsc{helm} run,
builds a grid of plausible local deployments for that run's model, evaluates a
small length-stratified prompt sample, and ranks each cell by closeness to the
official completions. Because Together stores no log-probs, the objective is
\emph{agreement with the official completion text} (a composite of quasi-exact-
match, first-word match, and similarity, with a model-agnostic
``prompt-independence collapse'' detector). The grid is two-tier: serve-time knobs
(dtype $\times$ tokenizer $\times$ \code{max\_model\_len}) become distinct
infer-stack endpoints; request-time knobs (\code{add\_special\_tokens},
\code{add\_generation\_prompt}, protocol) are probed per request against one
container. The tool has \code{dry-run} / \code{run} / \code{confirm} / \code{auto}
phases, and \code{confirm} produces a real one-endpoint catalog you can
\code{acquire}.

A key design principle, formalized after the fp16 mis-diagnosis, splits ``exclude
settings beforehand'' into two opposite treatments:
\begin{itemize}[nosep]
  \item \textbf{Feasibility pruning (safe).} A recipe that physically cannot serve
        on this hardware teaches nothing and wastes a full serve cycle, so it is
        dropped a priori with a typed reason (mirroring the Stage-1 filter
        pattern). Example: \code{float32} on a MoE model $\rightarrow$
        \code{infeasible:moe-fp32-shared-mem} (see \S\ref{sec:fp32}).
  \item \textbf{Relevance pruning (risky) is refused.} ``This knob probably won't
        change the output'' is exactly what the tool will not guess; such knobs
        stay in the sweep unless the operator narrows an axis explicitly.
\end{itemize}

\section{Chat-template version drift: \code{add\_generation\_prompt}}
\label{sec:addgenprompt}

The next finding is a subtle, real reproducibility trap on OLMoE-instruct, and it
is a clean example of a \emph{software-version} dependency hiding inside a recipe.
\textsc{helm}'s \code{get\_prompt} always calls
\code{apply\_chat\_template(\ldots, add\_generation\_prompt=True)} --- but that flag
only does anything if the chat template \emph{implements} it
(\code{\{\% if add\_generation\_prompt \%\}<|assistant|>\textbackslash n\{\% endif \%\}}).
The OLMoE template shipped with \textsc{helm}'s \emph{older} transformers
\emph{ignored} the flag, so \code{add\_generation\_prompt} was effectively
\emph{False} and the model was fed the prompt \emph{without} the trailing
\code{<|assistant|>\textbackslash n}. Modern transformers/vLLM ship an updated
template that honors it, so the same call now \emph{appends} the suffix
$\rightarrow$ a different prompt $\rightarrow$ a different output on every cell.
Rendering with \code{add\_generation\_prompt=false} on modern transformers moves
the output back toward the official, confirmed empirically. The tool now sweeps
\code{add\_generation\_prompt} $\in \{$True, False$\}$ as a cheap request-time knob
and lets the scorer decide.

\section{The float32 discovery: a flagship attribution pilot}
\label{sec:fp32}

\evidencebox{Authoritative status}{This section reports two distinct forms of
evidence. The source/configuration path makes fp32 the most likely historical dtype
under bounded pre-v5 assumptions. Separately, a 12-instance IFEval probe found
float32 candidate cells with the highest sampled agreement. Neither establishes the
complete historical deployment, and the candidate has not yet been confirmed through
the ordinary HELM path on held-out data.
\newline\newline
\textbf{Status at the 2026-07-27 revision.} The confirmation this box calls for was
run. Read this section as the \emph{discovery} half of the result and
Chapter~\ref{chap:substrate-closed} as the confirmation half. The short version:
float32 is necessary but not sufficient, the ordinary-path float32 run leaves a
$+0.07$ IFEval residual, and that residual is the inference engine. The deductive
dtype argument below was independently corroborated --- reproducing the official at
float32 on \textsc{helm}'s own client is byte-exact --- so the ``most likely fp32''
conclusion is now observed rather than inferred for OLMo-2 instruct.}


The single most consequential finding of the internship concerns model
\emph{precision}, and it is the one most worth foregrounding in the paper.

\textsc{helm}'s \code{HuggingFaceClient} loads a model with only the kwargs the
deployment config supplies. For \code{huggingface/olmoe-1b-7b-0125-instruct}, the
config gives exactly \code{args: \{device\_map: auto\}} --- \emph{no}
\code{torch\_dtype}. So \code{AutoModelForCausalLM.from\_pretrained(\ldots,
device\_map="auto")} is called with no dtype, and \emph{every transformers 4.x}
defaults that to \code{torch.float32} for backward compatibility, \emph{ignoring
the checkpoint's bf16 config}. (Reading the checkpoint's dtype is a transformers
v5 change; even late 4.x still hits ``\code{else: set fp32 as the default dtype for
BC}.'') The OLMoE architecture landed in transformers $\sim$4.45, so a January-2025
run is squarely in the fp32-default regime. This is a \emph{deductive} claim ---
read off the loader source and the library's version-dependent default, not fit to
outputs --- so it is stronger than an observational match. Its residual is the one
parameter it depends on: the exact \code{transformers} version, which the run dir
does not record (only bounds). Under that reading:

\begin{quote}\itshape
The code path \textsc{helm} used loads the weights in \code{float32} under every
\code{transformers} version consistent with the OLMoE run's date --- so the
official most likely ran \code{float32}. (Conditional on \code{transformers}$<$5;
a run produced under \,$\geq$5 would default to bf16 and move this conclusion.)
\end{quote}

A small oracle-scored \emph{probe} corroborates this: reproducing the official on
the HF side at \code{float32} (\code{transformers.generate()}) reaches quasi/exact
match $\approx 1.0$ on the sampled instances, versus $\approx 0.17$ for the best
fp16 vLLM cell. Table~\ref{tab:fp32} shows the overnight probe sweeps across the
four OLMo instruct models. \textbf{This is deployment-\emph{matching} evidence, not
a completed HELM reproduction} (see the epistemic-status note below): the scores are
composites over $n\approx12$ ifeval instances, scored against the sampled official
completions, and the winning cells additionally rely on a probe-only
\code{add\_special\_tokens=False} request knob that an ordinary \textsc{helm} run
does not send.

\begin{table}[H]
\centering
\small
\begin{tabular}{lll}
\toprule
\textbf{Model} & \textbf{HF best} & \textbf{vLLM best} \\
\midrule
OLMoE-1B-7B (MoE)  & \textbf{fp32 MATCH 0.971} (quasi 1.0) & fp16 PARTIAL 0.494 \\
OLMo-2-7B (dense)  & fp32 PARTIAL 0.175 (quasi 0.0) & \textbf{fp32 MATCH 0.915} \\
OLMo-2-13B (dense) & fp32 PARTIAL 0.324 (quasi 0.0) & \textbf{fp32 MATCH 0.904} \\
OLMo-2-32B (dense) & fp16 PARTIAL 0.234 (quasi 0.0) & \textbf{fp32-tp2 MATCH 0.961} \\
\bottomrule
\end{tabular}
\caption{Overnight deployment-match sweeps ($n=12$ ifeval instances, scored against
the sampled public completions). Float32 is the highest-agreement tested precision; the winning
\emph{engine} depends on architecture. OLMoE is MoE, so vLLM's Triton fused-MoE
kernel cannot serve fp32 (shared-memory limit) and only HF fp32 produced the exact sampled OLMoE match; the
dense OLMo-2 samples matched most closely under vLLM fp32. These are probe results pending ordinary-path confirmation.}
\label{tab:fp32}
\end{table}

\par\medskip\noindent\fbox{\parbox{0.97\textwidth}{\small
\textbf{Epistemic status of the fp32 result (interpretive revision, 2026-07-15).}
This table is a \emph{configuration-discovery} result, not a reproduction result.
Read it as: \emph{among the deployment configurations tested, \code{float32}
candidates produced the highest agreement with the sampled public OLMo outputs.}
It does \emph{not} yet establish ``the historical deployment was \code{float32}''
as an identified fact, nor ``\code{float32} reproduced the official \textsc{helm}
result,'' for three reasons: (i) matching a small oracle sample does not uniquely
identify the producing deployment --- multiple latent configurations can be
observationally equivalent on $n\approx12$; (ii) the winning cells use a probe-only
\code{add\_special\_tokens=False} knob not sent by an ordinary \textsc{helm} run;
(iii) the ``confirm'' step --- land the winning config on the real \textsc{helm}
execution path, run the exact public \file{run\_spec.json}, and compare through the
standard pair-report pipeline on \emph{held-out} instances --- has not been run for
any model. The \emph{deductive} dtype-default argument above is the stronger leg and
survives these caveats (it reads the loader, not the outputs), but it too is
conditional on the unrecorded \code{transformers} version. See
Appendix~\ref{app:revision} for the consolidated epistemic status and the corrected
disagreement taxonomy.}}

\par\medskip
Two aspects make this finding sharp:

\begin{itemize}[nosep]
  \item \textbf{The matching precision is the one cell that cannot run.} vLLM's
        Triton fused-MoE kernel needs $\sim$131~KiB of shared memory per block in
        fp32 --- over the $\sim$99~KiB cap on workstation cards --- so it OOMs at
        the \emph{kernel} level (not VRAM). Every runnable local vLLM cell for
        OLMoE is reduced-precision, and none matches the sampled high-precision reference as closely. That
        \emph{fp16 beats bf16} among those cells is itself corroboration: fp16's
        finer mantissa flips fewer greedy near-ties against a high-precision
        reference, which says the reference is higher-precision than bf16 ---
        consistent with fp32. Serving fp32 requires either an H100-class card (big
        shared memory) or, for dense models, tensor-parallel fp32
        (\code{--fp32-tensor-parallel-size 2}); for the MoE, HF \code{generate()}
        at fp32 is the only route.
  \item \textbf{It reframes the disagreement.} A bf16-vs-official OLMo-2
        disagreement is a \emph{deployment-boundary artifact} (a precision the
        local engine could not easily provide), \emph{not} a reproducibility
        failure --- fp32 is the leading tested candidate for closing it, pending ordinary-path confirmation. This is precisely the
        recipe/environment-vs-reproducibility distinction of \S\ref{sec:distinction},
        applied to the most consequential single case.
\end{itemize}

\section{The residual puzzle: Hugging Face divergence for dense OLMo-2 (resolved)}
\label{sec:residual-resolved}

\evidencebox{Status at the 2026-07-27 revision}{This section is retained as written
in the 2026-07-15c consensus revision, because the reasoning it records is the
reasoning that produced the follow-up experiment. It was resolved on 2026-07-23 in
favor of the interpretation the 15c review insisted on keeping alive: a defect in
the dense-model probe path. See the resolution note at the end of the section, and
Chapter~\ref{chap:substrate-closed} for the experiment.}


Rigor demanded following up an inconsistency. The HF probe exactly matched the sampled OLMoE outputs but gave near-zero quasi-match for the three dense OLMo-2 models,
even though vLLM fp32 best matched the sampled dense-model outputs. The forensic diagnosis localized the disagreement:
tokenizing the exact ifeval prompt showed byte-identical tokens across every path
(so it is \emph{not} the prompt); yet the HF probe's first token disagreed with the
official (first-token agreement 0.42) while vLLM fp32 matched the sampled official text character-for-
character (first-token 1.00). A first-token disagreement on a byte-identical prompt
is a clean litmus: at step 0 there is no accumulated history, so a different first
token can only mean a different \emph{forward-pass} logit vector. The divergence is
a fp32 forward-pass difference, not a prompt or precision-class difference ---
attributable to knobs \textsc{helm} never records: \code{attn\_implementation},
device placement (\code{device\_map=auto} shards a fits-on-one-GPU model across
GPUs, changing fp32 reduction order), and the decode path (\textsc{helm}'s
\code{do\_sample=True, temperature=1e-7} vs the probe's greedy). The probe was
extended to sweep these forward-pass axes. \textbf{Interpretive revision: this is a
finding, not merely a broken probe.} An earlier framing dismissed the OLMo-2 HF
divergence as ``not a science problem'' (a defect in our tool, since vLLM fp32 did
match). The sharper reading, which the paper should adopt: a \emph{current}
\code{transformers.generate()} at fp32 does \emph{not} reproduce a historical result
\emph{recorded as a \code{HuggingFaceClient} (in-process \code{transformers})
deployment} --- \emph{even under byte-identical prompt tokens}. That is precisely an
\emph{unresolved systematic divergence under prompt-token equivalence} --- direct
evidence of execution-stack sensitivity, and a core datum for the thesis, not an
artifact to wave away. Its cause is not yet isolated, and two classes of explanation
remain live until one is ruled out: a genuine execution-instrument divergence
(candidates: model/tokenizer revision, \code{transformers}/\code{torch} version,
attention implementation, generation defaults, device placement, numerical precision,
CUDA behavior) \emph{and} a defect in the dense-model HF probe path itself (e.g.\
tokenizer/template handling that differs from the OLMoE path on which the probe
succeeds). The reframing rejects only the premature dismissal of the divergence, not
the possibility that our own tooling contributes to it. That vLLM fp32 happens to match is a separate fact and
does not make the HF$\to$HF gap uninteresting. The reusable insight stands and is
sharper for it: \emph{``same engine family'' (transformers$\rightarrow$transformers)
does not imply reproducible} --- device sharding, attention-implementation drift, and
sample-vs-greedy each perturb fp32 logits enough to diverge greedy decoding.

\subsection*{Resolution (2026-07-23)}

The candidate list above named the cause. When the probe was re-run with every axis
it had been leaving at a sweep default pinned to the officials' configuration ---
float32, old-template rendering (\code{add\_generation\_prompt} neutralized), and
\textsc{helm}'s own decode --- \emph{all four} forward-pass cells,
\{eager,\,\textsc{sdpa}\} $\times$ \{\code{device\_map=auto},\,single\}, reproduced
the official dense OLMo-2 completions byte-exactly on both 7B and 13B ($32/32$
instances, composite $1.000$). The near-zero quasi-match of the original probe is
therefore attributed to the probe's own configuration: the sweep's default template
rendering was the modern behavior rather than the historical one, and
\code{device\_map=auto} on that host's smaller cards sharded a model that fits on
one GPU, changing fp32 reduction order. On 96~GB cards nothing shards, and attention
implementation turns out to be a non-factor for this model at fp32.

Two things follow, and they pull in opposite directions.

\begin{itemize}[nosep]
  \item \textbf{The retraction.} The framing this section adopted --- ``a current
        \code{transformers.generate()} at fp32 does not reproduce a historical
        \code{HuggingFaceClient} result even under byte-identical prompt tokens''
        --- is \emph{withdrawn}. It does reproduce it, exactly. The 15c review was
        right to refuse to promote this to a finding while a probe artifact
        remained on the candidate list, and right to refuse the earlier dismissal
        as well; ``keep both live until one is ruled out'' was the correct ruling
        and it is what ruled one out.
  \item \textbf{What survives, stronger.} The reusable insight above is unchanged in
        substance and now has a metric attached to it. The instrument sensitivity is
        real; it lives at the \emph{engine} boundary rather than inside the
        transformers family. Same weights, same float32, same prompt, same greedy
        decode, \textsc{vllm} instead of Hugging~Face: $+0.067$ and $+0.082$ on
        published IFEval accuracy (Chapter~\ref{chap:substrate-closed}).
\end{itemize}

\noindent The methodological lesson is worth more than either. A sweep tool
\emph{defaults} every axis its caller does not pin, and an unpinned axis is a silent
experimental choice --- the same failure genus as the unrecorded parameters this
chapter is about, committed by our own instrument. The probe was read as a
faithful reproduction attempt when it was a configuration search; a configuration
search that reports its best cell without reporting which axes it swept invites
exactly this misreading.

\section{Naming the gap: the unrecorded execution substrate}
\label{sec:substrate}

The investigation crystallized into a taxonomy (\file{helm-unrecorded-deployment-
params.md}, Appendix~\ref{app:unrecorded}): \textsc{helm} serializes the
\emph{recipe} and the model's \emph{name} but almost nothing about the
\emph{execution substrate}. Empirically, of 148 HuggingFaceClient deployments in
\textsc{helm}'s \code{model\_deployments.yaml}, only 19 pin \code{torch\_dtype}
(all \code{bfloat16}) and only 7 pin a model \code{revision}. Ranked by effect on a
greedy output, the unrecorded parameters are: (Tier 1) weight revision, load
precision, quantization, software-stack versions; (Tier 2) attention implementation,
matmul/TF32 precision, hardware, device topology, batch composition; (Tier 3)
tokenizer + chat-template version, \code{add\_special\_tokens}/BOS handling,
\code{generation\_config.json} defaults, stop-sequence/detokenization semantics.
The one high-leverage lever is \emph{pinning the model revision (commit SHA)},
because the tokenizer, chat template, and generation config all live inside the
model repo --- so \code{model\_revision} + \code{torch\_dtype} +
\code{transformers\_version} together close the majority of the gap. Crucially,
these must live in \code{client\_spec.args} / \code{catalog.yaml} / a
\code{recipe\_facts} block, \emph{not} in the run-spec name or
\code{model\_deployment} identity string --- a SHA in the pairing key would break
official$\leftrightarrow$local pairing, whereas a SHA in the provenance block is
invisible to it.

A companion effort (\code{hf\_inprocess.py} + infer-stack
\code{acquire --reserve-gpus N}) built the mechanism to reproduce a
HuggingFaceClient official \emph{the same way it was produced} --- in-process
\code{transformers.generate()} at pinned fp32, on a GPU held (but not served) by a
reserve-only lease. The mechanism landed; the routing switch that would make the
default replay path choose it for a HuggingFaceClient official was scoped but left
unwired (a replay still defaults to vLLM).

\evidencebox{Status at the 2026-07-27 revision}{The routing switch is \emph{still}
unwired, and the engine question was answered without it. The
2026-07-22/23 experiment reached the same place through the deployment-match
\code{hf-probe} path --- pure \code{transformers.generate()} on one GPU, no lease,
no serving stack --- which was the deliberately lower-risk choice for unattended
overnight use (Chapter~\ref{chap:substrate-closed}). The consequence is a real gap
in what was proved: the byte-exact result is at the \emph{completion} level on 32
instances per model, not a full \textsc{helm} metric run on the in-process client.
Because identical completions imply an identical IFEval score, this is a strong
inference rather than a measured metric; wiring the switch (or a full-$n$ probe
confirm) would close it. Note that this leaves the routing switch on the open-threads
list for a different reason than before: it is no longer the only way to answer the
question, only the way to answer it at the metric level.}

%==============================================================================
\chapter{Era-Pinned Containers: Reproducing the Pre-v0.5 Corpus}
\label{chap:eras}
%==============================================================================

\emph{Commits: 2026-07-08 through 2026-07-13; the
\code{impl/era-pinned-helm-containers} branch, the \code{docker/era\_shim/}
package, and gotcha G13.}

\section{The problem: the measurement instrument itself has versions}

An internal planning estimate placed roughly \textbf{59\%} of the then-indexed
audit corpus on the classic-track \code{v0.2.4}/\code{v0.3.0} path
(\file{docs/container-execution.md}; \file{dev/journals/claude.md}). The percentage
must be regenerated from a checked-in corpus manifest with an explicit numerator and
denominator before publication. Those runs \emph{cannot} be replayed by the modern
runner, which pins \textsc{helm} 0.5.x and Python 3.11 and whose from-spec CLI
imports v0.5+ module paths that do not exist pre-v0.5. This is not a model problem;
it is a \emph{harness} problem. And it matters for reproducibility science because
the harness is the measurement instrument: the technical report had already shown
that the exact \code{pandas} version (2.0.x vs 2.2+) flips per-instance row
ordering on \code{entity\_matching}, i.e.\ the harness's own dependency versions
change the recorded instances.

The solution is to freeze the instrument at the era that produced each run: give
each era its own \emph{CPU-only} Docker image whose \textsc{helm} harness is
checked out at that era's release commit (v0.2.4 = \code{626d8609}, v0.3.0 =
\code{8ea285f7}), with era Python (3.10) and era dependency pins, while keeping
model inference \emph{out-of-process} on a modern vLLM lease. The container thus
holds only the scoring instrument fixed; the model server is shared with the
modern path.

\section{The design invariants}

\begin{itemize}[nosep]
  \item \textbf{One manifest = one era = one image = one measurement instrument.} A
        mixed-era source set is a hard error at make-manifest time
        (\code{eval\_audit/eras.py}). Era resolution keys on the path-derived
        \code{(public\_track, suite\_version)} signal.
  \item \textbf{Verbatim by-name replay.} A pre-v0.5 \code{adapter\_spec} has no
        \code{model\_deployment} field, so nothing is rewritten; routing to vLLM is
        purely by-name (the era shim registers a deployment under the exact official
        model name), and the materializer \emph{refuses} to insert a
        \code{model\_deployment} field into an era spec. Consequently
        \code{same\_deployment} resolves \code{unknown} for era pairs (both sides
        lack the field) --- the correct behavior, not a bug.
  \item \textbf{A standalone shim.} \code{docker/era\_shim/helm\_era\_shim/} never
        imports magnet or eval\_audit; it provides a flag-compatible from-spec CLI
        (\code{replay.py}) plus a backported OpenAI-compatible completions client, so
        the era image depends only on era \textsc{helm}.
  \item \textbf{An era$\leftrightarrow$image guard.} The image carries an
        \code{org.aiq.era} OCI label, checked against the manifest era at schedule
        time; a wrong pin fails loud.
\end{itemize}

\section{The v0.2.4-versus-v0.3.0 API skew}

A multi-angle review (verifying every era-API claim against the era \textsc{helm}
source via \code{git show <commit>:<path>}) found ten issues, the load-bearing ones
being cases where the shim had been written against v0.3.0 semantics and broke on
the older v0.2.4 API. Two are especially instructive:

\begin{itemize}[nosep]
  \item \textbf{Silent Together fallthrough (v0.2.4).} v0.2.4 lacks the
        register-deployments-from-path helper, so without an explicit
        \code{register\_model\_deployments\_from\_path} the deployments route to
        \code{TogetherClient} --- silently reintroducing the very hosted-API
        artifact the audit exists to eliminate.
  \item \textbf{pyhocon dot-splitting (both eras).} pyhocon path-splits credential
        keys on \code{.}, so a dotted model name like \code{pythia-6.9b} dies at the
        credential lookup. The fix writes a nested-key \code{credentials.conf} and
        puts \code{api\_key} in the client-spec args.
\end{itemize}

\section{A turnkey validation runbook and the 8~GB subject}

The era gates were rebuilt as \code{dev/era-tests/}, mirroring the \code{dev/e2e-
tests/} conventions, so one model runs end-to-end through both classic eras with
near-zero setup.

\evidencebox{Packaged-evidence status}{By July 14 the journal refers to intact
\code{era-redpajama} runs and refreshable report artifacts for both proxy eras, and
the final slide says the end-to-end script works. Those report stores live under
\path{/data}, which the supplied source archive excludes. This master therefore
records the execution as locally supported but requires the report directories,
inputs, and hashes to be archived before publication.}
 The subject was swapped from \code{eleutherai/pythia-6.9b}
($\sim$14~GB, does not fit 8~GB) to \code{together/redpajama-incite-base-3b-v1},
the smallest corpus model with a \emph{full 74-run packet at both eras}
($\sim$2.8B params, $\sim$5.6~GB fp16). A ``validation ladder'' checks the
instrument in rungs: image sanity; instrument fidelity on a pandas-sensitive
\code{entity\_matching} run (byte-for-byte instance identity, no model); an
end-to-end single run; a full packet per era; and a per-scenario HF-fetch audit.
Several first-real-run bugs (a build-time import that named the wrong module for
the era; a \code{pandas}/\code{pyarrow} dependency-resolution skew; a fidelity rung
that diffed a \file{scenario\_state.json} the public corpus never ships --- switched
to the always-present \file{display\_requests.json}) reinforced the theme that
some bugs only surface at the first genuine build.

An important taxonomy point resurfaced here: rung 2 originally \emph{failed} on
scenarios whose data is no longer fetchable (e.g.\ MATH's script loader returns 401
on the 2026 Hub). That is an \emph{environment} filter, not an instrument failure,
so the rung was changed to classify data-fetch/auth crashes as SKIP --- the same
recipe-vs-reproducibility distinction, now applied to a historical-data gate.

\section{G13: a class path that resolves in no released \textsc{helm}}
\label{sec:g13}

The deepest archaeological finding of the internship concerns the original-paper
classic models (GPT-J 6B, GPT-NeoX 20B, OPT-66B) and is written up as gotcha G13
(\S\ref{app:gotchas}). Era-replaying these fails the shim's class preflight: their
\file{run\_spec.json} names \code{helm.benchmark.basic\_metrics.BasicMetric}, which
resolves in \emph{no} released \textsc{helm} version. An investigation over a
blob-less clone of the \textsc{helm} repository established that the stored path is
a once-migrated hybrid: a bulk \code{benchmark.}$\rightarrow$\code{helm.benchmark.}
prefix rewrite (at the \code{src/helm/} rename, 2022-11-16) applied to a run-spec
whose metric class was still \emph{flat} at production time. Reversing the naive
migration triangulates the producing code to a $\sim$4-week window
(2022-07-31 to 2022-08-26) of \emph{unreleased pre-v0.1.0 \textsc{helm}} --- after
scenarios were nested but before metrics were. The class-path \emph{shape} (flat
vs subpackage) is thus a sharper origin fingerprint than release dates. Because no
released version resolves these, and building the true origin instrument is
infeasible, the fix is a self-verifying declared class-path canonicalization in the
era shim: when a flat \code{helm.benchmark.*} class fails to resolve \emph{and} its
\code{helm.benchmark.metrics.*} relocation resolves to the same leaf, remap it on
the in-memory run-spec (recorded as a declared substitution). It fires only for the
known relocation --- a genuinely wrong era pin still fails loudly.

A related provenance fact was pinned along the way (memory
\emph{classic-officials-carried-forward}): the gptj/neox/opt/redpajama official
runs are \emph{byte-identical} across v0.2.4 and v0.3.0 --- carried forward, not
re-run --- so ``two eras agree'' is one public number measured by two instruments,
not two independent confirmations.

%==============================================================================
\chapter{New Model Families: GPT-OSS and Qwen}
\label{chap:newmodels}
%==============================================================================

\emph{Commits: 2026-07-11 through 2026-07-14; from-spec runbooks for
\code{openai/gpt-oss-20b} and a combined Qwen fan-out.}

With the from-spec, containerized, leased pipeline mature, two further model
families were brought in as reproduction subjects --- and as a test of whether the
machinery generalizes beyond OLMo.

\section{GPT-OSS-20B and the null-content crash}

The public GPT-OSS-20B rows were served via \code{together/gpt-oss-20b}, a chat
(harmony) client, so the faithful replay serves \emph{chat}, not the frozen older
preset's completions workaround. The design tension: GPT-OSS is a reasoning model,
and a chat turn can return \code{message.content=null} (reasoning-only,
\code{finish\_reason=length}), which un-patched \textsc{helm} crashes on. The root
cause was pinned precisely (memory \emph{gpt-oss-null-content-root-cause}): the
crash is a \emph{null-versus-empty-string serving difference} --- Together returns
\code{""}, vLLM returns \code{null} --- and normalizing
\code{null}$\rightarrow$\code{""} is faithful (\code{max\_tokens} and window were
ruled out as causes). The fix makes the null-safe chat clients normalize
\code{content:null}$\rightarrow$\code{""}, leaving \textsc{helm} untouched.
Discovery confirmed all four target run directories resolve 1:1 against the
84{,}966-run corpus with 0 ambiguities.

\evidencebox{Store status (collaborator-verified, not packaged)}{A live-store
review reports that the \code{gpt-oss-20b-from-spec} \emph{reports} were regenerated
on Jul 14 at 13:47 after the drift-plot fixes. This supports report freshness, but a
report timestamp alone does not establish that the underlying model outputs and
comparison inputs were generated by the final accepted execution path. Before these
numbers are cited, preserve and hash the store, frozen run specs, public targets,
model/deployment manifests, code commit, and comparison configuration, then run the
acceptance checks in the canonical store-status ledger.}

\finding{GPT-OSS-20B is a promising positive local case.} The headline aggregate score-drift plot
(Figure~\ref{fig:gptoss}) shows local and public headline scores agreeing closely
across all four benchmarks --- \code{bbq} $0.967/0.963$, \code{gpqa} $0.594/0.610$,
\code{ifeval} $0.732/0.754$, \code{mmlu\_pro} $0.740/0.748$ (public/local) --- with
squared errors bounded by $\sim\!4\times10^{-4}$, more than an order of magnitude
($\approx\!30\times$) tighter than the OLMo instruct drift on \code{ifeval}
($\sim\!1$--$1.6\times10^{-2}$). This is promising evidence that family-specific deployment and artifact semantics matter. It becomes a publication result only after the underlying run and comparison provenance are preserved, hashed, and audited; report regeneration alone is not a sufficient acceptance criterion.

\begin{figure}[H]
\centering
\fbox{\parbox{0.92\textwidth}{\centering\small\itshape
[Figure: ``Headline Aggregate Score Drift'' heatmap for
\code{experiment\_name=gpt-oss-20b-from-spec}. One model row
(\code{openai/gpt-oss-20b}) $\times$ four benchmark columns (bbq, gpqa, ifeval,
mmlu\_pro), each cell annotated \code{P <public> / L <local>}, colored by
$(\text{local}-\text{public})^2$. Source: 07-14 status deck.]}}
\caption{The locally reported GPT-OSS-20B store closely matches the sampled public headline metrics (squared error
$\le 4\times10^{-4}$). The rendered artifact is
\file{aggregate\_score\_diff\_headline.png} produced by
\code{build\_reports\_summary}.}
\label{fig:gptoss}
\end{figure}

\section{Qwen: generalizing the combined fan-out}

\evidencebox{Workflow status}{The combined Qwen presets, serving fan-out, discovery,
and report plumbing were implemented. The packaged source archive does not contain a
fresh, fully validated final exact-path report store for the complete eight-model
roster. Treat model-level performance claims as pending until the store is regenerated
and its raw run directories are preserved.}


A combined exact-spec workflow was prepared for the eight public \textsc{helm} Qwen text models as a single
\code{qwen-combined} multi-deployment from-spec fan-out --- the
\code{allenai-olmo-combined} analogue. The OLMo-specific ``combined preset''
helpers were generalized into \code{\_combined\_run\_entries} and
\code{\_build\_combined\_preset} so both families share one code path (safe to
generalize precisely because the OLMo combined tests pin the old output; regenerate
and diff the entry count).

Three decisions carried over as reusable rules:

\begin{itemize}[nosep]
  \item \textbf{Membership rule (the olmo-7b lesson).} Model size is \emph{never} a
        reason to exclude a member from the fan-out --- large models serialize under
        the lease while small ones co-host; the \emph{only} thing that forces a
        member split is discovery \emph{ambiguity} under the shared
        \code{precomputed\_root}. Propose all members and let the freeze preflight
        decide.
  \item \textbf{The whitelist is exactly reconstructible.} Filtering
        \code{official\_public\_index.csv} to classic-core + capabilities lands on
        the plan's per-model counts to the row ($85\times5 + 86 + 132\times2 = 775$),
        with two non-obvious exclusions (\code{banking77}, \code{bigcodebench}). All
        775 rows resolve with 0 ambiguities, so no member splits.
  \item \textbf{Protocol is answerable, not ``to confirm.''} \textsc{helm}'s own
        \code{model\_deployments.yaml} (read from the installed wheel) says base
        Qwen1.5 = \code{TogetherClient} (completions) and the chat/turbo models =
        \code{TogetherChatClient} (chat) --- so the serving protocol was resolved
        from the authoritative source, not guessed.
\end{itemize}

\section{The Qwen turbo window bug}
\label{sec:qwenwindow}

The last dated event of this chronology is a clean recurrence of the OLMo
chat-template-window interaction (Chapter~\ref{chap:olmo}), now on the Qwen turbo
endpoints. An \code{ifeval} smoke run on
\code{qwen/qwen2.5-72b-instruct-turbo} died with a
\code{ContextWindowExceededError}: the endpoint served \code{max\_model\_len:
4096}, but the two \code{*-instruct-turbo} endpoints are the only Qwen members that
replay the capabilities whitelist (\code{ifeval}, \code{mmlu\_pro}), whose official
\file{run\_spec.json} sets \code{max\_tokens=4096}. \textsc{helm}'s window service
truncates the \emph{prompt} but never clamps \code{max\_tokens}, so 4096 output
tokens fill the entire 4096 window and vLLM 400s with no room for a single prompt
token. The official Together serve ran a \code{max\_sequence\_length} of 128k, so
the overflow never surfaced upstream. The scoped fix raises those two endpoints to
\code{max\_model\_len: 8192} (4096 output + 4096 prompt headroom). This is a serving
configuration bug, not a reproducibility failure --- and it is a fitting last
example of the internship's recurring theme: the recipe under-specifies the
serving environment, and faithfully replaying the recipe surfaces exactly where
that under-specification bites.

%==============================================================================
\chapter{Reporting and Analysis Infrastructure}
\label{chap:reporting}
%==============================================================================

\emph{Commits throughout, concentrated 2026-07-08 through 2026-07-14. The reporting
layer is where reproducibility becomes legible, and much of the final fortnight was
spent making the aggregate story honest and comparable.}

\section{The pipeline the reports run on}

The analysis pipeline is read-only over audit results that already exist on disk
(``no model is run; no benchmark is downloaded''). Its stages are: (1) convert both
public and local runs to the \eeeshort artifact format
(\code{eval-audit-prepare-eee}); (2) compose a \emph{virtual experiment} from a
YAML manifest (\code{eval-audit-build-virtual-experiment}), applying a scope and
computing the coverage funnel; (3) per-packet core analysis
(\code{eval-audit-analyze-experiment}) through \code{NormalizedDiff}; (4) aggregate
publication (\code{build\_reports\_summary}). A \emph{virtual experiment} is the
organizing abstraction --- a declarative slice over existing audit data --- and the
\emph{packet} is the planner's canonical comparison unit (a bundle of normalized
components plus the comparisons to render between them).

\subsection{The three-level coverage funnel}
\label{sec:funnel}

The coverage funnel is the quantitative backbone of any ``fraction reproducible''
claim and reports three nested levels (Table~\ref{tab:funnel}). The reason all
three exist is gotcha G1: the raw \code{run\_spec\_hash} yields \emph{zero} matches
across \textsc{helm} releases even for semantically identical recipes, because
newer \textsc{helm} populates \code{adapter\_spec} fields older \textsc{helm} left
absent. The \emph{canonical} recipe hash collapses those known schema-evolution
fields before hashing; the gap between raw and canonical match is
\textsc{helm}-version churn, and the gap between canonical and logical match is
real recipe drift.

\begin{table}[H]
\centering
\small
\begin{tabularx}{\textwidth}{lY}
\toprule
\textbf{Level} & \textbf{Meaning} \\
\midrule
logical           & same scenario + model + augmentation (order-insensitive canonical key) \\
recipe-canonical  & + same scenario spec, prompt, decoding, \code{max\_train\_instances}, after schema-collapsing the run spec \\
recipe-identical  & byte-for-byte \code{run\_spec\_hash} match \\
\bottomrule
\end{tabularx}
\caption{The three-level coverage funnel. The report shows all three so that
\textsc{helm}-version churn (raw$\to$canonical) is separated from genuine recipe
drift (canonical$\to$logical).}
\label{tab:funnel}
\end{table}

\subsection{Sankey diagrams and the failure taxonomy}

Two Sankey diagrams render the funnel: Sankey~A (Universe$\to$Scope) narrows the
$13$k+ discovered runs to the in-scope set through ordered gates (structural,
metadata, open-weight, tag, deployment, size, manifest scope); Sankey~B
(Scope$\to$Reproduced$\to$Analyzed) carries in-scope rows through logical match,
recipe-canonical match, analyzed packet, and finally an agreement bucket. Crucially,
the failure-taxonomy colors and categories were made \emph{honest}: environment/
recipe exclusions and controlled-equivalence outcomes are distinct categories, so a
gated-model exclusion never reads as a reproducibility failure. This is the visual
realization of \S\ref{sec:distinction}.

\section{The aggregate score-drift heatmaps}

The final fortnight built the plots that summarize reproducibility at a glance ---
and, tellingly, most of the work was in making them \emph{correct and comparable}
rather than in drawing them.

\begin{itemize}
  \item \textbf{Per-metric aggregate-score-difference heatmap.} Instead of
        instance-level agreement, this shows, per core metric, the difference
        between the reproduced (local) and official (public) \emph{aggregate}
        scores; color is the signed difference on a symmetric diverging colormap
        (so $0$ is the white ``no drift'' midpoint), each cell annotated
        \code{P <public> / L <local>}. A deliberate data-source decision: read the
        signed means from the already-written run-level CSV
        (\code{core\_runlevel\_table.csv}) rather than re-serializing every packet's
        JSON --- it works on existing outputs with no re-render.
  \item \textbf{Headline holistic heatmap.} One plot showing every
        model$\times$benchmark pair, each benchmark using its \emph{headline}
        metric. \code{eval\_audit} had no primary-metric concept, but \textsc{helm}
        does (\code{run\_groups[].environment.main\_name} in the schema), so a
        curated \code{HEADLINE\_METRIC\_BY\_BENCHMARK} map (mmlu$\to$exact\_match,
        narrativeqa$\to$f1\_score, wmt\_14$\to$bleu\_4, \ldots) is used \emph{iff}
        that metric is actually present, falling back to a priority list otherwise;
        the plot names the actually-used metric on each axis tick, so the choice is
        transparent. The final version colors by squared error
        $(\text{local}-\text{public})^2$ (non-negative, emphasizes big drifts) and
        orients models-left / benchmarks-bottom to match the coverage matrix.
        Figures~\ref{fig:gptoss} and~\ref{fig:olmoheatmap} are instances.
\end{itemize}

\begin{figure}[H]
\centering
\fbox{\parbox{0.95\textwidth}{\centering\small\itshape
[Figure: ``Headline Aggregate Score Drift'' heatmap for
\code{experiment\_name=olmo-models-combined}. Six model rows (olmo-1.7-7b,
olmo-2-\{7b,13b,32b\}-instruct, olmo-7b, olmoe-1b-7b-instruct) $\times$ $\sim$12
benchmark columns. The base \code{olmo-7b} row reproduces the classic benchmarks
near-exactly (GSM $0.036/0.036$, MMLU $0.295/0.287$, narrative\_qa $0.597/0.595$,
wmt\_14 $0.097/0.098$); the instruct models show large drift on \code{ifeval}
(e.g.\ olmoe $0.628/0.731$; the darkest cells) and moderate drift on
\code{bbq}/\code{gpqa}/\code{mmlu\_pro}. The $\ddagger$ marks instance-level
fallback. Source: 07-14 status deck.]}}
\caption{The OLMo reproduction at a glance. Base-model classic benchmarks reproduce
near-exactly; instruct-model \code{ifeval} shows the largest drift --- consistent
with the chat-template-version (\code{add\_generation\_prompt}) and fp32-precision
findings of Chapter~\ref{chap:deploymatch} being most consequential on long-form,
instruct evaluations.}
\label{fig:olmoheatmap}
\end{figure}

\section{Correctness fixes that changed the numbers}

Several late fixes are worth recording because each one is a way the aggregate story
could have been \emph{silently wrong} --- and each is a lesson about honest
reporting:

\begin{itemize}[nosep]
  \item \textbf{Instance-level fallback for CoT benchmarks.} \code{ifeval},
        \code{gpqa}, and \code{mmlu\_pro} are chain-of-thought scenarios whose core
        metrics are \emph{instance-only} (\code{ifeval\_strict\_accuracy},
        \code{chain\_of\_thought\_correctness}); \textsc{helm} emits no run-level
        aggregate, so those benchmarks were silently \emph{absent} from the drift
        plot. The fix widens the data contract: fall back to the mean of the
        per-instance $0/1$ scores (which \emph{is} the accuracy), tagged with a
        $\ddagger$ so the provenance is visible. Widening the contract beats hacking
        the plot.
  \item \textbf{Metric-key granularity trap.} Run-level cells are keyed by full
        \textsc{helm} stat descriptions (``\code{f1\_score test on narrativeqa}'')
        while instance-fallback cells are keyed by bare ids
        (``\code{ifeval\_strict\_accuracy}''), so the curated headline map (bare
        names) silently failed to match run-level cells and every run-level
        benchmark fell to alphabetical-first. The fix normalizes to the metric
        \emph{family} at the matching boundary and selects the representative on the
        benchmark's \textsc{helm} \code{main\_split} (test vs valid differ, e.g.\
        narrative\_qa f1 $0.597$ test vs $0.661$ valid).
  \item \textbf{Stale-local dedupe.} An \code{olmo-7b}/mmlu cell showed a spurious
        $\sim\!0.15$ aggregate gap (public $0.295$ vs local $0.144$) despite
        near-exact instance agreement. The cause was not drift: every olmo-7b
        subject directory carried \emph{two} \code{official\_vs\_local} pairs --- a
        fresh run plus an un-demoted stale $0.0$ attempt --- so micro-averaging
        halved the local mean ($124$ rows instead of $62$). The fix restricts the
        accumulator to the canonical (first) comparison per report, with a warning
        per dropped attempt so a future genuine multi-local case surfaces instead of
        vanishing.
\end{itemize}

\section{The shared-gateway route registry}

A final infrastructure fix addressed a real operational failure: multiple runbooks
(olmo, qwen, gpt-oss, classic-together) share one standing \code{infer-stack} stack
but ship disjoint catalogs, and a converge under one catalog stripped another
runbook's still-live routes from the LiteLLM gateway --- presenting as a
\code{400 Invalid model name} and (because the minutes-long readiness wait silently
never passes once a route is stripped) as the operator's known ``failed to acquire
lease'' events. The fix persists an append-only \emph{route registry} in shared
state; every converge merges the invoking catalog with all live deployments and
renders the gateway route table from the \emph{whole} registry, so the rendered
config is a function of accumulated shared state rather than of which runbook
invoked converge --- byte-stable renders, so the gateway is never recreated. The
reusable insight: \emph{when a shared artifact is rebuilt by multiple writers,
render it from accumulated shared state, never from the invoking writer's
worldview.}

%==============================================================================
\chapter{Codebase Stewardship}
\label{chap:stewardship}
%==============================================================================

\emph{Commits: 2026-07-02, 2026-07-06, 2026-07-12. Three audit/simplification
passes interleaved with the research work.}

Reproducibility is a property of software, so the health of the codebase is part of
the scientific claim: a reviewer who cannot read the pipeline cannot trust its
verdicts. Three passes are worth recording as evidence of the ``code a researcher
can trust'' standard.

\begin{description}[style=nextline,leftmargin=1.25em]
  \item[Correctness audit (2026-07-02/03).] A full codebase audit produced a
        phased fix plan (P0/P1/P2 findings) implemented over $\sim$35 commits: e.g.\
        matching/joining on canonical keys rather than raw strings; correcting a
        perturbed/unperturbed agreement split; filtering non-finite scores at row
        construction; failure-taxonomy honesty (colors, categories, ordering);
        content-addressing generated \code{model\_deployments.yaml} filenames;
        tightening the HF-token permission window. Direct dependencies were declared
        and \code{transformers} pinned below 5 (a version whose \code{is\_offline\_
        mode} change breaks the \textsc{helm}$\to$\eeeshort rebuild).
  \item[Simplicity audit (2026-07-06).] An explicitly orthogonal ``reduce bloat''
        pass: $\sim$10{,}600 lines deleted (dead POC/oneoff trees, generated
        manifests, shims), \code{helm/diff.py} halved ($1{,}844\to844$),
        \code{PRESET\_CONFIGS} moved from a $1{,}154$-line dict to a YAML resource,
        the planning-doc set consolidated ($22\to11$ live docs), and the legacy
        agreement half of the comparison core finally retired (R-2), completing the
        migration onto \code{NormalizedDiff}. The methodological gem: pure-relocation
        refactors of plot-emitting code were proven correct by first
        \emph{characterizing the noise floor} (render \code{HEAD} twice, diff, and
        treat that file set as the only acceptable before/after diff --- necessary
        because plotly emits random div UUIDs).
  \item[Fourth simplification pass (2026-07-12).] A post-era ``finish what's in
        flight'' pass: Batches 1--3 landed in $22$ commits with the full
        \code{--run-slow} suite green ($638$ passed); \code{helm/} shrank
        $3{,}336\to\sim1{,}780$, \code{build\_reports\_summary.py} finished its split
        ($1{,}715\to\sim330$). A recurring review discipline shows through: audit-
        agent claims about \emph{why} code exists (``kept for tests'') were the ones
        to re-verify by grep --- liveness claims were reliable, purpose claims less
        so.
\end{description}

Two standing engineering disciplines recur across all three passes and are worth
lifting into the paper's reproducibility-methods note: (1) capture a
characterization baseline \emph{before} any behavior-changing refactor and gate on
\code{atol}$=10^{-9}$ equivalence; and (2) submodule hygiene --- land changes in the
submodule branch first, bump the gitlink as a deliberate separate commit, never
auto-commit a dirty gitlink.

%==============================================================================
\chapter{Thread B: Efficient Benchmarking}
\label{chap:effbench}
%==============================================================================

\emph{Repositories: \file{mrmr\_eval} and \file{dkps} (sibling repos, outside
\file{eval\_audit}); active $\sim$2026-06-01 through 2026-06-24; status decks 06-02,
06-09, 06-16, 06-23. Paused thereafter in favor of the HELM-reproduction thread.}

Running in parallel with the reproduction work for the first month was an
exploratory statistics thread on \emph{efficient benchmarking}: given that
evaluating a model on a full benchmark is expensive, can a small, well-chosen
\emph{coreset} of questions plus a regression predict the model's full-benchmark
score, and does a Data-Kernel Perspective-Space (\textsc{dkps}) representation help?
The experiment code lived in two sibling repositories forked from the relevant
prior work; the LoRA weight-space idea from the literature review
(Chapter~\ref{chap:litreview}) remained a proposal and produced no code.

\section{The two framings being compared}

\subsection{Feature selection + regression (Bowyer et al.)}

Following \emph{Efficient Benchmarking Is Just Feature Selection and Multiple
Regression}, the setup is a binary score matrix (models $\times$ questions). Each
question is a feature vector in $\mathbb{R}^M$ ($M$ source-model scores) and the
target is each model's overall accuracy. One \emph{selects} an informative coreset
of questions --- primarily with \textbf{mRMR} (Minimum-Redundancy-Maximum-
Relevance, in MIQ = relevance/redundancy and MID = relevance$-$redundancy variants)
--- and \emph{predicts} the full score with ridge or polynomial kernel-ridge
regression. The \code{benchpred} registry implements $\sim$100 method keys spanning
mRMR (ridge/kernel heads), IRT (\code{pirt}, \code{gpirt} = Gaussian-process IRT),
anchor points, Lasso, and random-sampling/search baselines, each optionally
``backfilled'' (append \code{+}) by re-fitting its coreset with ridge/kernel-ridge.

\subsection{DKPS: a model-centric, transductive predictor}

\textsc{dkps} (Data-Kernel Perspective-Space) is the model-centric alternative.
Given a pool of source models with known scores and a target model, all identified
only by their \emph{raw responses} to an aligned set of $N$ queries, the method:
(1) embeds every $(\text{model}, \text{query})$ response (via a text-embedding API,
or a one-hot encoding for multiple-choice tasks); (2) forms an inter-model
``data kernel'' --- a models$\times$models Euclidean distance matrix on the flattened
per-model response embeddings, normalized by $\sqrt{N}$; (3) computes the
\emph{perspective space} by classical MDS on that distance matrix (default 8
components); and (4) fits a linear regression from source coordinates to source
scalar scores and reads off the target's predicted score. It is \emph{transductive}
--- the target's responses participate in building the MDS space, so the space is
re-fit per prediction --- and it requires every model to answer the exact same
aligned queries under the same embedder.

\section{Experiment design}

The efficient-benchmarking sweeps concentrated on the HELM tasks that carry raw
responses (so \textsc{dkps} is applicable): \code{legalbench}, \code{math},
\code{med\_qa}, \code{wmt\_14}, plus \code{arc\_challenge}, \code{ifeval},
\code{commonsense}. The knobs swept were: number of source (reference) models
$\{20, 30, 40, 50\}$ (sometimes 60); coreset size $\{1,2,3,4,5,10,15\}\%$; three
seeds per cell. The reported metrics were MAE, RMSE, Kendall $\tau$, and Spearman
$\rho$ (Figure~\ref{fig:mrmr}).

\begin{figure}[H]
\centering
\fbox{\parbox{0.95\textwidth}{\centering\small\itshape
[Figure: \code{math} --- \code{mrmr\_eval} tutorial. A $2\times4$ grid of
RMSE(\%), MAE(\%), Kendall $\tau$, Spearman $\rho$; top row versus coreset size
(5/10/15\%), bottom row versus number of source models (20--50). Methods overlaid:
mRMR MIQ (k=5), gp-IRT, Search+, random\_sampling\_and\_learn, Random, AnchorPts,
Lasso. mRMR/kernel-anchor methods lead on RMSE/MAE; Lasso is the clear worst.
Source: 06-16 status deck.]}}
\caption{Representative efficient-benchmarking result on the \code{math} task.
Coreset-selection methods (mRMR, kernel anchor points) substantially outperform
Lasso and improve with coreset size; \textsc{dkps}-family predictors are competitive
only in specific regimes (many source models, tiny coresets).}
\label{fig:mrmr}
\end{figure}

\section{The DKPS+mRMR active-selection method}

The internship's own methodological contribution to this thread was an
\emph{adaptive} coreset selector, \code{dkps\_mrmr} (the realization of the
``kNN-mRMR online coreset selection'' proposal), building the coreset iteratively
and \emph{per target model}: seed with one global mRMR pick; then repeatedly
(a) \textsc{dkps}-embed source $\cup\ \{\text{target}\}$ on the current coreset,
(b) take the $k$ source models nearest the target in \textsc{dkps} space, and
(c) restrict the score matrix to those $k$ rows and take one greedy mRMR step to
append the next question; finally predict with a standard \textsc{dkps} transductive
regression. It subclasses both \code{MRMRPred} and \code{DKPSPred} to reuse the
mutual-information machinery and the embedding/regression, and required a new runner
capability so \code{predict} receives the full response matrix rather than a
pre-sliced coreset.

\section{Findings}

The concrete results (chiefly the \code{med\_qa} report) are candid, and their
honesty is itself worth preserving:

\begin{itemize}[nosep]
  \item \textbf{\textsc{dkps}-family methods are not general winners.} At $\sim$20
        source models and $\sim$50 queries, the feature-selection-paper methods
        (kernel anchor points, Search+, kernel-mRMR) outperform \textsc{dkps} and
        \code{dkps\_mrmr}; on the 40-source-model \code{med\_qa} row,
        \code{dkps\_mrmr} wins no coreset-size column.
  \item \textbf{\textsc{dkps} scales with source models where others do not.}
        \textsc{dkps} accuracy improves as the number of source models grows, which
        is not necessarily true of the other methods --- and this is where
        \code{dkps\_mrmr} shines: at a $1\%$ coreset its \code{med\_qa} MAE improves
        sharply from $9.06$ to $5.19$ as sources go $20\to50$, making it the
        \emph{single best method at 50 source models} in that corner.
  \item \textbf{\code{dkps\_mrmr} plateaus and is worst at large budgets.} It peaks
        early (MAE $\sim\!4.6\%$, $\rho\approx0.93$) and, unlike every other method,
        does not keep improving to $2.0$--$2.9\%$ MAE at $10$--$15\%$ coresets.
  \item \textbf{Cost and representation caveats.} Because the coreset is rebuilt per
        target, one trial costs $\sim$3~s (1\%) to $\sim$2.5~min (15\%), versus
        sub-second for plain \textsc{dkps}. mRMR is \emph{extremely slow} on
        continuous-score benchmarks (\code{wmt\_14}). And a one-hot response
        encoding (used when text embeddings are unavailable) exposes very limited
        information about a model compared to a text encoding --- for \code{med\_qa},
        with 5 choices but imperfect answer formatting, the one-hot dimension
        exceeds 5 and the representation is weak.
\end{itemize}

\section{What stayed a proposal}

Three lines from the ``directions'' slides were never implemented and are recorded
as open future work: \emph{variance-based early stopping} for adaptive
benchmarking; \emph{ensemble} predictors mixing feature-centric (AEA, Scales++) and
model-centric (\textsc{dkps}, Fluid) methods, e.g.\ re-weighting by \textsc{dkps}
distance from a reference model; and the entire \emph{LoRA weight-space} (W2T)
program (SVD canonicalization $\to$ transformer embedding $\to$ attribute
classification / performance prediction / adapter retrieval, and the exploratory
idea of learning a distribution over LoRA weights). This thread effectively
concluded at the end of June 2026 when attention returned fully to HELM
reproduction.

%==============================================================================
\chapter{Extending to Net-New Models: The Qwen3.5 Arc}
\label{chap:qwen35}
%==============================================================================

\emph{2026-07-15 through 2026-07-17. Led by Jon Crall
(\S\ref{sec:rev-authorship}); journal entries of 07-15 17:40 (with ten addenda)
and 07-17 09:40.}

Everything before this chapter reproduces \emph{someone else's} run. This chapter is
the first work in the record that \emph{computes} a run of our own: Qwen3.5-9B-Base,
a model with no public \textsc{helm} entry, evaluated on the same scenario roster as
the Qwen reproductions so that new-model numbers land in the same recipe frame as
reproduced ones. The runbooks label themselves \emph{compute instead of reproduce}
for exactly this reason. Three things came out of it that matter beyond the model:
an extension mechanism, two new gotchas, and a self-critique.

\section{Registering a model \textsc{helm} has never seen}

The first attempt registered \code{qwen/qwen3.5-9b-base} by editing the vendored
\textsc{helm} fork's \code{model\_metadata.yaml} and \code{tokenizer\_configs.yaml}.
That is the pattern the earlier Qwen planning document had blessed, and it was
\emph{reverted the same session} on the correct objection that it modifies the
instrument. \textsc{helm} already reads \code{model\_metadata.yaml},
\code{tokenizer\_configs.yaml}, and \code{model\_deployments.yaml} from its
\code{--local-path} registry directory; our stack simply never passed the first two
through. Building that passthrough (in the MAGNeT backend's
\code{prepare\_local\_helm\_config}, then in the manifest, bridge, and container
mounts, and finally in the bundle exporter's alias assertion) reduced ``add a
net-new model'' to \textbf{one preset block, two sidecar YAMLs, and one catalog
endpoint --- with no \textsc{helm} edit and no image rebuild}.

This is the extension seam a reproducibility instrument should have. The general
form of the lesson: when the rule is ``do not modify the instrument,'' look for the
instrument's own extension mechanism before building a shim around it. The prior
belief that ``we have done this before'' turned out to point at
\code{bundle\_export}'s \code{\_assert\_helm\_aliases\_exist}, which does the
\emph{opposite} --- it \emph{requires} registration in the vendored YAMLs. Every
local model until now had been upstream-registered, so only the deployments half of
the plumbing had ever been needed.

\section{A 0.0 cell that was not model weakness}

The first smoke run produced \code{mmlu:anatomy} \code{exact\_match}$=1.0$ (4/4) and
\code{boolq} \code{exact\_match}$=0.0$ (0/5). The instance-level forensics are the
point: the raw vLLM cache showed the model's first generated token was
\code{"\textbackslash n\textbackslash n"} (logprob $-0.28$), \textsc{helm}'s
canonical \code{stop=['\textbackslash n']} matched inside it, and every completion
truncated to the empty string. The prompt was well-formed with five inline
\code{Answer: Yes/No} shots; the tokenizer was exonerated by a separate probe.

Two hypotheses were written down before testing --- H1, the base model answers
paragraph-style (\code{Answer:\textbackslash n\textbackslash nYes}) and the recipe's
stop sequence zeroes it; H2, fp16-on-Turing numerics distort the hybrid attention
kernels --- and a discriminating probe replayed the canonical recipe alongside an
unstopped variant. \textbf{Verdict: H1.} The content behind the leading newlines was
correct on every case. The $0/5$ is a \emph{recipe-format artifact}, not model
failure: in this document's vocabulary, a deployment-boundary interaction caught by
instance-level forensics --- the project's central taxonomy operating on a
\emph{net-new} model rather than a historical one.

The fix is instructive because of how it was constrained. A newline-tolerant
completions client relaxes the stop server-side, strips leading newlines, then
re-applies the \emph{original} stops client-side; it fires only when a request
carries a \code{"\textbackslash n"} stop, so multiple-choice shapes pass through
byte-identically. Crucially, the exporter \emph{refuses} to enable it unless the
deployment name carries an \code{nlstrip} marker --- the substitution must be
visible in \code{run\_spec.json}, which is the artifact of record. Under it,
\code{boolq} went $0.0 \rightarrow 0.8$ and \code{mmlu:anatomy} stayed $1.0$. The
single remaining miss was a real finding: the base model emitted a \code{<think>}
tag on a plain completions prompt --- reasoning-format contamination in pretraining,
measurable per task once the recipe artifact is removed.

\section{Two new gotchas}

\paragraph{G14 --- a compiled-artifact cache keyed narrower than the config space.}
Disabling the vision modalities on a text-only serve changed the serving
configuration but \emph{not} vLLM's compile-cache hash, so the engine reloaded an
ahead-of-time-compiled graph traced under multimodal-enabled inputs and fed it
\code{None} where a tensor had been baked in. We got the loud failure
(\code{AttributeError} inside the compiled graph). The quiet one is the hazard worth
recording in a reproducibility document: \emph{a serving-config change that silently
reuses compiled artifacts from a different config can produce different numerics
with no error at all}, and the compiled-graph provenance is invisible to
\code{run\_spec.json}, invisible to vLLM's own cache key, and --- before this ---
invisible to our serving records. Mitigation: key the cache mount on a hash of
image, command, and generation environment, so any serve-argument change starts a
fresh cache while identical configs keep full reuse. The epistemics are also worth
keeping: this was observable \emph{only} because a config change happened between
two runs sharing a cache. A fresh machine would never see it; a long-running grid
would. Cross-machine reproduction runs are implicitly cache-freshness controls.

\paragraph{G15 --- a from-spec run entry and a compute run entry are not
interchangeable.} The first full grid was lifted verbatim from a reproduced Qwen1.5
grid with the model token swapped, and 57 of 86 jobs failed with
\code{get\_mmlu\_spec() got an unexpected keyword argument 'groups'}. In from-spec
mode the \code{eval\_split=} and \code{groups=} tokens are official-run-\emph{name}
matcher metadata; \textsc{helm} replays a frozen \code{run\_spec.json} and never
calls the constructor. In compute mode every run entry \emph{is} a constructor call.
A run key is a lossy name; lifting one from a replay grid into a compute grid
smuggles lookup-only kwargs into a function that rejects them. \textsc{helm} was not
modified --- the bug was entirely in the authored grid, and the rule held.

A further nine failures (7 \code{math}, 2 \code{natural\_qa}) were dataset-access
barriers: a \code{DatasetNotFoundError} on \code{hendrycks/competition\_math} and an
HTTP 403 on the \code{natural\_qa} source download. These are \emph{filtering
reasons}, not results, and were dropped from the grid with a note to restore the
entries if the datasets return. The corrected grid is 72 entries.

\section{The self-critique: we author what we criticize}
\label{sec:freeze-critique}

The compute runbooks author \code{run\_entries} as literal \textsc{helm} run-key
\textsc{dsl} strings in a preset file and hand them to \textsc{helm}'s expander at
execution time. That is the fragile artifact this project criticizes \textsc{helm}
for. The question was put directly: do we criticize \textsc{helm} for mutable
run-key names and expanders, then fall back to that pattern for our own runs?

The honest answer is yes, but the strong form overclaims and the useful form is
narrower. Three parts of the critique hold: a \emph{provenance asymmetry} ---
byte-exact frozen-spec replay for others' runs, a mutable name for our own; we
\emph{manufacture new fragility}, because the \code{run\_spec.json} a compute run
emits is a derivative of our pinned expander, so a future reader inherits the
version coupling; and it has already bitten, since G13 is precisely a stored run key
whose class path no released expander resolves. What blunts it: for a de-novo model
there is no prior spec to replay, so authoring is necessary rather than backsliding,
and the fragility is specifically \emph{cross-version} --- under a single pinned
\textsc{helm} build the key$\rightarrow$spec map is deterministic, so the risk is
deferred onto a future reader rather than incurred now.

The resolution reinforces the thesis instead of undermining it. A run key is a lossy
\emph{name}; the \code{RunSpec} is ground truth, and \textsc{helm}'s original sin is
treating the name as the durable handle and regenerating on demand. Our from-spec
discipline already fixes that for historical runs; the gap is that it does not yet
extend to our own freshly computed ones. The one-line rule is \textbf{expand once at
authoring, then freeze}: the expander touches each compute run exactly once, at
birth, and its content-addressed \code{run\_spec.json} becomes the archived handle.
\file{docs/planning/compute-run-spec-freeze-plan.md} specifies this, including a
re-expand-and-diff drift guard that would catch G13-class breakage. \textbf{It is
proposed, not implemented}; until it lands the presets still keep the key string as
the stored source of truth, so the discipline is aspirational and must be described
that way. Its blocking unknown is whether \textsc{helm} exposes
run-key$\rightarrow$\code{RunSpec} expansion without running inference.

\section{Status and what is not claimed}

The mechanism stack is proven end-to-end on GPUs: registry sidecars, world-pinned
leasing, per-config compile caches, modality-limited serving, the declared-
substitution client, and verified artifacts. What exists is a passing smoke run, a
corrected 72-entry core grid, three smaller sibling models
(0.8B/2B/4B) wired through a VRAM-aware placement planner that selects GPUs by
declared minimum VRAM rather than by index pinning, and a local-only virtual
experiment configuration. \textbf{No Qwen3.5 report has been produced or reviewed},
and the comparison it is designed for --- Qwen3.5-Base beside the \emph{reproduced}
Qwen1.5/2/2.5 numbers --- has not been run. Note also that the framing must stay
precise when it is: base-versus-instruct, not ``3.5 versus 2.5,'' because the clean
comparison is against the base Qwen1.5 officials.

%==============================================================================
\chapter{The Open-Judge Experiment}
\label{chap:openjudge}
%==============================================================================

\emph{2026-07-17 through 2026-07-21. Led by Jon Crall; plan at
\file{docs/planning/open-judge-plan.md} (drafted 07-17, reviewed and revised the
same day); journal entries of 07-17 through 07-21.}

\S\ref{sec:judge} recorded the judge-identity \emph{design}: judge models are
hard-coded inside \textsc{helm}'s annotator classes as a function of \textsc{helm}
version, so officials need a curated registry, and the \code{same\_judge}
comparability fact is emitted only for declared substitutions. This chapter is the
experiment that design was for, and it turns the closed-judge case from an
eligibility filter into a measurement.

\section{Why closed judges are the sharpest case in this document}

A model-graded benchmark scores its candidate with another model. When that judge is
a dated proprietary deployment, the substrate is not merely unrecorded --- it is
\emph{gone by construction}, and may no longer be servable at any price. Procedural
reproduction of such a run is not expensive but impossible, which is why this
document treats an unavailable judge as an eligibility gate on \emph{procedural}
reproduction while leaving the run fully in scope for \emph{artifact
reconstruction}: the judge's per-instance verdicts survive in the released record
even though the judge does not. The unobservable cell is itself evidence for the
thesis, and it is a rhetorical asset rather than a limitation --- but only if it is
stated as problem-evidence, never as solution-evidence.

\section{The design and its gate}

Candidate outputs are treated as \emph{immutable, content-addressed response
snapshots}; judgment is a separate fan-out stage over them. That boundary is what
makes substitution measurable: no candidate inference is rerun, so a judge arm
cannot silently change what is being judged. The hash covers only judging-relevant
facts, and each snapshot also carries the verbatim published \code{stats.json} and
\code{per\_instance\_stats.json} so the gate below survives a corpus move.

The load-bearing device is the \textbf{official-annotation identity-replay gate}:
reattach the \emph{original} annotations to a reconstructed scenario state, run the
\emph{real} official metric, and require the result to match the published
statistics to $10^{-12}$. It is a reconstruction test, not an experiment, and it
paid for itself during development by catching two fixture errors before any judge
request was sent --- official \code{safety\_score} aggregation is judge-count
weighted rather than a mean of instance means, and real published statistics carry
derived \code{computed\_on=worst} rows even for unperturbed runs.

\begin{table}[H]
\centering
\small
\begin{tabular}{llr}
\toprule
\textbf{Benchmark} & \textbf{Suite} & \textbf{Max.\ error vs published} \\
\midrule
xstest               & safety/v1.14.0        & $0$ \\
simple\_safety\_tests & safety/v1.14.0        & $0$ \\
harm\_bench          & safety/v1.14.0        & $0$ \\
anthropic\_red\_team & safety/v1.14.0        & $0$ \\
omni\_math           & capabilities          & $0$ \\
wildbench            & capabilities/v1.12.0  & $1.95\times10^{-14}$ \\
\bottomrule
\end{tabular}
\caption{Identity-replay gate on real public \code{gpt-oss-20b} runs. Six of six
benchmarks reproduce their published judge metrics from the reconstructed snapshot,
proving the reconstruction faithful and the metric denominators equal to the
official runs. This is \emph{artifact reconstruction} in the three-target
vocabulary of \S\ref{sec:targets-3}, executed at machine precision.}
\label{tab:replaygate}
\end{table}

One real-data fix landed between the offline fixtures and this table: official
WildBench runs use \code{adapter\_method=chat}, and the audit --- correctly and
loudly --- skipped them because it whitelisted generation only. Supported adapter
methods became a per-benchmark profile field with a curated allow-list rather than a
blanket accept, and the gate on the real chat run then \emph{proved} the
reconstruction faithful instead of the author asserting it.

\section{What the substitution measures}

The v1 experiment completed 12 of 12 attempts in $\sim$20.3 hours: full XSTest (450
instances) and WildBench (1000), two judge arms (Qwen3.5-27B and Qwen3.6-35B-A3B),
three replicates, temperature 0. A subsequent size sweep added the Qwen3.5 ladder
(0.8B/2B/4B/9B) and the remaining safety benchmarks, for 63 healthy artifacts with
zero request errors. Three findings, in ascending order of interest:

\begin{enumerate}
  \item \textbf{On label metrics, open judges are a drop-in replacement.} XSTest
        agreement with the official GPT-4o judge is $\kappa=0.936$ (Qwen3.5-27B) and
        $0.928$ (Qwen3.6-35B-A3B), while the two \emph{official} judges agree with
        each other at only $\kappa=0.829$. The open-versus-closed perturbation is
        smaller than the closed ensemble's own internal disagreement. Agreement
        saturates early --- $\sim$98--99\% from 4B upward, with the 27B buying
        nothing over the 9B.
  \item \textbf{On a graded rubric, the aggregate reproduces and the instances do
        not.} Open judges run $\sim$0.8 lower than officials on WildBench, but
        open-versus-open agreement ($\approx$0.66 absolute, Pearson 0.935) is
        indistinguishable from official-versus-official ($\approx$0.63, 0.936).
        The run-level mean is stable across replicates ($\sigma\approx0.22$ on
        $\approx6.7$, $\approx3\%$) while \textbf{43--46\% of individual instances
        change score}, with a maximum range of 5--6 points on a 1--10 scale.
  \item \textbf{The mechanism, and the generalization worth publishing.} Judge
        \emph{text} is non-deterministic at temperature 0 on both benchmarks:
        87--96\% of instances generate different raw output across replicates.
        The cause is server-side batching --- continuous batching makes the
        floating-point reduction order depend on batch composition, which flips the
        argmax at near-ties and cascades through a chain of thought. Whether that
        reaches the metric depends on \emph{metric granularity}: XSTest's near-binary
        label absorbs it (0.2--0.7\% score change), WildBench's 1--10 rubric has no
        attractor (43--46\%). So \textbf{judge non-determinism is universal; metric
        fragility is a property of the metric, not of the judge} --- which lands
        directly in this project's per-metric-fragility framing.
\end{enumerate}

\finding{Parse rate and calibration are independent axes, and a format-based health
check cannot separate them. WildBench parse rate climbs with judge size
(6.8\%~$\rightarrow$~51.6\%~$\rightarrow$~61.2\%~$\rightarrow$~84.6\%~$\rightarrow$~90.7\%)
while XSTest is $\sim$100\% throughout; and Qwen3.5-2B posts 99.9\% parse with 25.7\%
agreement on \code{anthropic\_red\_team}, flagging 740/1000 responses unsafe where
the official ensemble says 989/1000 are safe. A ``does the judge work?'' smoke test
based on output format would pass that arm. Report the two axes separately, always.}

\section{The check that makes finding 3 trustworthy}

Replicates are distinguished by varying \code{Request.random}, which would have been
a self-inflicted explanation for the divergence. It was verified to be cache-key-only
for these clients: the judges derive from \textsc{helm}'s \code{OpenAIClient}, which
has no seed handling, and the cache-key builder merely appends the field. The same
reasoning would be \emph{wrong} for Bedrock, Cohere, Mistral, Reka, or Google GenAI
clients, which do map \code{random} to a seed. The general rule --- before
attributing variance to the environment, prove your own experimental knobs are not
the cause --- is one this document has needed more than once.

\section{Honest caveats}

\begin{itemize}[nosep]
  \item \textbf{Truncation bias.} The WildBench + Qwen3.5-27B arm still truncates
        14.2\% of judgments at its token ceiling, so its replicate statistics rest on
        the 594/1000 instances parsed in all three replicates --- plausibly a biased
        subset, since long or hard instances drop out. The 35B-A3B arm (2\% parse
        failure) shows the same effect, which corroborates but does not substitute.
        Fix truncation before publishing those numbers.
  \item \textbf{One-class benchmarks.} The safety sets are $\sim$99\% one class, so
        ``agreement'' there is essentially a false-positive rate and must be stated
        as such.
  \item \textbf{A single candidate.} Every one of the 63 artifacts scores the same
        candidate model (\code{gpt-oss-20b}), while every interesting endpoint ---
        conclusion survival, ranking preservation --- is defined over a \emph{set} of
        candidates. This is the structural limitation of the experiment as it stands.
        It is unusually cheap to fix in this design, because official responses
        \emph{and} official judgments for every leaderboard model already sit in the
        public corpus, so candidate expansion costs zero candidate inference.
  \item \textbf{Written but not executed.} The kwdagger fan-out that would schedule
        the matrix across all four GPUs, and the Omni-MATH annotator's live path,
        had not run at this revision. The identity-replay gate does \emph{not}
        exercise the latter: replay reattaches original annotations and runs the
        official metric, so the new prompt-construction and parsing code for
        Omni-MATH has never faced a real judge.
\end{itemize}

\section{The contamination caveat}
\label{sec:contamination}

The Qwen3.5 family launched 2026-02-16 (the small sizes $\sim$2026-03-02). Every
benchmark dataset predates that by two to four years, and the candidate model is from
August 2025. What could \emph{not} be established is the publication date of HELM
Safety v1.14.0 and Capabilities v1.12.0 --- the corpus run directories carry no
execution timestamp, and public sources did not yield it. That date is the one that
matters, because \textsc{helm} publishes the official judge scores \emph{and} the
judges' reasoning text; if those were scraped before Qwen3.5's cutoff, ``the open
judge agrees with GPT-4o'' is partly a memorization result.

Our own data argues against \emph{wholesale} memorization --- the 2B arm inverts the
official labels, agreement is size-graded like a capability curve, and WildBench sits
at a systematic offset rather than converging --- but the mechanism worth worrying
about is \textbf{distribution shift, not recall}: a judge trained on this
distribution is plausibly better calibrated on it with no verbatim memory at all.
The honest framing, which every writeup of these numbers must carry, is that they
describe judge/benchmark pairs the judge was likely trained on and are an
\textbf{upper bound} on agreement for a novel benchmark or a private evaluation. The
cheapest decisive test is to swap in a candidate released after the judge's cutoff,
which makes the (prompt, response, judgment) triple unmemorizable even though the
prompts were public; the pipeline is already parameterized by candidate, so it is a
configuration change rather than new code.

A process note belongs with this: the judge release dates had been authored as
guesses with ``confirm against the repository'' comments, and sat unverified for days
until a question depended on them. Flagging uncertainty in a comment is not
resolving it. A date that a \emph{finding} depends on should be verified when the
finding is made.

\section{Recurring failure modes this thread contributed}

\begin{itemize}[nosep]
  \item \textbf{A behavior change invisible to a cache key serves stale results.} A
        parser fix that did not move its \code{parser\_version} returned a cache hit
        and served pre-fix output; the fix never ran. The identity guardrail was
        working as designed --- the honest record of a behavior change is exactly
        what must invalidate the cache. \emph{Rule: any change to parser or prompt
        behavior must move its version field}, and a spec \emph{hash} rather than a
        spec \emph{path} must ride in job identity, because a path does not change
        when the file is edited. This is G14's genus in a different subsystem.
  \item \textbf{Thinking judges draft the official answer tags inside their
        reasoning.} Hit three times, on three formats (safety
        \code{<reasoning>/<score>}, WildBench JSON, Omni-MATH headings). The standing
        fix is to split on the last reasoning terminator and parse only the suffix,
        which is a strict no-op on non-thinking official responses and so preserves
        parse parity with the official annotators.
  \item \textbf{Exit code 0 is not success.} A rejudge whose every request failed
        still exits 0. Attempts are now health-checked by parse status, and
        \emph{dead} (request error, an infrastructure fact) is distinguished from
        \emph{degraded} (malformed output --- the judge genuinely cannot produce the
        format, which is a finding to keep).
  \item \textbf{Shared mutable state defeats hand-partitioned concurrency.} A
        last-writer-wins sidecar directory clobbered a judge's deployment
        registration, and two workers on one cell deadlocked on the shared request
        cache. Content addressing protected the \emph{outputs} and that protection
        was wrongly generalized to safety. Prefer a scheduler to hand-partitioned
        terminal panes --- a lesson the substrate thread then had to learn again
        (\S\ref{sec:contention}).
\end{itemize}

%==============================================================================
\chapter{Closing the Substrate Question}
\label{chap:substrate-closed}
%==============================================================================

\emph{2026-07-21 through 2026-07-23. Led by Jon Crall, against Edward Wang's probe
results and runbooks. Pre-registered protocol at
\file{reproduce/olmo\_models\_combined/deployment\_match/overnight\_confirm\_plan.md};
scientific framing consolidated in
\file{docs/helm-reproduction-research-journal.md} under ``The fp32/agp0
substrate-recovery experiment.''}

The 2026-07-15 consensus made one demand above all others: run the confirm step.
This chapter is that experiment. Its result is the strongest single empirical claim
in this document, and it arrived by refuting the prediction that was registered
before it started.

\section{Setting up an experiment that could fail}

Three propositions were kept deliberately distinct rather than bundled into ``fp32
fixes it'':

\begin{enumerate}[nosep]
  \item[\textbf{A}] Unpinned dtype $\Rightarrow$ the official ran float32.
        \emph{Mechanism-verified from loader source, one family.}
  \item[\textbf{B}] Float32 recovers the official completions. \emph{Probe-only at
        $n=12$.}
  \item[\textbf{C}] The recovered configuration changes aggregates and conclusions.
        \emph{Untested.}
\end{enumerate}

Building the evidence ledger from primary artifacts rather than prose immediately
corrected the record: the phrase ``four OLMo models exactly recovered'' overstated
what existed. \emph{Exact} was literal only for OLMoE (HF float32, 12/12); the three
dense OLMo-2 models were vLLM float32 matches at 10/12, 10/12, and 11/12 --- all
$n=12$ IFEval probes, all probe-only, with the confirm step never executed for any
model. End-to-end local runs at the bf16 default put IFEval strict accuracy
\emph{above} official by $+0.098$ to $+0.126$ on all four instruct models.
Predictions for the night were then frozen in the plan file, with outcome classes
defined in advance (\textsc{exact} $\leq0.01$; \textsc{materially-improved}
$\leq0.5\times$ the bf16 gap; \textsc{unchanged} $\geq0.8\times$; otherwise
\textsc{anomalous}).

\finding{A zero-compute result from existing data, obtained while building the
ledger: over the aggregate score-drift headline, 4 of 25 OLMo model pairs flip their
official-versus-local ordering (\code{gpqa} 3/6, including 13B$\leftrightarrow$32B;
\code{bbq} 1/6), while \code{ifeval} --- which carries the \emph{largest} drift ---
flips none, because its $+0.10$ is close to uniform across models. The Qwen
experiment flips 1 of 201. Two things follow: procedural drift already damages
conclusions in our own data, and \textbf{drift magnitude does not predict conclusion
damage}. The \code{gpqa} official gaps are $\approx1.3\sigma$, so a paired bootstrap
must decide which flips are statistically real before this is published.}

\section{Three nights of infrastructure failure, and what they taught}
\label{sec:contention}

The experiment took three attempts. The failures are recorded because each is a
distinct, transferable class of error, and because two of them are the same error
this document has already catalogued elsewhere.

\begin{enumerate}
  \item \textbf{A template patcher matching a serialized form.} The script that
        neutralizes \code{add\_generation\_prompt} matched the exact string
        \code{\{\% if add\_generation\_prompt \%\}}, but OLMo-2 guards the assistant
        tag with a compound conditional (\code{\{\% if loop.last and
        add\_generation\_prompt \%\}}), so it matched nothing and both launches
        failed loudly at step one. Fixed by neutralizing the bare
        \emph{identifier} and printing each original context for the operator.
        \emph{Never match an exact serialized form of something a third party
        generates; match the invariant and show your work.} The fail-loudly step
        design did its job --- nothing half-launched.
  \item \textbf{A path that is a claim about the container's filesystem.} The
        rewritten templates were written to an audit-store path and passed verbatim
        into the vLLM container, which mounts only its cache directories. The
        engine crashed, the acquire timed out, and no run started. Fixed by writing
        to the host side of a mounted cache directory, referencing it by
        \emph{container} path, and setting the native template runtime key --- which
        is a structural key, so a changed template now makes a distinct deployment
        instead of coalescing onto the stale one. \emph{State you can see is not
        state the executor can see} --- the same genus as the earlier
        fresh-login-shell lease bug.
  \item \textbf{A held lease that was not eviction-proof.} The 13B run served
        successfully for 54 minutes and then died when a concurrent
        deployment-match sweep's per-cell converge tore down its leased endpoint;
        three workloads shared one serving controller. The recipe was never at
        fault. This is the third appearance of \emph{shared mutable state plus
        hand-partitioned concurrency} in this record, and it argues an
        infrastructure gap worth fixing upstream: a demand-refcount or placement
        re-solve should never stop a \emph{leased} container. The immediate remedy
        was to serialize; the run resumed from its request cache rather than
        restarting.
\end{enumerate}

A fourth failure was scientific rather than operational and is the more instructive
one. A cross-family sweep on a \code{marin-8b-instruct} official completed and
appeared to refute its registered prediction --- until its resolution record showed
\code{protocol\_resolved: false}. The resolver had failed to detect chat markers in
that model's template and \emph{silently defaulted} to raw completions, noting it
only in a free-text field, so all 32 cells probed wrong-shaped prompts. The
registered prediction was \textbf{untested, not falsified}, and the plan table was
corrected to say exactly that. \emph{A resolver that silently defaults on failure
converts ``I don't know'' into a wrong experiment}; for instruct models the sweep
should hard-fail on an unresolved protocol.

\section{Result 1: float32 is necessary and not sufficient}

Both end-to-end runs completed on the ordinary \textsc{helm} from-spec path, with
only dtype and template rendering moved relative to the bf16 baseline. Denominators
were verified equal --- 1082 per-instance rows on each side for each model --- and
IFEval has no judge or parse step, so there is no missing-data trap.

\begin{table}[H]
\centering
\small
\begin{tabular}{lccccc}
\toprule
\textbf{Model} & \textbf{Official} & \textbf{Local fp32+agp0 (vLLM)} &
$\mathbf{D_{\text{fp32}}}$ & \textbf{Prior bf16 gap} & \textbf{Closed} \\
\midrule
OLMo-2-7B-Instruct  & 0.6929 & 0.7597 & $+0.067$ & $+0.098$ & 32\% \\
OLMo-2-13B-Instruct & 0.7298 & 0.8121 & $+0.082$ & $+0.126$ & 35\% \\
\bottomrule
\end{tabular}
\caption{End-to-end \code{ifeval\_strict\_accuracy}, official versus a local float32
replay with old-template rendering, full $n$. The registered prediction
($D_{\text{fp32}}\rightarrow0$) is \textbf{refuted}: both land in the
pre-defined \textsc{anomalous}/partial class. The directional prediction (float32
helps) is confirmed; the magnitude prediction is not.}
\label{tab:fp32e2e}
\end{table}

Two tells said the residual was systematic rather than noise: the closure fraction is
$\approx1/3$ on \emph{both} sizes, and the residual is same-signed on both. The
leading hypothesis was the \emph{engine} --- the officials ran through \textsc{helm}'s
in-process \code{HuggingFaceClient}, and this replay ran through vLLM at the same
precision.

This was, at the time, a better paper result than ``float32 fixes it,'' and it is
worth recording that judgment because it was made before the next experiment was
run: precision and template rendering are \emph{necessary} (bf16 was worse; modern
template rendering was far worse) but not \emph{sufficient}, so substrate recovery
is layered and a single fix is not enough.

\section{Result 2: the official reproduces byte-exactly on its own engine}

The decisive test replaced vLLM with the engine that produced the officials: pure
\code{transformers.generate()} at float32, on one GPU, with no serving stack, no
lease, and no template file --- so the three failure modes of \S\ref{sec:contention}
could not recur.

It nearly failed for a fourth reason. The first wrapper pinned only dtype and decode
and left the probe's forward-pass axes at their sweep defaults, producing a 16-cell
$\times$ 541-instance job (a day or more, not a night) that additionally started on
the template rendering already known \emph{not} to match. \emph{When wrapping a tool
that sweeps by default, the wrapper must pin every axis it does not intend to
sweep}; an unpinned axis silently multiplies the job and silently chooses the
experiment. Corrected to a small-$n$ configuration search over the two axes that can
move greedy float32 logits, it finished in under an hour.

\begin{table}[H]
\centering
\small
\begin{tabular}{llccc}
\toprule
\textbf{Model} & \textbf{Cell (fp32, agp0, decode=helm)} & \textbf{exact} &
\textbf{quasi} & \textbf{first-40} \\
\midrule
OLMo-2-7B  & \code{device\_map=single} & 1.00 & 1.00 & 1.00 \\
OLMo-2-7B  & \code{device\_map=auto}   & 1.00 & 1.00 & 1.00 \\
OLMo-2-13B & \code{device\_map=single} & 1.00 & 1.00 & 1.00 \\
OLMo-2-13B & \code{device\_map=auto}   & 1.00 & 1.00 & 1.00 \\
\bottomrule
\end{tabular}
\caption{Hugging~Face float32 with old-template rendering reproduces the official
OLMo-2 instruct IFEval completions \textbf{byte-identically}, $32/32$ instances on
both models. All four forward-pass cells match at composite $1.000$, so attention
implementation (eager versus \textsc{sdpa}) and device placement are non-factors for
these models on 96~GB cards, where nothing shards. This also resolves
\S\ref{sec:residual-resolved}.}
\label{tab:hffp32}
\end{table}

An attribution probe then split the remaining ambiguity. Our vLLM replay had used the
run spec's \code{temperature=0.0}, which vLLM executes as true greedy, whereas
\textsc{helm}'s client maps \code{temperature==0} to
\code{do\_sample=True, temperature=1e-7} --- so the \emph{same} run spec produces
different decode behavior on different engines. Running the Hugging~Face probe with
true greedy decoding \emph{also} reproduced the official exactly, and the two decode
modes produced byte-identical completions on all 32 instances. Near-zero temperature
does not flip the argmax for these instances, so decode semantics are not the driver.

\finding{The residual is \textbf{pure engine numerics}. Same weights, same float32,
same prompt, same greedy decode --- only vLLM versus Hugging~Face differs --- and the
published metric moves by $+0.067$ and $+0.082$. To faithfully reproduce a
\code{HuggingFaceClient} official you must run Hugging~Face, or a numerically
matching engine. The inference engine is part of the experiment.}

\section{The decomposition}

\begin{table}[H]
\centering
\small
\begin{tabular}{llc}
\toprule
\textbf{Recipe} & \textbf{Substrate variables recovered} &
\textbf{IFEval drift vs official} \\
\midrule
bf16 + modern chat template & none & $+0.10$ \\
vLLM fp32 + old-template rendering & precision, template & $+0.07$ \\
HF fp32 + old-template rendering (either decode) & precision, template, engine &
\textbf{exact} \\
\bottomrule
\end{tabular}
\caption{OLMo-2 instruct IFEval, both 7B and 13B. Each unrecorded substrate layer
closes a defined slice of the drift; recover all three and the published number
reproduces byte-exactly. This is the cleanest available evidence for the
reproduction-first thesis: the number \emph{is} exactly recoverable, but only with
the full unrecorded execution substrate.}
\label{tab:decomposition}
\end{table}

\section{What this establishes, and what it does not}

\paragraph{Established.} For OLMo-2 7B and 13B instruct on IFEval: proposition~A is
confirmed observationally as well as deductively (the official \emph{is} float32,
to byte-exactness, not by inference from a loader default); proposition~B holds only
with the engine included; and the drift decomposes into three named, individually
measured layers.

\paragraph{Not established, and easy to overclaim.}
\begin{itemize}[nosep]
  \item \textbf{Scope.} One benchmark, one family, two sizes, instruct only. The
        cross-family and cross-benchmark generalizations are designed and queued
        (below) but not run at this revision.
  \item \textbf{Level of evidence.} The byte-exact result is a \emph{completion}-level
        match on 32 instances per model, not a \textsc{helm} metric run on the
        in-process client. Identical completions imply an identical IFEval score, so
        the metric claim is a sound inference --- but it is an inference, and a
        full-$n$ confirm or the unwired in-process routing switch
        (\S\ref{sec:substrate}) would make it a measurement.
  \item \textbf{Direction of the residual.} The local scores \emph{above} official in
        every unrecovered configuration. Nothing here explains why the engine gap is
        signed the way it is; it is characterized, not explained.
  \item \textbf{Preservation.} These artifacts are new, and none of them has been
        copied, hashed, or bundled (\S\ref{sec:consensus-closed}).
\end{itemize}

\section{From one cell to a result}

The method is now proven end-to-end on a single cell; the open work is making it
systematic. Three tasks are specified in
\file{docs/planning/edward-task-queue.md}:

\begin{itemize}[nosep]
  \item \textbf{Cross-family generality} --- a fan-out over Pythia-6.9B, Vicuna-7B,
        and Granite-4.0-Micro as \emph{treatment} (unpinned dtype: predict float32
        wins) with a dtype-\emph{pinned} official (Gemma-2-9B-it, bf16) as
        \emph{control} (predict bf16 wins, float32 loses). Predictions registered
        before execution; a bidirectional control is what makes the treatment
        result more than a coincidence.
  \item \textbf{Cross-benchmark generality} --- the same ladder on OLMo-2's
        \code{bbq}, \code{gpqa}, and \code{mmlu\_pro} officials. Multiple-choice
        benchmarks make completion-level exact match easy and therefore less
        diagnostic, so the \emph{metric} delta must be reported alongside it. This
        is also where the effect may or may not move a \emph{conclusion} --- recall
        that \code{gpqa} carried 3 of 6 pairwise flips at bf16.
  \item \textbf{A prospective reproduction census} --- the paper backbone, and the
        one that must be frozen before outcomes are seen: a stratified sample drawn
        by a fixed rule from the eligible corpus, one diagnostic ladder applied in
        the same order to every run (from-spec replay $\rightarrow$ dtype
        $\rightarrow$ tokenizer/template $\rightarrow$ engine $\rightarrow$
        residual), outcome buckets defined in advance (\textsc{exact} /
        \textsc{recovered-with-cause} / \textsc{unresolved} / \textsc{blocked}), and
        a per-run stopping rule. \emph{Blocked $\neq$ irreproducible.}
\end{itemize}

Only the third converts this chapter from a deep case study into a population claim,
and it is the one the eligibility-gate and denominator work of
Chapter~\ref{chap:cartography} was always for.

%==============================================================================
\chapter{Paper Direction and the Manuscript}
\label{chap:paperdirection}
%==============================================================================

\emph{2026-07-21 through 2026-07-27. Thesis assessment led by Jon Crall
(\file{docs/planning/tmlr-paper-thesis.md}); the manuscript itself
(\file{docs/papers/tmlr-2026/}) is Edward Wang's.}

\section{Three adversarial rounds on the thesis}

With two threads running --- substrate recovery and judge substitution --- the
open question became which is the paper. It was worked as three adversarial rounds
rather than a drafting exercise, and the record is worth keeping for two reasons: it
is where the single-candidate limitation of Chapter~\ref{chap:openjudge} was
diagnosed, and it is a case study in verifying claims across model-authored reviews.

\begin{itemize}
  \item \textbf{Round 1} concluded that the infrastructure was ahead of the science:
        every proposed headline framing for the judge work was unsupported for one
        structural reason --- all artifacts score a single candidate. The resulting
        priority ordering (candidates $>$ judge families $>$ benchmarks $>$ judge
        sizes) inverted the intuition it was tested against.
  \item \textbf{Round 2} adjudicated a rebuttal. Its load-bearing step was
        \emph{literature verification}: two papers were cited against our novelty
        claim, one of them postdating the reviewing model's knowledge cutoff with a
        chat-sourced URL. Both checked out on search --- JuStRank
        (arXiv:2412.09569) benchmarks judges by induced system rankings, and
        SLMJury (arXiv:2606.07810) sweeps small judges across ten benchmarks --- so
        the ``consumer-sized judge sweep'' framing is occupied and our size ladder
        can never be the headline. Neither touches substitution \emph{inside a
        published leaderboard's official pipeline under exact replay}, so the wedge
        survives, narrower. \emph{When two models argue about novelty, the citations
        are the part to verify first}: a hallucinated occupier would have wrongly
        shrunk the claim, and a real one ignored would have sunk the introduction.
  \item \textbf{Round 3} reopened the question against Edward's draft and a program
        brief, and landed --- weakly held --- on the position that the strongest
        object of reproduction is the \emph{experiment}, not the finding, with the
        judge work as the modern-failure-mode section rather than the spine. The
        hostile reviews of the two candidate directions each patch the other, which
        is the best argument that they are one paper. The decision remains gated on
        facts outside the repository: the program's historical scope, the boundary
        between what the inherited EEE paper already claims and what is new here,
        and staffing timelines.
\end{itemize}

\noindent A reversal was recorded honestly in the same window: a full-leaderboard
judge sweep recommended in the morning was paused by that evening on new
information. \emph{A decision optimized under an assumption must be re-costed when
the assumption moves, not carried as settled} --- the same lesson the deferred
scheduler in Chapter~\ref{chap:openjudge} taught, and holding a sweep is cheaper
than unwinding one.

\section{The manuscript}
\label{sec:paper-drafting}

The TMLR draft under \file{docs/papers/tmlr-2026/} was scaffolded from this master
reference on 2026-07-15 and substantially rewritten on 2026-07-27 (14 commits over
\file{introduction.tex}, \file{related\_work.tex}, \file{framework.tex}, and
\file{references.bib}). Four changes matter to \emph{this} document because they
sharpen framings it originated:

\begin{itemize}[nosep]
  \item \textbf{The trichotomy is credited, not replaced.} Related work now states
        the standard methods/results/inferential trichotomy as Goodman's, positions
        our claim-level robustness as inferential reproducibility and our procedural
        reproduction as methods reproducibility \emph{with its central premise
        removed} --- ``the same data and tools,'' where the tools are precisely what
        did not survive --- and narrows the taxonomy claim accordingly. Adjacent
        vocabularies (artifact badging, documentation-completeness grading) are
        described as cutting the space differently rather than restating it.
  \item \textbf{Artifact reconstruction is argued for, not merely defined.} The
        justification is the closed-judge case of Chapter~\ref{chap:openjudge}: when
        the judge is no longer served, re-execution is impossible and artifact
        reconstruction is the \emph{only} available target, yet the per-instance
        verdicts survive, so a scorer or migration defect remains detectable with no
        judge in the loop. A trichotomy in which every category re-runs something
        has nowhere to put that case.
  \item \textbf{The eligibility gate is target-relative.} An unavailable judge
        filters a run out of \emph{procedural} reproduction while leaving it in scope
        for \emph{artifact reconstruction}. This is the correct statement of a
        distinction this document has made since Chapter~\ref{chap:framing}, and
        it needed the judge experiment to become concrete.
  \item \textbf{The bibliography was verified against sources}, and four records
        were wrong and fixed. Given that the same window's thesis work turned on
        verifying two citations, this is not incidental hygiene.
\end{itemize}

\noindent The manuscript is a draft: introduction, related work, framework,
provenance, and system sections exist in revised form; the empirical sections await
the census of Chapter~\ref{chap:substrate-closed} and the candidate expansion of
Chapter~\ref{chap:openjudge}.

%==============================================================================
\chapter{Synthesis: What Is Reproducible, and Why}
\label{chap:synthesis}
%==============================================================================

\section{The headline scientific claims}

\emph{This section was revised on 2026-07-15 (see Appendix~\ref{app:revision}) to
scope its claims to the evidence actually in hand. The original draft asserted a
positive result --- ``a runnable subset of \textsc{helm} \emph{is} reproducible,
with residual drift that is small, concentrated, and attributable.'' That is a
plausible and motivating \emph{hypothesis}, but the experiments that would
establish it (a frozen population, freshly regenerated stores, end-to-end OLMo
confirmation on held-out instances, completed Qwen/classic runs, cross-machine
measurements, and a controlled ablation) are not all done. The claims below are the
defensible ones.}

\medskip
The more defensible thesis --- and the more scientifically interesting one --- is
that \emph{reproducing an \textsc{llm} benchmark is a layered system-identification
problem.} A public recipe is not a complete experimental record; the producing
measurement instrument spans the recipe, the model deployment, the software stack,
hardware-sensitive numerical behavior, and the artifact-processing history.
Controlled reconstruction can \emph{attribute} and sometimes \emph{close} apparent
reproducibility gaps; but when provenance has been lost or transformed, the original
experiment can be \emph{non-identifiable from the surviving evidence examined}. What the internship established, at the
evidence level actually reached:

\begin{enumerate}
  \item \textbf{The unrecorded execution substrate is real and consequential
        (established).} \textsc{helm} records the recipe and the model name but not
        the dtype, revision, tokenizer/template version, attention kernel, or
        software-stack versions (of 148 HF-client deployments, 19 pin dtype, 7 pin
        revision; Appendix~\ref{app:unrecorded}). Individual substrate parameters are
        \emph{identifiable to varying degrees}: the loader's dtype default and the
        deployment client class are readable from source/config; the tokenizer's
        special-token append is readable from the repo; the exact
        \code{transformers}/\code{torch}/CUDA versions and hardware are not.
  \item \textbf{Attribution on OLMo: established for one cell, open as a population
        claim} (revised 2026-07-27). The confirm step this entry called the top open
        experiment was run (Chapter~\ref{chap:substrate-closed}). For OLMo-2 7B and
        13B instruct on \code{ifeval}, the drift decomposes into three named,
        individually measured layers --- precision, chat-template rendering, and the
        inference engine --- and recovering all three reproduces the official
        completions \emph{byte-exactly}, while recovering only the first two leaves a
        $+0.067$/$+0.082$ residual attributable to engine numerics alone. The
        deductive dtype-default argument (\S\ref{sec:fp32}) is thereby corroborated
        observationally; the OLMo-7B prompt-independent gibberish remains attributable
        to a tokenizer EOS-append and OLMoE-instruct disagreement to
        \code{add\_generation\_prompt} drift. \textbf{What is still not established:}
        this is one benchmark, one family, two sizes, instruct-only, and the
        byte-exact evidence is at the completion level on 32 instances per model
        rather than a metric-level run. ``Precision is \emph{the} dominant variable''
        is now known to be \emph{false as stated} --- precision alone closes only
        about a third of the gap for this cell.
  \item \textbf{The inference engine is a first-class hidden variable
        (established).} Holding weights, precision, prompt, and decode fixed and
        changing only vLLM $\leftrightarrow$ Hugging~Face moves a published benchmark
        metric by $+0.07$--$0.08$. This was previously argued from kernel behavior
        and per-instance flip rates (G7); it is now measured at the level of the
        published number, with decode semantics explicitly excluded as the cause. A
        corollary worth stating separately: the \emph{same} run spec
        (\code{temperature=0.0}) means true greedy on one engine and
        near-zero-temperature sampling on another, so a nominally identical recipe
        does not fix the decode procedure.
  \item \textbf{Judge non-determinism is universal; metric fragility is not
        (established, single candidate).} Substituting open-weight judges into six
        model-graded benchmarks behind a machine-precision identity-replay gate
        (Chapter~\ref{chap:openjudge}) shows 87--96\% of judgments producing
        different text across replicates at temperature~0, from server-side batching.
        Whether that reaches the metric depends on metric granularity: 0.2--0.7\% on
        a near-binary safety label, 43--46\% on a 1--10 rubric, with the run-level
        mean stable in both cases. Scoped by two limits: every artifact scores a
        single candidate model, and the agreement figures are an upper bound under
        the contamination caveat of \S\ref{sec:contamination}.
  \item \textbf{A recoverable-vs-non-identifiable contrast (the strongest conceptual
        result).} OLMo is the \emph{partially recoverable} case: the missing provenance (dtype,
        tokenizer, template) can be reconstructed and its effect measured. The
        original-paper classic models (GPT-J, GPT-NeoX, OPT) are the
        \emph{non-identifiable} case: their run-specs name a class path that resolves
        in no released \textsc{helm} and triangulates to unreleased pre-v0.1.0 code
        (G13, \S\ref{sec:g13}), so the producing instrument cannot be uniquely
        identified or rebuilt. This contrast --- some lost provenance is
        experimentally recoverable, some is fundamentally underdetermined --- is a
        sharper contribution than a uniformly positive ``we closed the gaps'' story.
  \item \textbf{Non-runnable cases are a separate gate, not a reproducibility
        verdict.} Roughly half of the original broad grid could not be \emph{run at
        all} (gated datasets, revoked download access, missing credentials, packaging,
        resource contention). These are \emph{filtering/eligibility} reasons upstream
        of any disagreement measurement --- they answer ``can we run this?'', not
        ``did it reproduce?'' --- and must not be folded into a negative result.
\end{enumerate}

\section{The reusable methodology}

Independently of the specific numbers, the internship produced a reusable
methodology for reproducibility auditing that the paper can present as a
contribution:

\begin{itemize}[nosep]
  \item \textbf{Replay the resolved recipe, never reconstruct it}
        (Chapter~\ref{chap:fromspec}): from-spec replay of \file{run\_spec.json}
        eliminates a whole class of silent reconstruction drift.
  \item \textbf{Freeze the measurement instrument by era}
        (Chapter~\ref{chap:eras}): version-pinned harness containers make the
        scoring code a controlled variable, separable from the model.
  \item \textbf{Sweep, don't guess} (Chapter~\ref{chap:deploymatch}): for unrecorded
        execution parameters, an empirical sweep with feasibility-pruning (but not
        relevance-pruning) is the honest analogue of matching a recipe.
  \item \textbf{Classify every disagreement} (\S\ref{sec:distinction}): the
        recipe/environment-vs-reproducibility distinction, enforced through typed
        exclusion reasons, honest failure-taxonomy colors, and
        \code{comparability\_unknown:*} signals rather than silent fallbacks.
  \item \textbf{Report a distribution, not a point}: same-machine, cross-machine,
        and official-vs-local agreement, with per-metric tolerance sweeps, so that
        ``reproducible'' is quantified rather than asserted.
  \item \textbf{Register the prediction before the run} (added 2026-07-27,
        Chapter~\ref{chap:substrate-closed}): freeze outcome classes and a
        per-configuration prediction in a plan file before execution. The
        substrate experiment's headline value came from a \emph{refuted} prediction,
        which is only interpretable because the prediction and its classification
        thresholds were written down first. Pair treatments with a bidirectional
        control (a dtype-\emph{pinned} official predicted to reproduce at its pinned
        dtype and to fail at float32), so a confirmation cannot be a coincidence.
  \item \textbf{Gate a reconstruction before you experiment on it} (added
        2026-07-27, Chapter~\ref{chap:openjudge}): reattach the original artifacts,
        run the original metric, and require the published statistics back to machine
        precision. It validates the experiment, and during development it debugs the
        reconstruction --- it caught two wrong fixtures before any judge request was
        sent.
  \item \textbf{Pin every axis you do not intend to sweep}: a sweep tool defaults
        the rest, and a defaulted axis is a silent experimental choice. Two of this
        record's wrong conclusions (\S\ref{sec:residual-resolved}) and one 16$\times$
        over-scoped overnight run trace to unpinned axes --- the unrecorded-parameter
        problem committed by our own instrument.
\end{itemize}

\section{The provenance recommendation for the field}

The clearest actionable output for the broader community is that \textsc{helm} (and
any benchmark harness) should record the \emph{execution substrate}, not just the
recipe and the model name. Three fields --- \code{model\_revision} (commit SHA),
\code{torch\_dtype}, and \code{transformers\_version} --- transitively close the
majority of the reproducibility gap, because the revision fixes the
tokenizer/template/generation-config state and the transformers version fixes the
precision default and the chat-template behavior. These belong in a machine-readable
provenance block (the existing \code{recipe\_facts} slot suffices with no schema
change), kept \emph{out} of the run-spec identity string so they do not break
official$\leftrightarrow$local pairing (Appendix~\ref{app:unrecorded}).

\emph{Revised 2026-07-27:} make it four fields. The
2026-07-23 measurement showed the first three to be necessary and
\emph{insufficient} --- they closed about a third of the drift on the one cell taken
end to end, and the serving \textbf{engine and its version} closed the remainder
(Chapter~\ref{chap:substrate-closed}). For model-graded runs, add the judge model
and ensemble composition, which today live only inside the harness's annotator
source and change between harness versions (\S\ref{sec:judge}). And state the limit
honestly: server-side batch composition on a hosted endpoint is not recordable after
the fact by any of these means, which is precisely why runs served by external
providers belong to the irrecoverable-substrate category rather than to a list of
fields someone forgot to write down.


\section{Per-parameter identifiability map}

The evidence suggests that identifiability is not a binary property of an entire
run. It can be assessed separately for each latent execution parameter.

\begin{longtable}{@{}p{0.25\textwidth}p{0.20\textwidth}p{0.47\textwidth}@{}}
\toprule
\textbf{Parameter} & \textbf{Status in studied cases} & \textbf{Basis and limitation} \\
\midrule
\endhead
Benchmark recipe / run spec & Often identifiable & Frozen public
\file{run\_spec.json}; artifact migrations can still alter interpretation. \\
Client / deployment class & Often identifiable & Recorded deployment name and
HELM configuration/source; external hosted internals remain opaque. \\
Load dtype & Partly identifiable; \emph{recovered} in one case & Explicit when
pinned; otherwise inferable only conditional on model-loading code and
package-version bounds. For OLMo-2 instruct the conditional inference was
subsequently confirmed by byte-exact reproduction at float32
(Chapter~\ref{chap:substrate-closed}), which is the strongest form of
identification available: the parameter is recovered, not merely bounded. \\
Tokenizer special-token behavior & Often recoverable & Repository/configuration
inspection and token-level prompt comparison. Exact historical tokenizer revision
may remain absent. \\
Chat template / generation prompt & Often recoverable & Revisioned tokenizer files
or controlled rendering sweeps; effective behavior depends on package version. \\
Model/tokenizer revision & Frequently unidentified & Rarely pinned in HELM
metadata; aliases and mutable repositories prevent unique historical recovery. \\
Inference engine and version & Partly identifiable, and \emph{load-bearing} & Client
family may be known; exact serving version, kernels, scheduler, and generation
defaults often are not. Elevated by the 2026-07-23 result: the engine is not a
second-order detail but a variable that moves a published metric by $+0.07$ at fixed
weights, precision, prompt, and decode. Where the client family \emph{is} recorded
(as \code{HuggingFaceClient} was here), matching it is both necessary and
sufficient in the studied cell; where officials were served by an external provider,
this row is unidentifiable and the run belongs to the irrecoverable-substrate
category by design. \\
Judge model and its deployment (model-graded runs) & Registry-recoverable in
identity, unrecoverable in substrate & The judge model is not in
\file{run\_spec.json}; it is hard-coded per annotator class and \textsc{helm}
version, so a curated registry recovers \emph{which} judge
(\S\ref{sec:judge}) --- and the official ensemble composition is itself
version-dependent. The judge's own serving substrate is a dated proprietary
deployment that may no longer exist, so procedural reproduction is barred while
artifact reconstruction remains fully available (Chapter~\ref{chap:openjudge}). \\
Batching, cache, topology, numerical flags & Usually unidentified & Not represented
in public artifacts and often provider-specific. \\
Hardware / CUDA / driver & Usually unidentified historically & Can be recorded and
robustness-tested prospectively, but exact historical replay may be impossible. \\
Harness era & Recoverable for tagged runs; proxy-only for some classics & Tagged
release commits support era containers. Migrated paths can fingerprint an unreleased
origin without yielding a buildable exact instrument. \\
\bottomrule
\end{longtable}

\section{Open threads at the close of this chronology}
\label{sec:open-threads}

\emph{Rewritten 2026-07-27 against the state at this revision. Items discharged since
July 14 are marked so, because a reader comparing revisions needs to see what moved.}

\paragraph{Discharged since the previous revision.}
\begin{itemize}[nosep]
  \item \textbf{Ordinary-path OLMo confirmation} --- run, with a refuted prediction
        and a byte-exact resolution (Chapter~\ref{chap:substrate-closed}).
  \item \textbf{The dense OLMo-2 Hugging Face divergence} --- resolved as a probe
        artifact (\S\ref{sec:residual-resolved}).
  \item \textbf{The open-judge extension} --- built and executed
        (Chapter~\ref{chap:openjudge}).
\end{itemize}

\paragraph{Open, ordered by what the paper needs.}
\begin{enumerate}[nosep]
  \item \textbf{The prospective reproduction census} --- the sampling rule, the
        diagnostic ladder, the outcome buckets, and the stopping rule must be frozen
        \emph{before} outcomes are seen, and the first stratum bucketed. This is what
        converts one deep case into a population claim, and it is the single highest
        item on the list.
  \item \textbf{Cross-family and cross-benchmark generality} of the substrate
        decomposition, with registered predictions and a dtype-pinned control
        (Chapter~\ref{chap:substrate-closed}).
  \item \textbf{Candidate expansion for the judge experiment} --- every endpoint of
        interest is defined over a set of candidates, and official responses and
        judgments for the whole leaderboard already sit in the mirrored corpus, so
        this costs judge inference only.
  \item \textbf{The contamination control} --- rejudge a candidate released after the
        judge's training cutoff (\S\ref{sec:contamination}); until then every
        agreement figure is an upper bound. Establishing the HELM Safety v1.14.0 /
        Capabilities v1.12.0 publication dates would sharpen the caveat.
  \item \textbf{Preservation, now the most overdue item.} The 2026-07-15d audit found
        the two flagship modern stores pruned to their derived layers, so
        ``regenerate'' means \emph{re-run} for GPT-OSS and \code{olmo-models-combined};
        \code{era-redpajama} and the deployment-match sweeps are preservable by
        copy-and-hash as they stand; and the July 21--23 substrate artifacts are
        \emph{new} and equally unpreserved. Nothing in this thread is bundled.
  \item \textbf{Regenerate the corpus-wide denominator} from a checked-in manifest;
        the $\approx$59\% classic-track figure remains an internal estimate.
  \item \textbf{The HuggingFace-in-process routing switch} --- still unwired, and now
        needed for a narrower reason: to lift the byte-exact result from the
        completion level to a metric-level \textsc{helm} run (\S\ref{sec:substrate}).
  \item \textbf{Expand-once-then-freeze for compute runs} --- proposed, not
        implemented; until it lands our own new runs store a mutable run-key string
        as their source of truth (\S\ref{sec:freeze-critique}).
  \item \textbf{Pinning model revisions} through \code{eval\_audit}'s bundle boundary
        (a $\sim$3--4 file change) so generated bundles carry the SHA automatically.
  \item \textbf{Two infrastructure robustness gaps} surfaced by the failed nights: a
        held lease should be evict-proof against a concurrent placement converge, and
        a protocol resolver should hard-fail rather than silently default
        (\S\ref{sec:contention}).
  \item \textbf{Unexecuted code} --- the kwdagger rejudge fan-out and the live
        Omni-MATH annotator path; the Qwen3.5 report; the queued judge-benchmark
        cells.
\end{enumerate}

%==============================================================================
\appendix
\chapter{The \textsc{helm} Gotchas Catalogue (G1--G15)}
\label{app:gotchas}
%==============================================================================

This appendix condenses the running ledger of undocumented or hard-to-see
\textsc{helm} behaviors accumulated while building the audit pipeline
(\file{docs/helm-gotchas.md}). Each is a concrete, reviewer-anticipating pitfall
that a reproducibility paper's appendix should address.

\begin{description}[style=nextline,leftmargin=1em]
  \item[\textbf{G1. \file{run\_spec.json} schema evolves silently between releases.}]
    Newer \textsc{helm} populates \code{adapter\_spec} fields older \textsc{helm}
    omitted (\code{chain\_of\_thought\_prefix}, \code{global\_suffix},
    \code{num\_trials}, \code{model\_deployment}, top-level \code{metric\_specs} /
    \code{groups} / \code{annotators}), so byte-for-byte hashing yields $0/N$
    matches across releases even for identical recipes. Workaround: a
    schema-collapsed \emph{canonical recipe hash} (\S\ref{sec:funnel}).
  \item[\textbf{G2. The local \code{suite} field is the experiment name, not a
    public-track version.}] A versioned join on \code{logical\_key + suite\_version}
    yields $0$ matches; the funnel detects a non-version-shaped local suite and
    reports the versioned join as not meaningful (\code{N/A}) rather than $0$.
  \item[\textbf{G3. The \code{huggingface/<model>} deployment alias does not always
    exist.}] Many open-weight runs were executed through the HF API, but
    \textsc{helm}'s registry ships deployment YAMLs for only a subset; a missing
    alias forces a different stack (vLLM/etc.), introducing deployment-substitution
    drift. Workaround: register a custom local deployment.
  \item[\textbf{G4. \code{model\_deployment} semantics differ across versions.}]
    Older \textsc{helm} did not record the field (defaulting to ``the HF API'');
    newer does. A local \code{huggingface/<model>} and an official
    \code{<MISSING>} are semantically identical but hash differently ---
    schema-evolution, dropped from the canonical hash.
  \item[\textbf{G5. \code{adapter\_spec.instructions} differs between releases on
    identical scenarios.}] MMLU official prepends ``Answer with only a single
    letter.''; several LegalBench scenarios prepend a list-the-allowed-labels
    preamble. This is \emph{genuine} recipe drift (a different prompt), surfaced via
    the \code{same\_instructions} fact --- not to be silently normalized. (The BBQ
    \code{output\_format\_instructions=mcqa} drift of \S\ref{sec:canonkey} is the
    same species.)
  \item[\textbf{G6. The \code{classic} corpus moved bucket prefixes.}] Classic runs
    migrated from the bucket root to \code{classic/benchmark\_output/runs/<ver>/},
    mirroring every other suite; a special-case in the downloader was removed.
  \item[\textbf{G7. Same weights $\ne$ same output across serving stacks.}] Even
    greedy, HF \code{generate()} and vLLM differ in kernels, tokenizer batching,
    stop-sequence handling, unspecified sampling defaults, and EOS/pad handling ---
    a $3$--$6\%$ per-instance flip rate on multiple-choice, $\sim$30--40\% on
    long-form. Mitigations: match the original client, pin decoding, or document as
    \code{deployment\_drift}. (Chapter~\ref{chap:deploymatch} is the deep dive.)
  \item[\textbf{G8. \code{metric\_specs} schema evolution.}] Top-level
    \code{metric\_specs} differ on the majority of packets from restructuring, not a
    recipe change; excluded from the canonical hash --- compare the metric
    \emph{output}, not the declaration.
  \item[\textbf{G9. Local index \code{model\_deployment} is audit-specific.}] Custom
    local deployment names (e.g.\ \code{kubeai/vicuna-...-no-chat-template-local})
    do not appear in the public run spec; the logical-key join collapses across
    deployments and \code{same\_deployment} correctly reports \code{no}.
  \item[\textbf{G10. Public-track \code{suite\_version} is not the \textsc{helm}
    release version.}] The same \file{run\_spec.json} can appear identically in
    \code{v0.2.4} and \code{v0.3.0}; use \code{run\_spec\_hash} to detect
    recipe-identical duplicates across suites. (The era-replay containers key the
    \emph{instrument} on \code{(public\_track, suite\_version)}, which is the right
    granularity for the classic track --- see Chapter~\ref{chap:eras}.)
  \item[\textbf{G11. \file{per\_instance\_stats.json} corruption on giant runs.}]
    Certain msmarco runs were truncated mid-write during upload (consistent
    byte-offset \code{JSONDecodeError}); recently re-uploaded versions parse cleanly.
    Redownload on a size mismatch.
  \item[\textbf{G12. \code{run\_spec\_hash} canonicalization differs by release.}]
    \textsc{helm}'s own canonicalization is version-dependent, so semantically
    identical specs hash differently across releases; the coverage-side canonical
    hash applies a stable normalization on top.
  \item[\textbf{G13. Classic-corpus \code{class\_name} paths resolve in \emph{no}
    \textsc{helm} version.}] The stored \code{helm.benchmark.basic\_metrics.BasicMetric}
    (helm-prefix + flat) is a once-migrated hybrid --- a naive
    \code{benchmark.}$\to$\code{helm.benchmark.} rewrite of a still-flat production
    path --- pinning the producing code to unreleased pre-v0.1.0 \textsc{helm}
    (2022-07-31 to 08-26). Class-path \emph{shape} is the origin fingerprint; fixed
    by a self-verifying declared class-path canonicalization in the era shim
    (\S\ref{sec:g13}).
  \item[\textbf{G14. A compiled-artifact cache keyed narrower than the config space
    that shapes the artifacts.}] vLLM's ahead-of-time compile-cache key omits at
    least \code{limit\_mm\_per\_prompt}: the cache hash was identical before and
    after disabling the vision modalities, so the engine reloaded a graph traced
    under different inputs. We got the loud failure; the quiet one --- silently
    reusing compiled artifacts from a different serving config, producing different
    numerics with no error --- is the reproducibility hazard, and the compiled-graph
    provenance is invisible to \file{run\_spec.json} and to vLLM's own key alike.
    Workaround: key the cache mount on a hash of image, command, and generation
    environment. Note the epistemics: this is observable only when a config changes
    \emph{between two runs sharing a cache}, so cross-machine reproduction runs
    double as cache-freshness controls (Chapter~\ref{chap:qwen35}).
  \item[\textbf{G15. A from-spec run entry and a compute run entry are not
    interchangeable.}] In from-spec mode a run entry is a \emph{lookup key} into the
    official corpus and may carry name-metadata (\code{eval\_split=},
    \code{groups=}); \textsc{helm} replays a frozen \code{run\_spec.json} and never
    calls the scenario constructor. In compute mode the same string \emph{is} a
    constructor call, and those kwargs are rejected --- 57 of 86 jobs failed this way
    on the first Qwen3.5 grid. Any port of a replay grid to a compute grid must strip
    run-entry arguments to what the \code{run\_spec\_function} actually accepts
    (\S\ref{sec:freeze-critique}).
\end{description}

%==============================================================================
\chapter{The Unrecorded-Parameter Taxonomy}
\label{app:unrecorded}
%==============================================================================

Condensed from \file{docs/helm-unrecorded-deployment-params.md}. \textsc{helm}
serializes the recipe and the model's name but almost nothing about \emph{how} the
numbers were computed. Empirically, of 148 HuggingFaceClient deployments in
\textsc{helm}'s \code{model\_deployments.yaml}, only \textbf{19 pin
\code{torch\_dtype}} (all \code{bfloat16}) and only \textbf{7 pin a model
\code{revision}}. The parameters, ranked by effect on a greedy (temperature-0)
output:

\begin{longtable}{@{}p{0.02\textwidth}p{0.30\textwidth}p{0.60\textwidth}@{}}
\toprule
\textbf{Tier} & \textbf{Parameter} & \textbf{Why it matters} \\
\midrule
\endhead
1 & Model weight revision (SHA) & Same name $\ne$ same weights; authors re-upload
    checkpoints/configs/tokenizers. The single highest-leverage field --- fixes the
    tokenizer, chat template, and generation config as-of-then. \\
1 & Load precision / dtype & Unpinned $\Rightarrow$ transformers-version-dependent
    default (pre-v5 = float32). \emph{Confirmed 2026-07-23}: at float32 on the
    original engine, HF reproduces the official OLMo-2 instruct completions
    byte-exactly; at bf16 the published IFEval metric moves by $+0.10$
    (Chapter~\ref{chap:substrate-closed}). \\
1 & Inference engine (client family) & \emph{Promoted to Tier~1, 2026-07-27.} At
    fixed weights, precision, prompt, and greedy decode, vLLM versus HF moves
    published IFEval accuracy by $+0.067$/$+0.082$ --- larger than most Tier-2
    effects and comparable to the precision term itself. \\
1 & Quantization & GPTQ/AWQ/fp8/bnb change the weights outright. \\
1 & Software-stack versions & The meta-parameter: transformers sets the fp32-vs-bf16
    default \emph{and} whether the chat template honors \code{add\_generation\_prompt};
    engine/flash-attn versions change kernels. Nowhere in the run dir. \\
\midrule
2 & Attention implementation & eager/sdpa/flash/PagedAttention change reduction order
    $\Rightarrow$ different logits at temp 0. \\
2 & Matmul / TF32 precision & For fp32 especially, true-fp32 vs TF32 matmuls differ. \\
2 & Hardware / GPU model & Determines kernels, TF32 behavior, whether fp32 MoE fits. \\
2 & Device topology / TP & Multi-GPU all-reduce order differs from single-GPU. \\
2 & Batch composition & Kernels are not batch-invariant; same prompt, different
    batch neighbors, different logits. \\
\midrule
3 & Tokenizer + chat-template version & The \code{add\_generation\_prompt} drift
    (\S\ref{sec:addgenprompt}); the templated prompt is never stored. \\
3 & \code{add\_special\_tokens} / BOS & The OLMo-7B EOS-append case
    (\S\ref{sec:eosappend}). \\
3 & \code{generation\_config.json} defaults & Injects unspecified fields
    (repetition penalty, eos/pad ids). \\
3 & Stop-sequence / detokenization & HF stops on token ids, vLLM on string match
    $\Rightarrow$ different truncation and scored text. A stop sequence can also
    interact with a model's answer \emph{style} to zero a metric outright
    (Chapter~\ref{chap:qwen35}). \\
3 & Decode semantics at \code{temperature=0} & The same recipe field means true
    greedy on one engine and \code{do\_sample=True, temperature=1e-7} on another.
    Measured to be a non-factor for OLMo-2 IFEval (byte-identical completions), but
    it is an engine-dependent reading of a nominally identical recipe, so it must be
    recorded rather than assumed. \\
3 & Compiled-graph / kernel-cache provenance & Compiled artifacts can be reused
    across serving configs whose differences the cache key omits (G14). Invisible to
    the run dir, to the harness, and to the engine's own key. \\
\midrule
4 & RNG seed; KV-cache dtype & Mostly for sampling (temperature $>0$) and
    serving-engine reproductions. \\
\midrule
--- & \emph{Model-graded runs:} judge model, ensemble composition, judge deployment &
    Not in \file{run\_spec.json} at all; hard-coded per annotator class and
    \textsc{helm} version, and the official ensemble changes between versions. The
    judge's serving substrate is typically a dated proprietary deployment ---
    unrecoverable rather than merely unrecorded (Chapter~\ref{chap:openjudge}). \\
--- & \emph{Any hosted endpoint:} server-side batch composition & Continuous
    batching makes a hosted judge or candidate non-deterministic at temperature~0;
    87--96\% of judgments differ across replicates. Unrecorded and, for a historical
    run, unrecoverable in principle. \\
\bottomrule
\end{longtable}

\noindent\textbf{The one non-obvious lever.} Several Tier-3 gaps collapse into
Tier-1: the tokenizer, chat template, and generation config all live \emph{inside}
the model repo, so pinning the model revision recovers all of them as-of-then.
\code{model\_revision} + \code{torch\_dtype} + \code{transformers\_version} close
the majority of the gap. Both weight-loading engines already accept \code{revision}
as a data-only edit (infer-stack \code{catalog.yaml}; \textsc{helm}'s
\code{client\_spec.args}); the only code gap is that \code{eval\_audit}'s bundle
boundary drops it (a $\sim$3--4 file fix). These fields must live in the
\code{recipe\_facts}/provenance block, \emph{not} the run-spec identity string, or a
SHA in the pairing key would break official$\leftrightarrow$local pairing.

\noindent\textbf{A fourth field, added 2026-07-27.} The three-field recommendation
was derived before the engine effect had been measured. It should now read
\code{model\_revision} + \code{torch\_dtype} + \code{transformers\_version} +
\textbf{the serving engine and its version}. In the one cell measured end to end,
the first three are \emph{necessary and insufficient} --- they close roughly a third
of the drift --- and the engine closes the rest. \textsc{helm} already records the
client \emph{class}, which is what made the recovery possible at all; what is missing
is the engine \emph{build}, and, for hosted endpoints, an acknowledgement that
server-side batch composition is not recoverable at any level of diligence. The
practical form of the recommendation is unchanged: these are data-only fields in a
provenance block, and three of the four are already accepted by both weight-loading
engines today.

%==============================================================================
\chapter{Deliverables and Timeline}
\label{app:deliverables}
%==============================================================================

\section{Internship timeline at a glance}

\begin{longtable}{@{}p{0.16\textwidth}p{0.80\textwidth}@{}}
\toprule
\textbf{Dates (2026)} & \textbf{Work} \\
\midrule
\endhead
May 15 -- Jun 2  & Onboarding; literature review (LoRA/W2T, efficient benchmarking).
    Thread B begins. \\
Jun 3 -- Jun 4   & First commits: e2e adapter config and tests (Chapter~\ref{chap:e2e}). \\
Jun 9            & Corpus cartography / filter statistics (Chapter~\ref{chap:cartography}). \\
Jun 11 -- Jun 12 & Repo refactor (Phases 0--2) and the unified comparison core
    (Phase 3 / \code{NormalizedDiff}, judge registry) (Chapter~\ref{chap:core}). \\
Jun 12 -- Jun 18 & OLMo campaign begins: tokenizers, chat-template window budget,
    the 114-packet canonical-key fix (Chapter~\ref{chap:olmo});
    Thread~B \code{dkps\_mrmr} lands (Jun 24). \\
Jun 16 -- Jun 24 & Containerized HELM execution and the infer-stack leasing pipeline
    (Chapter~\ref{chap:containers}). \\
Jun 25 -- Jun 30 & From-spec replay: run-spec verbatim replay, deployment rewrite,
    rel-path addressing, content-addressed caching (Chapter~\ref{chap:fromspec}). \\
Jun 30 -- Jul 10 & Deployment matching: OLMo-7B EOS-append, the deployment-match
    tool, \code{add\_generation\_prompt} drift, the float32 discovery, the
    unrecorded-parameter taxonomy (Chapter~\ref{chap:deploymatch}). \\
Jul 2 / 6 / 12   & Three codebase audit / simplification passes
    (Chapter~\ref{chap:stewardship}). \\
Jul 8 -- Jul 13  & Era-pinned containers, the shim, redpajama-3b runbook, classic
    Together models, the G13 archaeology (Chapter~\ref{chap:eras}). \\
Jul 11 -- Jul 14 & GPT-OSS-20B and Qwen fan-outs (Chapter~\ref{chap:newmodels});
    aggregate score-drift reporting fixes, route registry
    (Chapter~\ref{chap:reporting}). \\
\midrule
\multicolumn{2}{@{}l}{\emph{Evidence cutoff of the 2026-07-15c revision. Work below
    is largely Jon Crall's; see \S\ref{sec:rev-authorship}.}} \\
\midrule
Jul 15 -- Jul 17 & The Qwen3.5 extension: \textsc{helm} registry sidecars, the
    newline-tolerant declared-substitution client, the compile-cache wart (G14), the
    compute-versus-from-spec grid bug (G15), VRAM-aware placement, and the
    expand-once-then-freeze self-critique (Chapter~\ref{chap:qwen35}). \\
Jul 17 -- Jul 21 & The open-judge experiment: response snapshots, the
    identity-replay gate (6/6 benchmarks), the live smoke, the v1 full run
    (2 judges $\times$ 2 benchmarks $\times$ 3 replicates, $\sim$20.3~h), the
    judge-size ladder, and the contamination caveat
    (Chapter~\ref{chap:openjudge}). \\
Jul 21 & Three adversarial rounds on the paper thesis; the zero-compute
    pairwise-flip analysis; the pre-registered overnight protocol
    (Chapter~\ref{chap:paperdirection}, Chapter~\ref{chap:substrate-closed}). \\
Jul 21 -- Jul 23 & Substrate recovery closed: three failed nights diagnosed, the
    float32 end-to-end confirm (prediction refuted), the byte-exact Hugging~Face
    float32 reproduction, and the pure-engine attribution
    (Chapter~\ref{chap:substrate-closed}). \\
Jul 23 & Cross-family generality tooling and the task queue that turns one cell into
    a systematic result. \\
Jul 27 & TMLR manuscript: introduction and related-work rewrite, Goodman's
    trichotomy credited, the taxonomy claim narrowed, the eligibility gate made
    target-relative, and the bibliography verified against sources
    (\S\ref{sec:paper-drafting}). \\
\bottomrule
\end{longtable}

By commit count: 452 authored commits through the July 14 cutoff, distributed as 181
in June and 271 in July; the busiest days were 2026-07-06 (50 commits, the simplicity
audit) and 2026-07-12 (39, the fourth simplification pass). The continuation window
adds \textbf{127 commits} (2026-07-16 through 2026-07-27), of which 107 are Jon
Crall's and 20 are Edward Wang's; its busiest days were 2026-07-17 (25, open-judge
Milestone~A), 2026-07-16 (22, the Qwen3.5 arc) and 2026-07-21 (22, the thesis rounds
and the overnight protocol).

\section{Key artifacts produced}

\paragraph{Reproduction runbooks} (\file{reproduce/}, \file{dev/}):
\code{olmo\_models\_combined}, \code{gpt\_oss\_20b\_from\_spec},
\code{qwen\_models\_combined}, \code{classic\_together\_combined},
\code{dev/e2e-tests} (phi-2), \code{dev/era-tests} (redpajama-3b). \emph{Added in
the continuation window}: \code{qwen35\_vllm} and \code{qwen35\_small\_vllm} (the
first \emph{compute} runbooks --- new models rather than replays),
\code{open\_judge\_gpt\_oss} (audit $\rightarrow$ snapshot $\rightarrow$ identity
replay $\rightarrow$ rejudge $\rightarrow$ analyze), and the
\code{deployment\_match} scripts \code{60}--\code{73} (float32 end-to-end confirm,
Hugging~Face float32 config probe, cross-family sweep).

\paragraph{Tools and libraries}: \code{dev/tools/deployment\_match/} (the
serve-and-probe deployment matcher with \code{hf-probe}, \code{hf-match},
fp32-TP, and \code{compare\_prompt.py}); \code{docker/era\_shim/} (the standalone
era from-spec replay CLI + OpenAI-compatible client); \code{eval\_audit/eras.py}
and \code{docker/eras.yaml} (era registry); \code{eval\_audit/hf\_inprocess.py}
(reserve-GPU in-process reproduction); \code{eval\_audit/normalized/diff.py}
(\code{NormalizedDiff}, the unified comparison core);
\code{eval\_audit/judge\_registry.py} and \code{eval\_audit/metrics\_taxonomy.py}.
\emph{Added in the continuation window}: the \code{eval\_audit/judging/} package
(content-addressed response snapshots, \code{JudgeSpec}/\code{JudgmentAttemptSpec},
configurable annotators for six benchmarks, the annotation-only rejudge runner,
judge-attributed metrics, indexing and variance analysis) with its CLIs
(\code{eval-audit-audit-judge-sources}, \code{eval-audit-verify-judge-replay},
\code{eval-audit-judge-prompt-lengths}, \code{eval-audit-analyze-judges},
\code{eval-audit-schedule-rejudge}); \textsc{helm} registry-sidecar passthrough
through the manifest, bridge, container mounts, and bundle exporter; and the
newline-tolerant completions clients.

\paragraph{Documentation} (the intended paper appendix material):
\file{helm-gotchas.md} (G1--G13; G14--G15 added 2026-07-27), \file{vllm-vs-huggingface-deployment-match.md},
\file{helm-unrecorded-deployment-params.md}, \file{eee-vs-helm-metadata.md},
\file{pipeline.md}, \file{architecture.md}, \file{container-execution.md}, and the
planning documents under \file{docs/planning/} --- to which the continuation window
added \file{open-judge-plan.md}, \file{tmlr-paper-thesis.md},
\file{compute-run-spec-freeze-plan.md}, and \file{edward-task-queue.md}, plus the
pre-registered
\file{reproduce/olmo\_models\_combined/deployment\_match/overnight\_confirm\_plan.md}
and the TMLR manuscript sources under \file{docs/papers/tmlr-2026/}.

\paragraph{Thread~B} (sibling repos): \file{mrmr\_eval} (the \code{benchpred}
efficient-benchmarking harness with the \code{dkps}/\code{dkps\_mrmr} methods and the
\code{med\_qa} report) and \file{dkps} (the vendored DKPS library and HELM example
runners).

\section{A note on method and provenance of this document}

This chronology was reconstructed after the fact from the primary sources listed in
Chapter~\ref{chap:framing}. Where numbers are quoted (agreement ratios, token
counts, corpus statistics, the fp32 sweep table, the efficient-benchmarking results)
they are taken directly from the status decks, the developer journals, or the
committed reports; where a claim is a design decision or a lesson, it is drawn from
the corresponding journal entry or planning document. The development was carried
out with heavy use of an AI pair-programming harness (recorded in the first-person
developer journals); the research direction, design decisions, GPU-host execution,
and interpretation are the author's. Readers preparing the academic manuscript
should treat the \S-cross-referenced figures and tables here as pointers into those
primary artifacts, which remain the ground truth.

\emph{On the 2026-07-27 continuation.} Chapters~\ref{chap:qwen35}--\ref{chap:paperdirection}
were written from the same class of primary sources for the July 15--27 window: the
developer journals (which for this window are unusually detailed, including
registered predictions and same-session corrections), the Git history, the
pre-registered plan files, and the planning documents. Two provenance caveats apply
specifically to them. First, the numeric results are read from the journals'
records of runs executed on Kitware GPU hosts; the underlying result stores were
\emph{not} re-inspected while writing this revision, and --- per
\S\ref{sec:consensus-closed} --- none of them is preserved or hashed, so every figure
in Chapter~\ref{chap:substrate-closed} and Chapter~\ref{chap:openjudge} should be
regenerated from artifacts before it appears in a manuscript. Second, the work in
this window was led by a second engineer (\S\ref{sec:rev-authorship}), so the
first-person framing of the earlier chapters does not carry over.

%==============================================================================
\chapter{Interpretive Revisions (2026-07-15, 2026-07-27): Epistemic Status of the Central Claims}
\label{app:revision}
%==============================================================================

The body of this chronology reconstructs, in a single after-the-fact narrative
voice, what was believed and reported across the internship. Some of those claims
--- especially around the \code{float32} result --- were stated more strongly than
the evidence currently supports. Rather than rewrite the dated narrative (which
would erase what was believed at each point), this appendix records the corrected
epistemic status in one place. It was prompted by an external review of the
chronology and the paper plan, and its conclusions are folded into
\path{docs/planning/tmlr-paper-plan.md}.

\section{A workflow-status vocabulary}

Claims in the body should be read against \emph{where in the pipeline} their
evidence sits. These are not interchangeable, and the body sometimes chose the
strongest available interpretation:

\begin{enumerate}[nosep]
  \item \textbf{Infrastructure implemented} --- code exists, unit-tested on a CPU host.
  \item \textbf{Smoke-tested execution} --- ran end-to-end on a tiny input, once.
  \item \textbf{Deployment-match probe} --- an oracle-scored config sweep over a small
        sample, possibly using non-default request knobs. \emph{Discovery, not
        confirmation.}
  \item \textbf{Exact-spec run} --- the official \file{run\_spec.json} executed
        verbatim on the real path.
  \item \textbf{Complete \textsc{helm} artifact generation} --- normal
        \file{scenario\_state.json}/\file{stats.json} produced.
  \item \textbf{Final comparison report} --- run through the standard pair-report
        pipeline.
  \item \textbf{Fresh vs.\ stale store} --- regenerated with current code vs.\ a store
        predating a relevant fix.
\end{enumerate}

The single most important correction: the \code{float32} result (\S\ref{sec:fp32})
sits at level~3 (probe / discovery), not levels~4--6. The ``confirm'' step that
would move it to a reproduction result --- land the winning config on the ordinary
\textsc{helm} path, run the exact public \file{run\_spec.json}, generate normal
artifacts, compare via pair-report, and evaluate on \emph{held-out} instances --- has
not been run for any model. Configuration discovery and held-out confirmation must be
separated, or a config fit to a small sample is reported as if it were an identified
fact.

\section{A finer disagreement taxonomy}

\S\ref{sec:distinction}'s binary --- ``recipe/environment failure'' vs.\ ``true
reproducibility failure'' --- is too coarse, and the phrase ``true reproducibility
failure'' should be \emph{avoided} whenever execution provenance is incomplete. Keep
the upstream run-vs-not-run gate (non-runnable = a filtering/eligibility reason, not
a disagreement), and for observed disagreements use:

\begin{enumerate}[nosep]
  \item \textbf{Recipe mismatch} --- the prompts/decoding/scenario differ (e.g.\ the
        BBQ \code{output\_format\_instructions} drift, G5).
  \item \textbf{Deployment mismatch} --- serving engine / client class differs (vLLM
        vs.\ hosted API; \code{HuggingFaceClient} vs.\ vLLM).
  \item \textbf{Execution-instrument mismatch} --- precision, tokenizer/template
        version, attention kernel, device placement, software-stack versions.
  \item \textbf{Artifact-interpretation or migration mismatch} --- the pandas
        row-order / instance-id case (EEE Case~A); the G13 class-path migration.
  \item \textbf{Residual disagreement under controlled equivalence} --- the only
        category that resembles conventional repeatability failure: everything
        controllable is matched and a gap remains.
  \item \textbf{Historically non-identifiable reproduction} --- the surviving
        artifacts and releases do not uniquely specify what faithful reproduction
        would even mean (GPT-J / GPT-NeoX / OPT via G13). Qualitatively different from
        (5): not ``it failed to repeat'' but ``the target is underdetermined.''
\end{enumerate}

Most of what the body called ``deployment-boundary artifacts'' are categories
(2)--(4) --- recoverable by controlling the substrate. Category (6) is the negative
counterpart and is a first-class result, not a gap to be closed.

\section{Corrected-claims table}

\begin{longtable}{@{}p{0.42\textwidth}p{0.54\textwidth}@{}}
\toprule
\textbf{As stated in the body} & \textbf{Corrected status} \\
\midrule
\endhead
``The official OLMoE run --- and every OLMo HF-client run --- executed in
\code{float32}.'' & The \emph{loader code path} defaults to fp32 under every
\code{transformers} version consistent with the run date (a deductive claim, the
stronger leg), \emph{conditional} on the unrecorded \code{transformers}$<$5. Likely,
not proven. \\
``fp32 reproduces the official completions \emph{exactly}.'' & Among tested configs,
fp32 candidates \emph{best matched the sampled} public outputs ($n\approx12$ ifeval),
using a probe-only \code{add\_special\_tokens=False} knob. Not an end-to-end HELM
reproduction; not held-out. \\
``all four officials reproduce at fp32.'' & Probe-level, ifeval-only, instruct-only.
OLMoE via HF fp32; dense OLMo-2 via vLLM fp32; the HF probe does \emph{not} reproduce
OLMo-2 (an open execution-sensitivity finding, below). \\
``Precision is the dominant hidden variable.'' & Supported \emph{within} the OLMo
ifeval pilot; not established across benchmarks or model families. \\
``the OLMo-2 HF divergence is the probe mis-generating, not science.'' & Reframed: a
\emph{current} HF stack failing to reproduce a historical result \emph{recorded as a
\code{HuggingFaceClient} deployment} under byte-identical prompt tokens is itself an
unresolved systematic divergence under prompt-token equivalence --- evidence for the
thesis (with a dense-model probe artifact still on the candidate-cause list). \\
``residual drift is small, concentrated, and almost always attributable.'' & A
motivating \emph{hypothesis}, pending the confirm step, fresh stores, cross-machine
data, and a controlled ablation. Not yet established. \\
OLMo/GPT-OSS heatmaps ``exist''. & \emph{Split by store (collaborator-verified
2026-07-15).} GPT-OSS report files were regenerated Jul 14 at 13:47 after reporting
fixes, but the run and comparison provenance must still pass the preservation and
acceptance gate. \code{olmo-models-combined} is stale: the reported on-disk headline
still shows the un-deduped \code{olmo-7b} halving (MMLU \textbf{0.295/0.144}, GSM
0.036/0.018, narrative\_qa 0.597/0.311), whereas the provisional corrected value cited
in the figure is 0.295/0.287. Regenerate before citing any OLMo base-model number. \\
\bottomrule
\end{longtable}

None of these corrections weaken the \emph{methodological} contributions (from-spec
replay, era-pinned instruments, sweep-don't-guess attribution, the provenance
taxonomy) or the \emph{conceptual} contribution (recoverable vs.\ non-identifiable
provenance). They scope the \emph{empirical} claims to what has actually been run.

\section{Second interpretive pass (2026-07-27)}
\label{sec:revision-2}

The 2026-07-15 pass scoped claims \emph{down} to the evidence in hand. This pass
records what the subsequent experiments did to those scoped claims --- three of them
move up, one moves down, and one turns out to have been wrong in a way that is worth
naming precisely.

\begin{longtable}{@{}p{0.30\textwidth}p{0.16\textwidth}p{0.50\textwidth}@{}}
\toprule
\textbf{Claim (as scoped on 2026-07-15)} & \textbf{Movement} & \textbf{Basis at this
revision} \\
\midrule
\endhead
The fp32 result is configuration \emph{discovery} at workflow level~3, awaiting the
confirm step. & \textbf{Up}, to levels 4--6 for one cell & The exact public run spec
was replayed on the ordinary path at full $n$, normal artifacts were produced, and
the pair comparison was run (Table~\ref{tab:fp32e2e}); a separate probe reproduces
the official completions byte-exactly (Table~\ref{tab:hffp32}). Scope: OLMo-2 7B and
13B instruct, \code{ifeval} only. \\
``Precision is the dominant hidden variable,'' supported within the OLMo pilot. &
\textbf{Down / refuted as stated} & Precision plus template rendering closes only
$\approx$1/3 of the drift. The registered prediction that float32 would recover the
official was refuted before the engine result arrived. The honest replacement:
\emph{substrate recovery is layered, and no single variable is dominant}. \\
The dense OLMo-2 Hugging Face divergence is an unresolved systematic divergence
under prompt-token equivalence. & \textbf{Withdrawn} & It was the dense-model probe
artifact the 15c review kept on the candidate list: unpinned template rendering and
device sharding in our own sweep. With the axes pinned, current
\code{transformers.generate()} at float32 reproduces the official exactly
(\S\ref{sec:residual-resolved}). \\
Execution-stack sensitivity is real and consequential. & \textbf{Up, on different
evidence} & Now measured at the published-metric level as a
\textsc{vllm}$\leftrightarrow$HF engine gap of $+0.067$/$+0.082$ at fixed weights,
precision, prompt, and decode --- rather than inferred from a divergence that turned
out to be ours. \\
An unavailable judge is an eligibility filter. & \textbf{Sharpened} & It filters a
run out of \emph{procedural} reproduction only. The per-instance verdicts survive,
so artifact reconstruction remains fully available and was executed at machine
precision on six benchmarks (Table~\ref{tab:replaygate}). \\
``residual drift is small, concentrated, and almost always attributable.'' &
\textbf{Unchanged: still a hypothesis} & One cell is now fully attributed, which is
an existence proof rather than a prevalence claim. The frozen-denominator census
that would test prevalence is specified and unstarted. \\
Store status (GPT-OSS fresh reports; OLMo stale). & \textbf{Sharpened, and worse
than assumed} & The 2026-07-15d disk audit found both stores pruned to derived
layers with no raw run artifacts, so neither can be preserved by copy-and-hash;
\emph{re-run} is required. The July 21--23 substrate artifacts are new and equally
unpreserved. \\
\bottomrule
\end{longtable}

\noindent Two claims are \emph{new} at this revision and should be read with their
scope attached: that the inference engine is a first-class hidden variable
(established on one cell, characterized but not explained), and that judge
non-determinism is universal while metric fragility is metric-dependent (established
across six benchmarks and six judge arms, but on a single candidate model and under
a contamination caveat that makes every agreement figure an upper bound).

\noindent One methodological correction applies to this document's own instrument.
Two of the erroneous conclusions in this record --- the dense OLMo-2 divergence and a
cross-family sweep that appeared to refute its prediction --- were produced by
\emph{our} tools silently defaulting an unspecified setting: an unpinned sweep axis
in the first case, a protocol resolver that defaulted on failure in the second. A
project whose thesis is that unrecorded execution parameters corrupt published
numbers should hold its own instruments to the standard it recommends; both cases
are documented in place rather than quietly fixed, because the failure mode is the
finding.



%==============================================================================
\chapter{Collaborative Claim--Evidence--Status Matrix}
\label{app:master-claim-matrix}

\noindent\textbf{Status note.} This appendix is a human-readable snapshot. The
canonical, machine-readable source is
\file{chronicle/data/master\_claim\_evidence\_ledger\_2026-07-27.csv}; update that
file first and regenerate this table before submission. Rows marked
\emph{(rev.\ 2026-07-27)} changed in the continuation revision; see
\S\ref{sec:revision-2}.

%==============================================================================

\begin{longtable}{@{}p{0.35\textwidth}p{0.17\textwidth}p{0.40\textwidth}@{}}
\toprule
\textbf{Claim} & \textbf{Status} & \textbf{Qualification / principal evidence} \\
\midrule
\endhead
Name-only matching can misclassify comparable and incomparable runs. &
\statusEstablished & Early \code{groups=} and omitted-temperature cases; exact run
specs later became the authoritative recipe source. \\
The inherited EEE study detects gaps but does not attribute the residual substrate. &
\statusEstablished & NeurIPS Case Study 3 and its appendix; Pythia/Vicuna/Falcon
scope predates the internship. \\
Exact public \file{run\_spec.json} replay removes observed prompt-reconstruction
drift. & \statusEstablished & BBQ output-format instruction case and the June 25
from-spec implementation. \\
OLMo-7B prompt-independent gibberish was caused by appended special-token behavior. &
\statusEstablished & Token-level prompt diagnosis and tokenizer override. \\
IFEval output is sensitive to chat-template and
\code{add\_generation\_prompt} behavior. & \statusEstablished & Controlled prompt
rendering and deployment-match sweeps. \\
Float32 candidates best matched sampled OLMo instruct outputs in the 12-instance
IFEval probe. & \statusEstablished{} \emph{(rev.\ 2026-07-27)} & Confirmed and
superseded by a stronger result. Ordinary-path replay at full $n$ plus a byte-exact
Hugging~Face float32 reproduction (Chapter~\ref{chap:substrate-closed}) establish
float32 as the official load precision for OLMo-2 instruct, and show it to be
necessary but not sufficient: precision and template rendering close
$\approx$1/3 of the drift, the engine closes the rest. Scope: one benchmark, one
family, two sizes. \\
The published OLMo-2 instruct IFEval number is exactly reproducible once precision,
template rendering, and engine are recovered. & \statusEstablished{} \emph{(rev.\
2026-07-27)} & $32/32$ byte-identical completions on 7B and 13B across all four
forward-pass cells. Completion-level evidence; the metric claim follows by
implication, not by a metric-level in-process run. \\
The \textsc{vllm}$\leftrightarrow$Hugging~Face engine gap moves a published metric at
fixed weights, precision, prompt, and decode. & \statusEstablished{} \emph{(rev.\
2026-07-27)} & $+0.067$ (7B) and $+0.082$ (13B) on \code{ifeval\_strict\_accuracy};
decode semantics excluded by a greedy-versus-\textsc{helm}-decode probe that produced
byte-identical completions. Characterized, not mechanistically explained. \\
Open-weight judges recover the published conclusions of model-graded benchmarks. &
\statusPreliminary{} \emph{(rev.\ 2026-07-27)} & Agreement on label metrics exceeds
the official ensemble's internal agreement, and run-level rubric means are stable;
but all artifacts score a \emph{single} candidate, 43--46\% of rubric instances
change score across replicates, and every figure is an upper bound under the
contamination caveat (\S\ref{sec:contamination}). \\
Judge non-determinism at temperature~0 is universal; metric fragility is
metric-dependent. & \statusEstablished{} \emph{(rev.\ 2026-07-27)} & 87--96\% of
judgments differ in raw text across replicates; score impact 0.2--0.7\% on a
near-binary label versus 43--46\% on a 1--10 rubric. Verified not to be an artifact
of our own per-replicate cache key. \\
Reconstructed snapshots reproduce published judge metrics. & \statusEstablished{}
\emph{(rev.\ 2026-07-27)} & Identity-replay gate, six of six benchmarks, max error
$0$ to $1.95\times10^{-14}$ (Table~\ref{tab:replaygate}). \\
A net-new model can be added without modifying \textsc{helm}. & \statusImplemented{}
\emph{(rev.\ 2026-07-27)} & Registry-sidecar passthrough, proven end to end on GPUs
for Qwen3.5-9B-Base; no \textsc{helm} edit, no image rebuild
(Chapter~\ref{chap:qwen35}). No Qwen3.5 report has been produced. \\
Our own compute runs follow the freeze-the-spec discipline we recommend. &
\statusPlanned{} \emph{(rev.\ 2026-07-27)} & They do not: a mutable run-key string
is still the stored source of truth. \file{compute-run-spec-freeze-plan.md} is
proposed and unimplemented (\S\ref{sec:freeze-critique}). \\
The public OLMoE load most likely used fp32. & \statusPreliminary & Conditional
source-level inference: unpinned dtype, HuggingFaceClient path, bounded Transformers
4.x behavior. Exact historical package environment is absent. \\
The dense OLMo-2 current-HF divergence has a known single cause. &
\statusResolved{} \emph{(rev.\ 2026-07-27)} & Resolved as a probe artifact: unpinned
template rendering and \code{device\_map=auto} sharding in our own sweep. With those
axes pinned, current \code{transformers.generate()} at float32 reproduces the
official byte-exactly in all four forward-pass cells
(\S\ref{sec:residual-resolved}). \\
RedPajama-3B ran through both later classic proxy eras. & \statusLocalOnly &
Journal/slide and report-refresh references support local completion; raw
\path{/data} reports are not included in this source archive. \\
GPT-J/GPT-NeoX/OPT public artifacts identify a unique released producing harness. &
\statusUnresolved & Evidence supports the opposite: migrated class paths point to an
unreleased source interval and later v0.2.4/v0.3.0 images are declared proxies. \\
Public classic artifacts were carried forward between v0.2.4 and v0.3.0. &
\statusEstablished & Byte-identical public artifacts for shared runs; not two
independent officials. \\
The compatible classic source interval is July 31--August 26, 2022. &
\statusEstablished & Source-history triangulation under the migration interpretation;
it is an inferred compatibility interval, not an observed run timestamp. \\
The complete OLMo/Qwen/GPT-OSS/classic performance grids are publication-ready. &
\statusUnresolved{} \emph{(rev.\ 2026-07-27: worse than assessed)} & Stores and
workflows exist at different levels, and the 2026-07-15d disk audit found the two
flagship modern stores pruned to derived layers with no raw run artifacts --- so
GPT-OSS and \code{olmo-models-combined} require a \emph{re-run}, not a report
refresh. The July 21--23 substrate artifacts are new and unpreserved. \\
A position/systems paper is supportable now. & \statusEstablished & The conceptual
model, forensic cases, system implementation, and provenance recommendations are
substantial if empirical claims remain scoped. \\
A broad technical claim that most residual gaps are small and attributable is
supportable now. & \statusPlanned & Requires a frozen denominator, fresh stores,
ordinary-path confirmation, held-out validation, ablations, repeatability baselines,
and archived artifacts. \\
\bottomrule
\end{longtable}

\chapter{Source Inventory and Reproducibility Notes}

\begin{itemize}
  \item Repository archive: \path{eval_audit-source-2026-07-14T144915-5-b8b858697e7c.tar.gz}; supplied HEAD \code{b8b858697e7c2491af163c8f5d757fe6ae6f3c1a}.
  \item Weekly slide archive: \path{internship-powerpoints.zip}, seven decks from June 2 through July 14.
  \item Detailed Claude chronology: preserved under \path{sources/claude_chronology_2026-07-15.tex}.
  \item Independent evidence-led chronology: preserved under \path{sources/openai_chronicle_2026-07-15.tex}.
  \item Machine-readable commit and slide ledgers are included under \path{data/}.
  \item The source archive manifest excludes untracked, ignored, and \path{/data}
        result stores. Any claim based on those stores must be accompanied by a
        separately preserved artifact bundle before publication.
\end{itemize}

\noindent\textbf{Added for the 2026-07-27 continuation revision.}
\begin{itemize}
  \item Consensus acceptance, furnished source hashes, and the disk-level store
        provenance audit: \path{validation/consensus_accepted_2026-07-15d/}
        (letter, \path{SOURCE_HASHES_claude_2026-07-15d.txt},
        \path{store_provenance_audit_2026-07-15d.csv}). The furnished immutable
        reference is commit
        \code{86ec84af09c0b8f4442afa1ce87b56e1f8b3dc61} (tree
        \code{a3734de5b0178077bf241b6026e7d68d74247853}); no Git bundle was exported,
        because the repository is private.
  \item Developer journal entries 2026-07-15 through 2026-07-23
        (\path{dev/journals/claude.md}) and the 2026-07-17 entry of
        \path{dev/journals/codex.md}. These are the primary source for
        Chapters~\ref{chap:qwen35}--\ref{chap:paperdirection}.
  \item Pre-registered experiment protocol:
        \path{reproduce/olmo_models_combined/deployment_match/overnight_confirm_plan.md}
        (predictions and outcome classes frozen before execution), and the
        substrate-recovery section of
        \path{docs/helm-reproduction-research-journal.md}.
  \item Planning documents added in the window: \path{docs/planning/open-judge-plan.md},
        \path{tmlr-paper-thesis.md}, \path{compute-run-spec-freeze-plan.md},
        \path{edward-task-queue.md}.
  \item Manuscript sources: \path{docs/papers/tmlr-2026/}.
  \item Git history for the window: 127 commits, 2026-07-16 through 2026-07-27, on
        \code{main} / \code{jons/qwen35-extension} / \code{impl/run-from-run-spec}.
        \textbf{No commit ledger for this window is packaged here}; the appendices
        under \path{chronicle/appendices/} still cover only the period through
        2026-07-14 and should be regenerated before submission.
  \item \textbf{Not packaged, and load-bearing}: every result store cited in
        Chapters~\ref{chap:openjudge} and~\ref{chap:substrate-closed}. The
        substrate and rejudge artifacts live on Kitware GPU hosts and a workstation
        mirror; none has been copied, hashed, or bundled.
\end{itemize}

\chapter{Publication-Grade Preservation Checklist}

A paper artifact should preserve source commit and dirty state; container digests;
Python, Transformers, PyTorch, CUDA and driver versions; GPU model and topology;
model and tokenizer commit SHAs; effective dtype, quantization, attention backend,
TF32/matmul flags, batching, caching, and tensor parallelism; raw, rendered, and
tokenized prompts; complete generation configuration; frozen public run specs and
hashes; raw local and public outputs; every deployment-search candidate rather than
only the winner; compatibility rewrites; report-generation version; seeds and
repeated-run manifests; and all analysis/figure scripts.

\noindent Three additions from the continuation revision, each earned by a specific
failure or result. \textbf{(i)~The serving engine and its build}, not only the
client family --- it is worth $+0.07$ on a published metric at fixed everything else
(Chapter~\ref{chap:substrate-closed}). \textbf{(ii)~Compiled-artifact and
kernel-cache provenance}, because a compile cache keyed more narrowly than the
serving-config space can silently supply artifacts built under a different
configuration (G14). \textbf{(iii)~For model-graded runs}, the judge model, the
ensemble composition, the harness version that fixes them, and the judge's decoding
budget --- plus an explicit note that a hosted judge's batch composition is
\emph{not} recoverable, so replicate counts and observed judgment variance should be
preserved in place of a reproducibility promise.

\noindent A procedural addition of the same weight: \textbf{preserve the raw run
artifacts, not only the derived layers}. The 2026-07-15d audit found the two
freshest report stores sitting on pruned inputs, which is the failure this checklist
exists to prevent and the reason a report timestamp must never serve as an
acceptance criterion.

% Guarded so the master document compiles standalone even when the commit-ledger
% appendices are not shipped alongside it (packaging fix, 2026-07-15b).
\IfFileExists{appendices/main_commit_ledger.tex}{\chapter{Main-repository commit ledger}\label{app:main-commit-ledger}
\small
This appendix lists every commit authored as \texttt{Edward Wang <edward.wang@kitware.com>} in the supplied \texttt{eval\_audit} history from 2026-06-03 through 2026-07-14. Insertions and deletions are Git numstat totals and should not be interpreted as net authored lines: they include refactors, moves, generated fixtures, and regenerated lock files.
\section*{2026-06-03 (1 commits)}
\begin{longtable}{@{}p{0.105\textwidth}p{0.665\textwidth}r r r@{}}
\toprule Hash & Subject & + & - & Files \\ \midrule\endhead
\texttt{ca46b11e} & Add adapter config for the e2e test and update matching logic & 319 & 74 & 2 \\
\bottomrule\end{longtable}
\section*{2026-06-04 (1 commits)}
\begin{longtable}{@{}p{0.105\textwidth}p{0.665\textwidth}r r r@{}}
\toprule Hash & Subject & + & - & Files \\ \midrule\endhead
\texttt{ed9e40ed} & Add e2e tests & 263 & 0 & 6 \\
\bottomrule\end{longtable}
\section*{2026-06-11 (19 commits)}
\begin{longtable}{@{}p{0.105\textwidth}p{0.665\textwidth}r r r@{}}
\toprule Hash & Subject & + & - & Files \\ \midrule\endhead
\texttt{609e02bb} & Fix stale prioritized-example repair test to match no-regeneration contract & 48 & 34 & 1 \\
\texttt{59fbdc2d} & Add repo refactor plan and journal entry & 355 & 0 & 2 \\
\texttt{f3953650} & Phase 0a: clean up repo root (refactor plan) & 18 & 12 & 5 \\
\texttt{6f294d03} & Phase 0b: finish vllm\_service -> infer\_stack rename & 2 & 2 & 3 \\
\texttt{3507539d} & Phase 0c: consolidate root docs & 5 & 2 & 4 \\
\texttt{4a0ae5b7} & Phase 1: every console script resolves to eval\_audit.cli.* & 71 & 12 & 9 \\
\texttt{6d5035aa} & Phase 1: retire grouped-CLI dispatchers and dead alias module & 34 & 37 & 7 \\
\texttt{4473bfc9} & Phase 1: role docstrings for core\_* modules; fix stale runbook command & 46 & 1 & 6 \\
\texttt{b3a4b272} & Phase 2 prep: hoist @profile shim into infra/profiling.py & 44 & 102 & 15 \\
\texttt{7f64806b} & Phase 2: split build\_reports\_summary into reports/summary/ subpackage & 4570 & 4235 & 12 \\
\texttt{72bfeb3d} & Journal: refactor plan implementation session (Phases 0-2) & 85 & 0 & 1 \\
\texttt{9b7ddb74} & Phase 2: extract Stage 1 library logic from cli/index\_historic\_helm\_runs & 639 & 568 & 3 \\
\texttt{ca6f5d7b} & Phase 2: extract diff primitives from helm/diff.py & 715 & 668 & 2 \\
\texttt{dd742728} & Phase 2: split core\_metrics into curves / plots / tables modules & 2165 & 2003 & 5 \\
\texttt{9a377838} & Phase 2: split filter\_analysis into tables / text / charts / io modules & 1313 & 1206 & 5 \\
\texttt{c201938f} & Phase 2: split eee\_only\_heatmap into data / render modules & 1243 & 1173 & 3 \\
\texttt{9fcdfefe} & Phase 2: split helm/analysis into instance\_stats / analysis\_report & 1075 & 1028 & 3 \\
\texttt{3822d676} & Mark Phase 2 secondary targets done in plan; journal entry & 91 & 6 & 2 \\
\texttt{5b05562f} & Add Phase 3 design docs (comparison-core unification + equivalence matrix) & 522 & 27 & 3 \\
\bottomrule\end{longtable}
\section*{2026-06-12 (28 commits)}
\begin{longtable}{@{}p{0.105\textwidth}p{0.665\textwidth}r r r@{}}
\toprule Hash & Subject & + & - & Files \\ \midrule\endhead
\texttt{76a320c3} & Revise Phase 3 design docs for the new research program & 389 & 261 & 3 \\
\texttt{352d55d0} & Journal: Phase 3 design + research-program revision & 73 & 0 & 1 \\
\texttt{94c7fb49} & Phase 3 / 4.0: lift metric taxonomy out of helm/, add judge-dependence class & 283 & 115 & 4 \\
\texttt{4a36e6d4} & Phase 3 / 4.1: recipe\_facts accessor + judge-identity fields & 453 & 1 & 3 \\
\texttt{1836be3e} & Phase 3 / 4.2: normalized/diagnose.py -- framework-free diagnosis port & 559 & 0 & 2 \\
\texttt{6656048b} & Phase 3: behavior baseline capture harness + committed F3/F4 snapshots & 1919 & 0 & 5 \\
\texttt{426a211a} & Phase 3 / 4.4: EEE synthesis library home + zero-HELM EEE entry points & 507 & 436 & 11 \\
\texttt{12623aa5} & Journal: Phase 3 implementation session (4.0-4.2, baseline, 4.4) & 75 & 0 & 1 \\
\texttt{a347eca0} & Phase 3 / 4.3: NormalizedDiff -- the unified comparison core (additive) & 616 & 102 & 3 \\
\texttt{cc6c9e8b} & Phase 3: judge-identity inventory across the closed-judge benchmarks & 86 & 0 & 1 \\
\texttt{dfe2c941} & Journal: Phase 3 continuation (4.3 + judge inventory) & 62 & 0 & 1 \\
\texttt{5c8f7a48} & Phase 3 / 4.6: route the renderer through NormalizedDiff; single diagnosis impl & 59 & 271 & 3 \\
\texttt{9e40aa92} & Phase 3 / 4.5 (loader half): declared instance-source policy + F6 probe & 206 & 22 & 2 \\
\texttt{e5f905c7} & Phase 3 / 4.5 (renderer half): --instance-source flag + recorded provenance & 75 & 0 & 6 \\
\texttt{0299ed03} & Journal + design-doc status: Phase 3 sub-stages 4.0-4.6 complete & 75 & 0 & 2 \\
\texttt{3ef5233e} & Phase 3 / 4.9a: curated judge registry + identity resolution & 166 & 3 & 3 \\
\texttt{900ab05c} & Phase 3 / 4.9b+c: declared judge substitutions through planner + renderer & 363 & 18 & 5 \\
\texttt{5b34364f} & Phase 3 / 4.9d: Stage-1 --allow-closed-judge-benchmarks relax & 147 & 13 & 3 \\
\texttt{ca83c051} & Phase 3 / 4.7: draft the upstream EEE recipe\_facts issue & 89 & 0 & 1 \\
\texttt{f099bb2a} & Phase 3 / 4.8: docs pass + legacy-bridge marking + final status & 100 & 27 & 5 \\
\texttt{ac3445ee} & Journal: Phase 3 completed (4.9 + 4.7 draft + 4.8) & 69 & 0 & 1 \\
\texttt{aaa2c928} & Fix infer\_stack adapter: drop removed root= kwarg from load\_profile\_contract & 6 & 1 & 1 \\
\texttt{bcb0e16f} & Add OLMo reproduction smoke suite: adapter presets, run scripts, and virtual experiment config & 1385 & 0 & 14 \\
\texttt{ec4960aa} & Update deprecated hf auth call & 1 & 1 & 1 \\
\texttt{92f9582a} & Switch profiles without prompting & 1 & 1 & 1 \\
\texttt{31961a84} & Fix OLMo tokenizer aliases to match HELM registry & 3 & 3 & 1 \\
\texttt{517cbbf5} & Add sentencepiece/tiktoken tokenizer backends as deps & 6 & 0 & 1 \\
\texttt{2aea115d} & Fix olmo-7b tokenizer load via a HELM entry-point plugin & 58 & 6 & 2 \\
\bottomrule\end{longtable}
\section*{2026-06-15 (4 commits)}
\begin{longtable}{@{}p{0.105\textwidth}p{0.665\textwidth}r r r@{}}
\toprule Hash & Subject & + & - & Files \\ \midrule\endhead
\texttt{f85b1817} & Add OLMO\_FORCE\_RERUN to re-run OLMo smoke grid over prior results & 25 & 1 & 2 \\
\texttt{38698653} & Add OLMo full-run grid; rename runbook to olmo\_models; group on full runs & 226 & 57 & 13 \\
\texttt{dd6575b6} & Add updated config to work on aiq-gpu & 82 & 12 & 1 \\
\texttt{bef5ccc0} & Restrict to gpus 2,3 & 5 & 0 & 1 \\
\bottomrule\end{longtable}
\section*{2026-06-16 (7 commits)}
\begin{longtable}{@{}p{0.105\textwidth}p{0.665\textwidth}r r r@{}}
\toprule Hash & Subject & + & - & Files \\ \midrule\endhead
\texttt{5c0490c9} & Reserve chat-template headroom for OLMo openai-compatible runs & 98 & 0 & 2 \\
\texttt{8cb5eb75} & Stage natural\_qa data via gcloud auth + helm-run shim (script-only) & 363 & 0 & 6 \\
\texttt{01a68880} & Drop OLMo math runs (competition\_math disabled on HF Hub) & 14 & 7 & 1 \\
\texttt{ddc40407} & Drop OLMo natural\_qa runs (GCS bucket denies authenticated reads) & 40 & 364 & 7 \\
\texttt{4354ef68} & Default OLMo smoke grid to force-rerun; keep full grid opt-in & 10 & 5 & 2 \\
\texttt{5d02e12b} & Add containerized HELM execution path for the kwdagger pipeline & 1484 & 7 & 17 \\
\texttt{7b028102} & Restructure phi-2 e2e tests into the olmo\_models runbook shape & 746 & 121 & 17 \\
\bottomrule\end{longtable}
\section*{2026-06-17 (9 commits)}
\begin{longtable}{@{}p{0.105\textwidth}p{0.665\textwidth}r r r@{}}
\toprule Hash & Subject & + & - & Files \\ \midrule\endhead
\texttt{4d5c411e} & Add containerized phi-2 e2e example (vLLM served on host, HELM via --network host) & 327 & 11 & 10 \\
\texttt{63f8c8cd} & Enable public comparisons for the olmo full virtual experiment & 5 & 5 & 1 \\
\texttt{0d4a4183} & docs(planner): add robust order-insensitive matching plan & 200 & 0 & 1 \\
\texttt{b7373676} & feat(run\_entries): add order-insensitive canonical\_logical\_key & 147 & 0 & 2 \\
\texttt{a25aac93} & feat(planner): group runs by order-insensitive canonical key & 252 & 100 & 2 \\
\texttt{d8e6b698} & docs(planner): document canonical-key matching + journal entry & 91 & 0 & 2 \\
\texttt{ceefb460} & Merge branch 'main' into olmo-reproduction & 0 & 0 & 0 \\
\texttt{36b7860a} & Merge branch 'olmo-reproduction' into main & 0 & 0 & 0 \\
\texttt{f46445d9} & Merge branch 'main' into e2e-refactor & 0 & 0 & 0 \\
\bottomrule\end{longtable}
\section*{2026-06-18 (9 commits)}
\begin{longtable}{@{}p{0.105\textwidth}p{0.665\textwidth}r r r@{}}
\toprule Hash & Subject & + & - & Files \\ \midrule\endhead
\texttt{f23bf696} & Run the phi-2 container e2e example by default (opt-out) & 73 & 28 & 4 \\
\texttt{0e8617a7} & Fix dangling venv python in the HELM-runner image (multi-stage uv) & 59 & 1 & 2 \\
\texttt{fa37fd6c} & Run HF e2e scenario first on a freed GPU; spin vLLM down at grid start & 59 & 7 & 5 \\
\texttt{1cb228a5} & Enable public comparison in the non-smoke e2e experiments & 67 & 26 & 3 \\
\texttt{de439c56} & Split the e2e-phi2 report into one virtual experiment per scenario & 282 & 170 & 11 \\
\texttt{89695fe6} & Force matplotlib Agg backend in report rendering (fix Tk core dump) & 68 & 0 & 6 \\
\texttt{55fb5ecf} & Merge branch 'e2e-refactor' into main & 0 & 0 & 0 \\
\texttt{cc57d5d7} & Run OLMo reproduction through the containerized HELM pipeline & 430 & 4 & 8 \\
\texttt{c8b3d525} & Align OLMo BBQ run entries with public recipe (output\_format\_instructions) & 391 & 7 & 4 \\
\bottomrule\end{longtable}
\section*{2026-06-22 (5 commits)}
\begin{longtable}{@{}p{0.105\textwidth}p{0.665\textwidth}r r r@{}}
\toprule Hash & Subject & + & - & Files \\ \midrule\endhead
\texttt{10c3db97} & Migrate infer\_stack integration to the catalog/leasing API & 387 & 286 & 5 \\
\texttt{f779d596} & Reschema the e2e + OLMo infer-stack config dirs to catalog.yaml & 169 & 548 & 8 \\
\texttt{f4236d61} & Migrate the e2e + OLMo runner scripts to the leasing CLI verbs & 167 & 108 & 8 \\
\texttt{e8e76818} & Update the e2e + OLMo runbook docs for the leasing CLI & 73 & 52 & 3 \\
\texttt{ab8af6c8} & journal: infer-stack CLI/API migration session & 73 & 0 & 1 \\
\bottomrule\end{longtable}
\section*{2026-06-23 (8 commits)}
\begin{longtable}{@{}p{0.105\textwidth}p{0.665\textwidth}r r r@{}}
\toprule Hash & Subject & + & - & Files \\ \midrule\endhead
\texttt{4afe6a66} & Merge branch 'olmo-docker-pipeline' into main & 0 & 0 & 0 \\
\texttt{22395e21} & Merge branch 'main' into infer-stack-cli-api-migration & 0 & 0 & 0 \\
\texttt{d259d248} & docs: plan for infer-stack <-> kwdagger high-throughput integration & 499 & 0 & 1 \\
\texttt{dac25978} & journal: leasing-pipeline-lifecycle planning + implementation session & 132 & 0 & 1 \\
\texttt{ee34857c} & Messy state & 952 & 36 & 19 \\
\texttt{da930f4c} & Remove old-schema config.yaml resurrected by a bad merge resolution & 39 & 162 & 2 \\
\texttt{de9919db} & docs: drop I4/I5, add eval\_audit-integration hand-off & 212 & 28 & 2 \\
\texttt{24991d17} & Wire per-run infer-stack GPU leasing into the HELM fan-out & 960 & 18 & 11 \\
\bottomrule\end{longtable}
\section*{2026-06-24 (11 commits)}
\begin{longtable}{@{}p{0.105\textwidth}p{0.665\textwidth}r r r@{}}
\toprule Hash & Subject & + & - & Files \\ \midrule\endhead
\texttt{4158bea0} & Require containerized HELM execution; leasing as an orthogonal bracket & 358 & 259 & 7 \\
\texttt{c5ff9a35} & Containerize the e2e presets; drop the redundant -container scenario & 72 & 129 & 5 \\
\texttt{03dc08e5} & Lease + containerize by default across the OLMo and e2e grids; drop the toggles & 312 & 292 & 9 \\
\texttt{25811227} & docs: containerization mandatory + leasing default; session journal & 332 & 105 & 4 \\
\texttt{5317e14f} & Update gitmodules and bump commit hashes & 6 & 6 & 5 \\
\texttt{ac444241} & Merge branch 'main' into infer-stack-cli-api-migration & 0 & 0 & 0 \\
\texttt{31dda44c} & infer\_stack: make protocol\_mode a required preset fact (no silent chat default) & 134 & 8 & 3 \\
\texttt{2fd3fd6c} & e2e/olmo grids: pass --yes to the bootstrap 'release --evict' & 4 & 4 & 4 \\
\texttt{94980dda} & journal: leasing CLI bringup notes (\_load\_catalog AttributeError fix) & 47 & 0 & 1 \\
\texttt{d9f075db} & infer\_stack: bump submodule to the \_load\_catalog catalog-field fix & 1 & 1 & 1 \\
\texttt{b2aff8fe} & infer\_stack: bump submodule to the dev/lease-polish merge (995f45e) & 1 & 1 & 1 \\
\bottomrule\end{longtable}
\section*{2026-06-25 (30 commits)}
\begin{longtable}{@{}p{0.105\textwidth}p{0.665\textwidth}r r r@{}}
\toprule Hash & Subject & + & - & Files \\ \midrule\endhead
\texttt{fde72004} & infer\_stack: bump submodule to the dev/lease-polish merge (0cde11d) & 1 & 1 & 1 \\
\texttt{be2a72a6} & catalogs: declare protocol: completions for base-model endpoints & 80 & 7 & 3 \\
\texttt{23a20459} & e2e: always force-rerun both grids; remove the E2E\_FORCE\_RERUN opt-out & 21 & 19 & 2 \\
\texttt{248c3b4b} & e2e: recap every scenario's report path at the end of 40\_build\_summary & 26 & 0 & 1 \\
\texttt{43d3c2b0} & e2e: drop the broken "start here" README link from the summary recap & 1 & 2 & 1 \\
\texttt{212289e5} & docs(planning): bring the run-from-run-spec replay plan onto the impl branch & 404 & 0 & 1 \\
\texttt{1cd44778} & manifests: add from\_run\_spec + --precomputed-root to make-manifest & 41 & 0 & 2 \\
\texttt{622881b2} & pipeline+bridge: route from\_run\_spec runs to the replay docker node & 64 & 0 & 2 \\
\texttt{f54d29c4} & test+journal: from\_run\_spec pipeline wiring tests; session journal & 284 & 0 & 2 \\
\texttt{7ee5ecee} & submodules: bump aiq-magnet gitlink to the from-spec replay CLI (4b10e1b) & 1 & 1 & 1 \\
\texttt{d5180845} & submodules: bump infer\_stack gitlink + track its working branch & 2 & 2 & 2 \\
\texttt{bd097610} & gitmodules: track aiq-magnet impl/run-from-run-spec branch & 1 & 1 & 1 \\
\texttt{5e1edb4b} & docs(planning): plan to migrate the phi-2 e2e to run\_spec.json replay & 237 & 0 & 1 \\
\texttt{353e74b5} & docs(planning): correct the phi-2 e2e from-spec plan after a code review & 59 & 27 & 1 \\
\texttt{f18e5036} & e2e(from-spec): HF override + sibling from-spec manifests (Change 1) & 113 & 0 & 3 \\
\texttt{1d90f453} & e2e(from-spec): teach export-benchmark-bundle the from-spec path (Change 2) & 78 & 2 & 2 \\
\texttt{d81ca3ca} & e2e(from-spec): gate the grid on E2E\_FROM\_SPEC, carve out incomparable (Change 3+4) & 56 & 10 & 3 \\
\texttt{2368fe9a} & e2e(from-spec): tests for the phi-2 from-spec exporter + discovery (Change 6) & 223 & 0 & 1 \\
\texttt{ea0141ae} & journal: session entry for the phi-2 e2e from-spec review + implementation & 59 & 0 & 1 \\
\texttt{daf6dfbb} & e2e(from-spec): make faithful replay the default, drop the opt-out & 81 & 136 & 8 \\
\texttt{33a0606f} & docs(e2e): reflect from-spec-as-default (no toggle) in plan + README & 84 & 47 & 2 \\
\texttt{05d2932e} & journal: from-spec becomes the e2e default (toggle removed) & 36 & 0 & 1 \\
\texttt{e2543fb5} & fix(e2e from-spec): override dropped phi-2's context window -> window-service crash & 15 & 0 & 2 \\
\texttt{05a704d3} & docs(planning): plan to record the local deployment on from-spec replays & 262 & 0 & 1 \\
\texttt{1a119919} & docs(planning): lock in explicit --model-deployment (reject auto-derive) & 31 & 19 & 1 \\
\texttt{f65997d3} & from-spec(deployment-rewrite): record the LOCAL deployment, unmask same\_deployment & 449 & 71 & 12 \\
\texttt{b783a95c} & docs(from-spec): mark deployment-rewrite implemented; amend superseded plans & 165 & 17 & 7 \\
\texttt{a4b5bac5} & docker(helm-runner): split helm/magnet install so magnet edits skip the heavy layer & 27 & 6 & 1 \\
\texttt{bad2755f} & gitignore: ignore docker/build.sh staging context (.build-staging/) & 3 & 0 & 1 \\
\texttt{abc00b3c} & submodules: bump aiq-magnet gitlink to the --model-deployment rewrite (50489f0) & 1 & 1 & 1 \\
\bottomrule\end{longtable}
\section*{2026-06-26 (1 commits)}
\begin{longtable}{@{}p{0.105\textwidth}p{0.665\textwidth}r r r@{}}
\toprule Hash & Subject & + & - & Files \\ \midrule\endhead
\texttt{1eb52704} & docker(helm-runner): raise build-step nofile ulimit to survive uv bytecode compile & 14 & 0 & 1 \\
\bottomrule\end{longtable}
\section*{2026-06-29 (19 commits)}
\begin{longtable}{@{}p{0.105\textwidth}p{0.665\textwidth}r r r@{}}
\toprule Hash & Subject & + & - & Files \\ \midrule\endhead
\texttt{8d96a471} & e2e(infer-stack): resolve data\_dir as env > settings.yaml pin > /data default & 54 & 18 & 2 \\
\texttt{0f09fe87} & submodules: bump every\_eval\_ever gitlink to tmp-branch (b1c5a0f0e) & 1 & 1 & 1 \\
\texttt{15020825} & infer\_stack: bump submodule to the dev/lease-polish catch-up merge (045536c) & 1 & 1 & 1 \\
\texttt{b5c4cfed} & fix(eee-convert): replay run prod\_env so local deployment aliases resolve & 45 & 0 & 1 \\
\texttt{03b98e9f} & feat(e2e): aiq-gpu rsync fetch utilities (17 + 42 + shared lib) & 292 & 1 & 4 \\
\texttt{eb904ea7} & docs(journal): session entries for aiq-gpu rsync utilities + eee-convert fix & 141 & 0 & 1 \\
\texttt{ecd1db15} & docs(e2e): add NOTES.md for the venv pin + from-spec conversion gotchas & 54 & 0 & 2 \\
\texttt{1b75bd0c} & docs(planning): olmo from-run-spec migration plan & 252 & 0 & 1 \\
\texttt{5c31a05f} & feat(olmo-from-spec): Change 4 discovery dry-check + measured baseline & 347 & 30 & 3 \\
\texttt{b2ebad72} & feat(olmo-from-spec): Change 1 -- from-spec preset fields + run-entry reconciliation & 180 & 79 & 6 \\
\texttt{ab5bdc2e} & docs(planning): handoff prompt to continue the olmo from-spec migration & 143 & 0 & 1 \\
\texttt{99bdc0ee} & feat(olmo-from-spec): Change 3 -- wire --from-spec into the smoke/full grids & 20 & 2 & 2 \\
\texttt{037ba682} & test(olmo-from-spec): Change 6 -- discovery dry-check + comparability proof & 174 & 0 & 1 \\
\texttt{1ad68a8b} & fix(olmo infer-stack): Change 7 -- resolve data\_dir as env > settings.yaml pin > /data default & 61 & 7 & 2 \\
\texttt{90d95811} & docs(olmo-from-spec): Change 8 -- README from-spec default + annotate NOTES as resolved & 145 & 43 & 9 \\
\texttt{91413bac} & docs(olmo-from-spec): mark Changes 2/3/6/7/8 done; only Change 5 (GPU) remains & 148 & 49 & 3 \\
\texttt{5accb6cd} & fix(olmo container): install eval\_audit in runner image so in-container helm-run loads the olmo-7b tokenizer override & 71 & 1 & 2 \\
\texttt{536b0607} & docs(journal): in-container olmo-7b tokenizer override (eval\_audit missing from runner image) & 23 & 0 & 1 \\
\texttt{2e1c40e3} & docs(planning): multi-model from-spec via local-deployment-scoped discovery strip & 233 & 0 & 1 \\
\bottomrule\end{longtable}
\section*{2026-06-30 (29 commits)}
\begin{longtable}{@{}p{0.105\textwidth}p{0.665\textwidth}r r r@{}}
\toprule Hash & Subject & + & - & Files \\ \midrule\endhead
\texttt{ae14bc77} & fix(container): deliver HF token via mounted cache, not env-forwarding & 81 & 0 & 2 \\
\texttt{ebcc5b7d} & docs(container): correct HF-token delivery mechanism & 49 & 14 & 6 \\
\texttt{c763a6ea} & docs(journal): HF token on-disk delivery for containerized runs & 24 & 0 & 1 \\
\texttt{9ce64278} & fix(olmo container): install crfm-helm[all] + pin huggingface\_hub==0.36.2 & 25 & 10 & 3 \\
\texttt{0fdc7457} & test(olmo smoke): preflight probes container python env (langdetect + hf\_hub pin) & 57 & 5 & 2 \\
\texttt{528d341a} & docs(journal): container env fix (crfm-helm[all] + hf\_hub pin) + smoke probe & 25 & 0 & 1 \\
\texttt{bf2b9767} & test(olmo smoke): add wmt\_14 + ifeval recipe canaries to smoke grids & 23 & 3 & 2 \\
\texttt{676746d0} & docs(journal): wmt\_14 + ifeval smoke canaries (from-spec-constrained selection) & 16 & 0 & 1 \\
\texttt{99e27c51} & test(e2e): preflight probes container python env (langdetect + hf\_hub pin) & 53 & 3 & 2 \\
\texttt{3b32c1eb} & docs(journal): e2e container-env probe (no canary, per request) & 14 & 0 & 1 \\
\texttt{8d422490} & docs(journal): resolve OLMo-canary question (kept; "no canary" was e2e-only) & 1 & 1 & 1 \\
\texttt{8dd861b5} & docs(planning): address reproductions by (public\_root, relative\_path) & 369 & 0 & 1 \\
\texttt{b7b115cc} & docs(planning): rel-path plan -- kwdagger cross-products, broadcast map mandatory & 15 & 17 & 1 \\
\texttt{e94d2c86} & docs(planning): rel-path plan -- host-side resolution + kwdagger submatrix & 284 & 300 & 1 \\
\texttt{1d1097d0} & docs(planning): rel-path plan -- materialize substituted copies host-side & 135 & 92 & 1 \\
\texttt{77f13574} & test(kwdagger): pin the submatrix-zip contract the rel-path plan relies on & 123 & 0 & 2 \\
\texttt{515bddbd} & feat(from-spec): host-side run\_spec resolver + materializer (rel-path plan §4.1) & 452 & 0 & 2 \\
\texttt{6305d294} & feat(from-spec): wire exact-path replay through the bridge + docker node (§4.2-4.6) & 443 & 54 & 4 \\
\texttt{22026464} & test(from-spec): pin run\_spec\_json on the from-spec node algo\_params & 6 & 3 & 1 \\
\texttt{337ac8d4} & feat(from-spec): make-manifest --run-spec-sources-fpath for exact-path replay (§4.5) & 157 & 4 & 2 \\
\texttt{c8fe9b76} & docs(from-spec): mark rel-path plan partially implemented + journal entry & 24 & 1 & 2 \\
\texttt{bbbeaeb7} & fix(from-spec): gate materialized per-run lease\_endpoint on --lease & 25 & 3 & 2 \\
\texttt{5f568bf5} & feat(from-spec): exporter auto-freeze of rel-paths (rel-path plan §4.5) & 287 & 4 & 3 \\
\texttt{a9a1a16b} & feat(from-spec): dry-check existence mode for frozen manifests (rel-path plan §4.7) & 202 & 0 & 2 \\
\texttt{249735d3} & docs(from-spec): mark rel-path plan host-side-complete + journal entry & 24 & 10 & 2 \\
\texttt{cea323f2} & fix(from-spec): content-address the materialized run\_spec filename (kwdagger identity) & 69 & 11 & 3 \\
\texttt{d9b60ec8} & docs(from-spec): document kwdagger caching model + content-addressed identity & 35 & 5 & 2 \\
\texttt{115a9997} & debug(olmo): MWE deployment matrix for OLMo-7B fp16-collapse diagnosis & 1170 & 0 & 9 \\
\texttt{74ba33de} & fix(olmo-7b): serve OLMo-1.7 tokenizer to stop spurious EOS append & 16 & 0 & 1 \\
\bottomrule\end{longtable}
\section*{2026-07-01 (15 commits)}
\begin{longtable}{@{}p{0.105\textwidth}p{0.665\textwidth}r r r@{}}
\toprule Hash & Subject & + & - & Files \\ \midrule\endhead
\texttt{ab9ccce8} & feat(deployment-match): general HELM deployment-match search tool (Phase 1) & 1687 & 0 & 11 \\
\texttt{86ad9598} & feat(deployment-match): Phase 2 -- two-tier serve+probe driver (`run`) & 169 & 0 & 2 \\
\texttt{aee205b4} & feat(deployment-match): Phase 3 -- confirm winner via compare-pair (`confirm`) & 174 & 0 & 2 \\
\texttt{792ce211} & test(deployment-match): Phase 4 -- pytest coverage + testable tokenizer predicate & 196 & 14 & 2 \\
\texttt{c978dea6} & docs(deployment-match): mark plan IMPLEMENTED + document run/confirm & 60 & 16 & 3 \\
\texttt{88667c69} & tooling(catalog): catalog core-report packets by decoding temperature & 234 & 0 & 1 \\
\texttt{d703ad4c} & feat(deployment-match): `auto` -- one-shot dry-run->run->score->confirm & 73 & 1 & 2 \\
\texttt{c8d8417d} & fix(deployment-match): let infer-stack place GPUs (acquire --queue), don't pin & 38 & 9 & 4 \\
\texttt{aaab9a35} & tooling(catalog): detect chain-of-thought alongside temperature & 143 & 54 & 1 \\
\texttt{f04edd85} & fix(from-spec): thread freeze\_rel\_paths through export->materialize & 54 & 3 & 3 \\
\texttt{b998f53f} & feat(olmo): combined multi-model preset for parent-root fan-out & 209 & 7 & 2 \\
\texttt{bb2d959d} & feat(olmo): combined multi-model fan-out runbook & 613 & 0 & 13 \\
\texttt{c6d7a294} & chore(olmo): remove superseded fp16-collapse debug MWE & 0 & 1154 & 8 \\
\texttt{d6534170} & fix(from-spec): strip model\_deployment token from exact-path locator run\_entry & 67 & 1 & 2 \\
\texttt{a2879335} & chore(olmo): remove default GPU restriction (serve across all detected GPUs) & 12 & 7 & 3 \\
\bottomrule\end{longtable}
\section*{2026-07-02 (6 commits)}
\begin{longtable}{@{}p{0.105\textwidth}p{0.665\textwidth}r r r@{}}
\toprule Hash & Subject & + & - & Files \\ \midrule\endhead
\texttt{88262878} & fix(olmo): enable public comparison in combined vexp (was local-only) & 12 & 12 & 1 \\
\texttt{065b683a} & docs(olmo): remove stale "local-only by default" claims & 17 & 18 & 3 \\
\texttt{a14f2faa} & docs(planning): plan for from-spec adapter\_spec decode-param overrides & 158 & 0 & 1 \\
\texttt{71249b9d} & Merge branch 'jons/fable-commits' into impl/run-from-run-spec & 0 & 0 & 0 \\
\texttt{3014cea3} & docs(journal): session entry for fable-commits review + merge & 12 & 0 & 1 \\
\texttt{52513fc9} & feat(olmo): fold base olmo-7b into the combined fan-out vexp (all six models) & 127 & 38 & 8 \\
\bottomrule\end{longtable}
\section*{2026-07-03 (34 commits)}
\begin{longtable}{@{}p{0.105\textwidth}p{0.665\textwidth}r r r@{}}
\toprule Hash & Subject & + & - & Files \\ \midrule\endhead
\texttt{a90a084f} & docs(planning): full codebase audit -- findings catalog + phased fix plan & 587 & 0 & 2 \\
\texttt{63cac3fc} & test(kwdagger): guard collection with importorskip("kwdagger") (P0-9) & 13 & 0 & 5 \\
\texttt{fa8b745a} & fix(cli): repair analyze\_backlog import + canonical report layout (P0-8) & 9 & 8 & 2 \\
\texttt{a5a2a2f5} & build(deps): pin transformers<5 / hf-hub, declare direct deps (P1-24) & 13 & 2 & 2 \\
\texttt{3e98614f} & fix(summary): tol010 sankeys bucket on the 0.01 curve point (P0-1) & 68 & 8 & 3 \\
\texttt{6d911222} & fix(index): resolve unstamped local index first; single-source resolver (P0-2, R-6) & 108 & 22 & 4 \\
\texttt{dbe3b68f} & fix(planning,virtual): match/join on canonical keys, not raw strings (P0-3, P0-4) & 171 & 7 & 5 \\
\texttt{05c9dd62} & fix(helm): correct perturbed/unperturbed agreement split (P0-6) & 97 & 4 & 3 \\
\texttt{b42064c0} & fix(planning): unknown-on-partial facts; extract official max\_eval\_instances (P1-1, P1-2) & 134 & 6 & 4 \\
\texttt{09fc3294} & fix(normalized): filter non-finite scores at row construction (P1-15) & 181 & 33 & 5 \\
\texttt{d05126d0} & fix(reports): cross-machine curve uses rel\_tol=0 highlights (P1-13) & 93 & 1 & 3 \\
\texttt{0201ec06} & fix(eee): derive experiment name from local/<experiment>/ layout (D-1) & 59 & 28 & 7 \\
\texttt{86d41ae5} & fix(bundle): content-address model\_deployments.yaml filename (P0-5) & 52 & 3 & 2 \\
\texttt{5a7d270d} & fix(pipeline): precomputed\_root is algo identity on from-spec node (P1-21) & 30 & 5 & 2 \\
\texttt{2896be59} & fix(manifests): missing-entries NameError; reject duplicate labels (P1-22, P1-23) & 25 & 2 & 3 \\
\texttt{54bf35f4} & fix(docker): image-id identity for unpinned images; --name + teardown (P2) & 144 & 13 & 6 \\
\texttt{c5ac0203} & fix(summary): failure-taxonomy honesty -- colors, categories, ordering (P0-7, P1-8, P1-11) & 152 & 11 & 4 \\
\texttt{7642a1b6} & fix(summary): not\_analyzed status level; surface unreadable bundles (P1-9, P1-10) & 87 & 19 & 4 \\
\texttt{c5d7cdb5} & fix(filter): one sankey row per excluded run, primary reason (P1-12) & 38 & 2 & 2 \\
\texttt{cdac08ec} & fix(eee): status.json gate on cached/local resolvers; unconditional stale-scope delete (P1-7, P1-6) & 94 & 13 & 4 \\
\texttt{6dac5e99} & refactor: delete the no-op --allow-single-repeat flag (D-3) & 16 & 32 & 6 \\
\texttt{ffac9b4e} & fix(cli): analyze-many exits non-zero on failure; warns on incomplete summary (P1-19) & 74 & 0 & 2 \\
\texttt{41e5b962} & fix(core-metrics): --no-plots preserves previously rendered figures (P1-16) & 26 & 7 & 2 \\
\texttt{2537b388} & fix(summary): thread invocation flags into reproduce.sh (P1-20) & 72 & 5 & 3 \\
\texttt{b3130f42} & fix(stale-state): prune stale packet + breakdown dirs on re-run (P1-4, P1-5) & 37 & 1 & 2 \\
\texttt{0ae927fb} & fix(determinism): stable benchmark discovery + HelmRunDiff tie-breaks (P1-17, P1-18) & 57 & 8 & 4 \\
\texttt{a9a20979} & fix(determinism): sort remaining nondeterministic outputs (P2, IM-12) & 28 & 5 & 4 \\
\texttt{bdf94a70} & fix(cli): --dry-run vs --run 1 errors; correct stale argv test (P2) & 12 & 1 & 2 \\
\texttt{688831a3} & fix(helm-diff): f-string prefix + None-safe group sort (P2) & 7 & 3 & 1 \\
\texttt{14ce8042} & fix(cli,loaders): drop debug blocks; count skipped samples lines (P2, IM-3) & 13 & 13 & 2 \\
\texttt{99396f3f} & fix(secrets,coverage): close HF-token perms window; return sankey paths (P2, IM-9) & 14 & 2 & 2 \\
\texttt{235b555e} & fix(loaders): tolerate non-numeric timestamp; preserve EEE provenance (P2) & 14 & 5 & 1 \\
\texttt{10c8f32d} & docs(journal): session entry for Phases A--F audit-fix implementation & 22 & 0 & 1 \\
\texttt{9303ec5e} & docs(preset,runbook): representative-subset wording (D-2); hashed bundle filename in tree listing & 7 & 5 & 2 \\
\bottomrule\end{longtable}
\section*{2026-07-06 (50 commits)}
\begin{longtable}{@{}p{0.105\textwidth}p{0.665\textwidth}r r r@{}}
\toprule Hash & Subject & + & - & Files \\ \midrule\endhead
\texttt{0608b29c} & test(planner): update stale slow-suite assertions to canonical-key behavior & 15 & 9 & 1 \\
\texttt{15487f8e} & fix(secrets): chmod 0600 the model\_deployments bundle YAML holding the live key & 24 & 0 & 2 \\
\texttt{6da780de} & fix(reports): route Stage-6 writers through atomic fs\_publish helpers (IM-1) & 18 & 14 & 3 \\
\texttt{607189c8} & fix(reports,cli): surface mpl fallback error; fix eee caveats indentation & 48 & 7 & 3 \\
\texttt{b127aa48} & fix(cli,adapter): fail fast on backlog/preset/base\_url misconfig (P2) & 179 & 12 & 7 \\
\texttt{20a16f5f} & fix(workflows): dedupe kwdg re-attempts, bucket load errors, preserve CSV dtypes (P2) & 151 & 14 & 3 \\
\texttt{ed576083} & fix(reports,cli,loaders): honesty fixes IM-6/IM-8/IM-5 & 208 & 3 & 7 \\
\texttt{0e324ed1} & perf(reports,loaders): behavior-preserving IM-2/IM-7/IM-4 & 33 & 28 & 4 \\
\texttt{e9cf6c44} & build(deps): bump python floor to 3.12; annotate aiq-magnet dep pins; relock & 537 & 1318 & 2 \\
\texttt{dc8f680e} & test(eee): centralize EEE-demo fixture path + skip guard in conftest (IM-11) & 70 & 30 & 7 \\
\texttt{e09c9176} & docs(eee-vs-helm): document non-comparable run-level agreement definitions (IM-13) & 34 & 0 & 1 \\
\texttt{8a9933c7} & docs(journal): session entry for Tier-1/Tier-2 audit-backlog implementation & 30 & 0 & 1 \\
\texttt{ef6903fb} & fix(secrets): warn on stderr when bundle key perms cannot be tightened & 9 & 0 & 1 \\
\texttt{bab7df2c} & refactor(reports): retire reports/quantiles.py (D-4) & 2 & 193 & 3 \\
\texttt{e0d51e6b} & feat(manifests): --max-eval-instances official sentinel for verbatim replay (D-5) & 136 & 4 & 5 \\
\texttt{5b1c69f2} & refactor(indexing): single classify\_model\_eligibility policy (R-7) & 211 & 94 & 3 \\
\texttt{833a84bb} & refactor(summary): single \_classify\_filter\_gates; delete dead sankey builders (R-4a/R-4b) & 85 & 392 & 4 \\
\texttt{99ff4494} & feat(reproduce/olmo): deployment-match runbook for olmoe-1b-7b/ifeval & 130 & 0 & 1 \\
\texttt{c8cffa42} & fix(summary): align the two filter-funnel families on metadata + judge gates (R-4c) & 85 & 1 & 5 \\
\texttt{fe7fe271} & refactor(summary): rename ambiguous agreement-tolerance key family (R-4d) & 69 & 58 & 8 \\
\texttt{b6bcd245} & refactor(normalized): dedupe EEE conversion core + one aggregate predicate (R-5) & 185 & 165 & 5 \\
\texttt{ba730045} & refactor(normalized): one metric\_handle helper; drop dead members (R-9) & 21 & 58 & 5 \\
\texttt{899e0fd3} & refactor(core-metrics): dead-surface sweep + assert-membership tol guard (R-10) & 41 & 254 & 7 \\
\texttt{143d3fac} & refactor: misc dead-code sweep (R-11) & 70 & 150 & 10 \\
\texttt{e5c1bdb4} & refactor: one run-entry model\_deployment token utility (R-8) & 165 & 34 & 5 \\
\texttt{7e1f5f12} & refactor(infer\_stack): extract PRESET\_CONFIGS to presets.py (R-3, pure relocation) & 1162 & 1140 & 2 \\
\texttt{d895a446} & refactor(infer\_stack): promote discovery core to public discovery.py (R-3, R-3d) & 105 & 78 & 4 \\
\texttt{e7eab858} & refactor(infer\_stack): extract serving\_facts.py (R-3b, pure relocation) & 165 & 128 & 2 \\
\texttt{9a808120} & refactor(infer\_stack): extract freeze.py (R-3d, pure relocation) & 93 & 76 & 2 \\
\texttt{ac5ca26a} & refactor(infer\_stack): extract bundle\_export.py; adapter is now a facade (R-3c) & 667 & 642 & 3 \\
\texttt{79bc030c} & docs(journal): session entry for Tier-3 structural refactors (D-4/D-5, R-3..R-11) & 31 & 0 & 1 \\
\texttt{560cdda0} & docs(planning): stamp audit doc IMPLEMENTED (except R-2) with pass summary & 17 & 0 & 1 \\
\texttt{cfe7d6a6} & docs(planning): simplicity \& bloat audit findings + phased implementation plan & 147 & 0 & 1 \\
\texttt{5b7b62bb} & chore(dev): delete dead POC and oneoff trees; fix docs that referenced them & 22 & 6901 & 25 \\
\texttt{1457a269} & chore(configs): drop checked-in generated manifests and dead examples/ & 0 & 1470 & 25 \\
\texttt{8b4c69d5} & refactor: remove helm/hashers shim, dead agreement profile, if-1 block & 13 & 178 & 8 \\
\texttt{90af80ba} & docs: fix first-run and accuracy issues in README, reproduce index, setup & 79 & 82 & 3 \\
\texttt{77e126b4} & chore(docs): move paper analysis scripts out of docs/, archive session logs & 32 & 17 & 9 \\
\texttt{425e05a2} & docs(planning): archive 12 implemented plans; consolidate olmo/infer-stack clusters & 193 & 132 & 67 \\
\texttt{24390f26} & chore(journals): rotate 544K of append-only journals to archive & 5542 & 5526 & 4 \\
\texttt{c804cdc1} & refactor(utils): finish R-6 helper consolidation into utils/coercion & 159 & 137 & 14 \\
\texttt{32622e0b} & refactor(reproduce): merge three divergent \_lib.sh copies into one shared lib & 437 & 374 & 4 \\
\texttt{51797aa9} & refactor(workflows): split \_render\_scope\_summary into per-stage render helpers & 458 & 197 & 1 \\
\texttt{a07736c2} & refactor(reports): split core\_metrics.main into per-section helpers & 182 & 57 & 1 \\
\texttt{3bd43bf5} & refactor(infer\_stack): move PRESET\_CONFIGS catalog to a YAML resource & 965 & 1009 & 3 \\
\texttt{a849ea32} & refactor(reports)!: migrate pair\_report/pair\_samples onto NormalizedDiff (R-2 stage A) & 284 & 117 & 4 \\
\texttt{a0cadf26} & refactor(helm)!: delete the legacy agreement half of diff.py (R-2 stage B) & 59 & 986 & 6 \\
\texttt{a34241ca} & docs: stamp simplicity audit IMPLEMENTED; journal entry for the session & 69 & 2 & 2 \\
\texttt{066cc520} & refactor(reproduce/olmo): self-contained combined runbook; remove stale single-run runbook & 323 & 1655 & 24 \\
\texttt{7ba01845} & chore: remove configs/docs orphaned by the single-model OLMo runbook removal & 28 & 481 & 13 \\
\bottomrule\end{longtable}
\section*{2026-07-07 (22 commits)}
\begin{longtable}{@{}p{0.105\textwidth}p{0.665\textwidth}r r r@{}}
\toprule Hash & Subject & + & - & Files \\ \midrule\endhead
\texttt{a1f893fa} & fix(deployment\_match): clamp max\_model\_len to the model's derived ceiling & 85 & 9 & 4 \\
\texttt{f4241a87} & docs(deployment\_match): vLLM-vs-HuggingFace deployment-match knob inventory & 166 & 0 & 1 \\
\texttt{2a9fb007} & feat(deployment\_match): hf-match grid profile to match a HuggingFaceClient run & 126 & 4 & 5 \\
\texttt{008987b1} & docs(deployment\_match): record full sampling replay as a deliberate non-goal & 9 & 3 & 1 \\
\texttt{e4e4b7c7} & feat(deployment\_match): sweep/pin attention\_backend via the infer-stack option & 109 & 35 & 6 \\
\texttt{d5cb46e8} & docs: instruct agents to proactively commit logical units of work & 11 & 0 & 1 \\
\texttt{acf6ca44} & docs(deployment\_match): correct the confirm-catalog fidelity gap description & 15 & 6 & 1 \\
\texttt{0a755c73} & fix(deployment\_match): confirm catalog reproduces the winning serving runtime & 69 & 18 & 5 \\
\texttt{070f986c} & feat(deployment\_match): sweep attention\_backend in hf-match instead of pinning & 55 & 36 & 4 \\
\texttt{7a0908ea} & feat(reproduce/olmo): default the OLMoE deployment-match to hf-match & 18 & 10 & 1 \\
\texttt{3fd80ddb} & fix(deployment\_match): raise max\_num\_batched\_tokens >= max\_model\_len for hf-match & 46 & 2 & 2 \\
\texttt{4b8eb37d} & chore(submodules): bump cmd\_queue to dev/0.3.2, kwdagger to dev/0.2.6 & 4 & 4 & 3 \\
\texttt{c436308b} & chore(submodules): bump infer\_stack to latest infer-stack-cli-api-migration & 1 & 1 & 1 \\
\texttt{2640e108} & feat(deployment\_match): endpoint/cell counters + ETA in the serve loop & 41 & 9 & 1 \\
\texttt{f855dea9} & feat(deployment\_match): --dtypes/--attention-backends knobs; OLMoE skips fp32 & 83 & 0 & 4 \\
\texttt{608bb45f} & feat(deployment\_match): preflight feasibility filter (generalizes fp32-MoE prune) & 151 & 27 & 6 \\
\texttt{a8dd792b} & fix(deployment\_match): hf-match pins protocol=completions (HuggingFaceClient repro) & 36 & 1 & 3 \\
\texttt{fbc2040d} & fix(deployment\_match): revert wrong completions pin; HuggingFaceClient IS a chat client & 55 & 31 & 3 \\
\texttt{1fb91576} & feat(deployment\_match): --log-requests / --extra-serve-args to expose the vLLM prompt & 73 & 0 & 5 \\
\texttt{d32abd3b} & feat(deployment\_match): compare\_prompt.py -- diff HELM get\_prompt vs vLLM /tokenize & 266 & 5 & 4 \\
\texttt{b3ae0168} & feat(deployment\_match): sweep add\_generation\_prompt (OLMoE chat-template version drift) & 96 & 12 & 5 \\
\texttt{19c0ec1c} & bump infer\_stack pointer & 1 & 1 & 1 \\
\bottomrule\end{longtable}
\section*{2026-07-08 (26 commits)}
\begin{longtable}{@{}p{0.105\textwidth}p{0.665\textwidth}r r r@{}}
\toprule Hash & Subject & + & - & Files \\ \midrule\endhead
\texttt{1690b7ac} & feat(deployment\_match): per-model ifeval runbook scripts for the 3 OLMo-2 instruct models & 495 & 0 & 3 \\
\texttt{c2c3978d} & Merge branch 'impl/simplicity-audit' into integrate/simplicity-audit & 0 & 0 & 0 \\
\texttt{37a411a5} & docs: fix dangling markdown links after the simplicity-audit merge & 81 & 81 & 9 \\
\texttt{9059dad2} & Merge branch 'integrate/simplicity-audit' into impl/run-from-run-spec & 0 & 0 & 0 \\
\texttt{2698389c} & feat(deployment\_match): --fp32-tensor-parallel-size to serve fp32 MoE on N GPUs & 151 & 11 & 6 \\
\texttt{426710fd} & docs: inventory HELM's unrecorded deployment params + HF-revision pinning effort & 231 & 0 & 2 \\
\texttt{a8c573a3} & docs: use precise terms for the provenance machinery & 47 & 5 & 1 \\
\texttt{37d91aff} & feat(deployment\_match): hf-probe -- reproduce the official at fp32 via HF and score vs oracle & 401 & 12 & 4 \\
\texttt{f52afee7} & docs: confirm OLMoE fp32 exact-match + scope the dtype gap to all OLMo-2 HF runs & 35 & 2 & 2 \\
\texttt{79ab77b6} & feat(deployment\_match): DM\_HF\_FP32 + DM\_FP32\_TP on the OLMo-2 7B/13B/32B runbooks & 127 & 33 & 3 \\
\texttt{2d792012} & feat(hf-inprocess): reproduce HuggingFace-deployed officials on reserved GPUs & 765 & 18 & 6 \\
\texttt{19cf7d40} & bump infer-stack submodule pointer & 1 & 1 & 1 \\
\texttt{9773233e} & Add plan to allow reproductions of old HELM runs in docker containers with package versions pinned to the ones used during the era & 90 & 0 & 1 \\
\texttt{545357e7} & feat(heatmap): aggregate-score-diff heatmap (color = local-public, cells annotated with both scores) & 631 & 0 & 4 \\
\texttt{59587f29} & docs(journal): aggregate-score-diff heatmap session & 61 & 0 & 1 \\
\texttt{c8707ead} & feat(summary): auto-emit aggregate-score-diff heatmaps in build\_reports\_summary & 234 & 62 & 5 \\
\texttt{f87e92ea} & docs(journal): auto-emit aggregate-score-diff in build\_reports\_summary & 48 & 0 & 1 \\
\texttt{ce63eb40} & docs: mark HF-in-process plan implemented (routing pending) + journal entry & 77 & 2 & 2 \\
\texttt{5b352269} & docs(deployment-match): point to the built HF-inprocess resolver + mechanism & 19 & 0 & 1 \\
\texttt{0c8d0eb2} & feat(summary): spell out agreement-bucket thresholds on plot legends & 93 & 12 & 4 \\
\texttt{0846fb99} & docs(journal): agreement-bucket threshold labels on legends & 50 & 0 & 1 \\
\texttt{564b087a} & Bump infer\_stack pointer & 1 & 1 & 1 \\
\texttt{c00394ee} & feat(deployment\_match): overnight batch -- HF fp32 + vLLM sweep for all OLMo instruct models & 124 & 0 & 1 \\
\texttt{3ff73c5c} & feat(deployment\_match): hf-probe sweeps dtypes (--dtype comma list) + DM\_HF\_DTYPES & 91 & 39 & 8 \\
\texttt{70c867da} & feat(heatmap): holistic headline-metric aggregate-score-diff heatmap & 376 & 10 & 7 \\
\texttt{4c2de25b} & docs(journal): holistic headline-metric aggregate-score-diff heatmap & 55 & 0 & 1 \\
\bottomrule\end{longtable}
\section*{2026-07-10 (16 commits)}
\begin{longtable}{@{}p{0.105\textwidth}p{0.665\textwidth}r r r@{}}
\toprule Hash & Subject & + & - & Files \\ \midrule\endhead
\texttt{29fc76c4} & feat(heatmap): headline plot uses modelsxbenchmarks layout + squared error & 116 & 56 & 3 \\
\texttt{236b3289} & docs(journal): headline plot layout + squared-error encoding & 40 & 0 & 1 \\
\texttt{50a7117b} & docs(journal): fp32 reproduces all OLMo instruct officials + HF-probe OLMo-2 fidelity bug & 97 & 0 & 1 \\
\texttt{4516b704} & feat(hf-probe): sweep forward-pass numerics (attn\_impl/device\_map/decode) to reproduce OLMo-2 officials & 225 & 54 & 5 \\
\texttt{807144de} & docs(journal): OLMo-2 HF divergence is a fp32 forward-pass diff + hf-probe numerics-sweep axes & 118 & 0 & 1 \\
\texttt{d6cc37fd} & feat(overnight): sweep every HF + vLLM parameter across all OLMo models & 107 & 7 & 5 \\
\texttt{60921441} & docs(journal): overnight full-sweep follow-up (HF+vLLM axes across all OLMo models) & 23 & 0 & 1 \\
\texttt{25b00db3} & feat(eras): era-pinned HELM registry + resolver (commit 1/6) & 475 & 0 & 5 \\
\texttt{836bbb52} & feat(eras): era image build path in build.sh + era dockerfile (commit 2/6) & 507 & 25 & 4 \\
\texttt{8d54db18} & feat(eras): helm\_era\_shim replay CLI + OpenAI-compatible client (commit 3/6) & 834 & 0 & 5 \\
\texttt{3670050f} & feat(eras): host-side era yaml + materializer guard (commit 4/6) & 252 & 19 & 5 \\
\texttt{38940ff2} & feat(eras): thread era through manifest + pipeline + bridge (commit 5/6) & 379 & 1 & 6 \\
\texttt{a520062b} & docs(eras): runbook + docs + journal (commit 6/6) & 358 & 1 & 9 \\
\texttt{5f17e3e3} & docs(eras): multi-angle code-review findings for the era implementation & 479 & 0 & 1 \\
\texttt{11fa971d} & feat(eras): portable validation-ladder harness + static era-import checker & 738 & 5 & 10 \\
\texttt{39c020b2} & docs(era-tests): plan the pre-v0.5 gates' migration to a turnkey dev runbook & 431 & 0 & 2 \\
\bottomrule\end{longtable}
\section*{2026-07-11 (32 commits)}
\begin{longtable}{@{}p{0.105\textwidth}p{0.665\textwidth}r r r@{}}
\toprule Hash & Subject & + & - & Files \\ \midrule\endhead
\texttt{0c58ee96} & fix(era-shim): v0.2.4 compatibility + credentials + run-dir safety (findings 1,2,3,7,8) & 300 & 49 & 4 \\
\texttt{befd490d} & fix(era-export): api\_key in args, explicit base\_url, multi-endpoint lease (findings 2,5,10) & 133 & 14 & 4 \\
\texttt{f549943e} & fix(era-guard): digest-pinned image era<->image guard (finding 4) & 107 & 16 & 3 \\
\texttt{9bf2228b} & fix(era-image): persist era identity into runtime ENV (finding 6) & 14 & 0 & 1 \\
\texttt{f09051be} & fix(eras): fail loud on track-less era-suite paths (finding 9) & 48 & 2 & 2 \\
\texttt{0486bb79} & refactor(era): process-context timing, session reuse, setdefault (below-the-cap) & 53 & 2 & 4 \\
\texttt{4bba8dfa} & docs(era): record findings resolution + deferred cleanups & 18 & 0 & 1 \\
\texttt{67c0b53f} & feat(era-tests): checked-in serving config + era presets + vexp manifests & 192 & 0 & 5 \\
\texttt{99ff4dfd} & feat(era-tests): dev runbook mirroring dev/e2e-tests (moves the pre-v0.5 gates) & 959 & 429 & 24 \\
\texttt{3046c8ea} & docs(journal): full era-path remediation + dev/era-tests runbook & 71 & 0 & 1 \\
\texttt{4c181be7} & fix(era-tests): review fixes -- v0.2.4 master-key path + rung mirror-root join & 117 & 11 & 8 \\
\texttt{981dfa1f} & feat(era-tests): swap subject pythia-6.9b -> redpajama-3b to fit an 8GB GPU & 204 & 111 & 13 \\
\texttt{c2598b49} & fix(era-build): import run\_benchmarking from helm.benchmark.run, not .runner & 22 & 2 & 2 \\
\texttt{2e8338ff} & fix(era-build): pin pyarrow==11.0.0 in era constraints (datasets 2.5.2 compat) & 10 & 0 & 2 \\
\texttt{4dc3c4a4} & fix(era-build): adopt era frozen requirements.txt as constraints (enriched seed) & 451 & 23 & 3 \\
\texttt{16ee1eeb} & fix(era-build): revert enriched freeze; setup.cfg \textasciitilde{}= ranges already era-pin & 47 & 425 & 3 \\
\texttt{e6a21ed1} & docs(era-build): warn against re-adopting the stale era requirements.txt & 16 & 0 & 2 \\
\texttt{37677132} & fix(era-gate): rung-2 diffs display\_requests.json, not the unshipped scenario\_state.json & 108 & 35 & 3 \\
\texttt{cb14ec80} & fix(era-gate): sanitize colons from rung-2 dry-run mount path & 5 & 1 & 1 \\
\texttt{1c552be2} & fix(era-build): install wget + unzip in the era image for scenario downloads & 9 & 1 & 1 \\
\texttt{38260572} & feat(reproduce): from-spec gpt-oss-20b runbook (bbq, ifeval, mmlu\_pro, gpqa) & 1006 & 0 & 19 \\
\texttt{e40ee63d} & fix(era-build): pin fsspec==2022.8.2 (datasets 2.5.2 LocalFileSystem skew) & 14 & 0 & 2 \\
\texttt{6b23e9d8} & fix(era-gate): forward HF token into rung-2 dry-run (script-dataset auth) & 4 & 0 & 1 \\
\texttt{d6172b1e} & docs(planning): qwen text-family combined fan-out implementation plan & 377 & 0 & 1 \\
\texttt{973b5b7c} & feat(qwen-fanout): qwen-combined preset + 8 from-spec members (generalize combined builder) & 1395 & 92 & 4 \\
\texttt{fcd14973} & feat(reproduce): qwen\_models\_combined runbook + infer-stack catalog + virtual experiment & 1111 & 0 & 15 \\
\texttt{531bc4f8} & docs(planning): mark qwen fan-out plan implemented (T1 confirmed; V2/V3/V5/V6 deferred) & 17 & 1 & 1 \\
\texttt{47200f34} & docs(journal): qwen combined fan-out implementation session & 61 & 0 & 1 \\
\texttt{7620c246} & fix(era-gate): rung-2 skips unfetchable-data probes as environment filters & 32 & 0 & 2 \\
\texttt{629c4541} & docs(planning): plan for producing NEW Qwen 3.6 core results (compute, not reproduce) & 339 & 0 & 1 \\
\texttt{faef9ea4} & docs(planning): rewrite Qwen 3.6 new-results plan around leased fan-out & 212 & 253 & 1 \\
\texttt{17f02583} & fix(provenance): double-quote the label key in image\_label's Go template & 24 & 1 & 2 \\
\bottomrule\end{longtable}
\section*{2026-07-12 (39 commits)}
\begin{longtable}{@{}p{0.105\textwidth}p{0.665\textwidth}r r r@{}}
\toprule Hash & Subject & + & - & Files \\ \midrule\endhead
\texttt{eb1b9faf} & Merge impl/era-pinned-helm-containers into impl/run-from-run-spec & 0 & 0 & 0 \\
\texttt{34929ae4} & docs(journal): era-pinned -> run-from-run-spec merge session & 16 & 0 & 1 \\
\texttt{ea74ac38} & docs(planning): repo simplification/refactor plan (4th pass, post-era) & 479 & 0 & 1 \\
\texttt{bffb84a6} & docs(planning): second-pass review revisions to the simplification plan & 141 & 96 & 1 \\
\texttt{dee2ba95} & docs(journal): fourth simplification-pass audit + plan review session & 61 & 0 & 1 \\
\texttt{b77bdf82} & docs(planning): A4 gate caveat -- committed phase3 baseline is EEE-only & 6 & 1 & 1 \\
\texttt{bf9172f4} & docs(planning): live-code read pass -- deltas R-g..R-m + new item E5 & 100 & 15 & 1 \\
\texttt{911a993f} & refactor(E1): delete five grep-verified dead symbols & 3 & 135 & 6 \\
\texttt{f59060cb} & refactor(E2): de-duplicate the copied plotting helpers & 97 & 78 & 8 \\
\texttt{b7300a19} & refactor(E4a): prune facade re-exports with zero internal+external refs & 20 & 78 & 3 \\
\texttt{cc487686} & refactor(E5a): retire the deprecated EVAL\_AUDIT\_EEE\_STRICT env override & 19 & 19 & 2 \\
\texttt{88c6e389} & refactor(B3): era config/build dedup -- single source of truth & 55 & 36 & 5 \\
\texttt{96b88300} & chore(F2): trim vestigial norecursedirs entries & 5 & 10 & 1 \\
\texttt{4dae513c} & refactor(D1): shared index-resolution + reproduce-script helpers & 111 & 88 & 5 \\
\texttt{7926209e} & refactor(D2): shared EEE render helpers; single-file meta extractor & 150 & 104 & 4 \\
\texttt{1f06e9a5} & refactor(A1): move run\_entries out of helm/ to eval\_audit/run\_entries.py & 46 & 23 & 20 \\
\texttt{9f7c6abd} & refactor(B1): one shared benchmark\_output path parser & 50 & 30 & 3 \\
\texttt{7e8b527d} & fix(E4a): restore \_single\_run\_core\_stat\_index facade re-export & 1 & 60 & 3 \\
\texttt{8078cc1f} & refactor(A2): delete the prod-dead half of helm/ & 42 & 655 & 8 \\
\texttt{eb313c08} & refactor(A3): one json\_compatible implementation (utils/jsonify.py) & 67 & 60 & 3 \\
\texttt{09449188} & refactor(C1): finish the build\_reports\_summary split (1,715 -> \textasciitilde{}330) & 1550 & 1397 & 6 \\
\texttt{364925b6} & refactor(B2): single named era\_mode flag in the bundle exporter & 18 & 5 & 2 \\
\texttt{648193d3} & refactor(D4): portfolio\_status logic -> reports/portfolio.py & 414 & 390 & 2 \\
\texttt{4d4085ef} & refactor(E3): local artifact writers in the filter-report emitters & 71 & 45 & 1 \\
\texttt{3d077a69} & refactor(C3): move the repair/publish half out of breakdown.py & 214 & 203 & 4 \\
\texttt{e630a576} & refactor(D3): layer the failure classifiers -- one generic taxonomy & 23 & 2 & 1 \\
\texttt{c42e8159} & docs(E5b): --skip-diagnosis kept -- phase3 4.8 closed as permanent & 11 & 0 & 2 \\
\texttt{c9d6c38d} & refactor(C2): extract the shared heatmap grid scaffold & 212 & 221 & 1 \\
\texttt{950ae99b} & docs(planning+journal): mark plan Batches 1-3 implemented & 81 & 2 & 2 \\
\texttt{36bb74e0} & refactor(D5): deprecate compare\_batch (phase 1, no behavior change) & 29 & 0 & 3 \\
\texttt{e1b8547d} & test(A4): extend phase3 baseline capture to the HELM render path (F1/F2) & 2952 & 18 & 5 \\
\texttt{4638f788} & docs(planning): record 2026-07-12 operator decisions + session journal & 106 & 0 & 2 \\
\texttt{75366f42} & chore(review): two lint fixes from the subagent-work review & 3 & 3 & 2 \\
\texttt{e295c13b} & fix(era): set tokenizer\_name + max\_sequence\_length on the era model deployment & 78 & 11 & 3 \\
\texttt{a3a33010} & fix(era-build): pin nltk==3.7 + provision punkt so BiasMetric tokenizes & 23 & 0 & 3 \\
\texttt{3651eeb1} & fix(runbooks): drop the deleted --allow-single-repeat flag from all callers & 9 & 13 & 8 \\
\texttt{3a9aa30b} & feat(presets): merge runbook-local presets via INFER\_STACK\_EXTRA\_PRESET\_FILES & 96 & 0 & 2 \\
\texttt{504c6b9f} & feat(reproduce): classic Together open-weight combined era-pinned runbook & 2546 & 0 & 17 \\
\texttt{b344671d} & fix(era-shim): render credentials.conf from the baked deployments key, not the env & 74 & 2 & 2 \\
\bottomrule\end{longtable}
\section*{2026-07-13 (16 commits)}
\begin{longtable}{@{}p{0.105\textwidth}p{0.665\textwidth}r r r@{}}
\toprule Hash & Subject & + & - & Files \\ \midrule\endhead
\texttt{47658cf7} & fix(olmo): pass a placeholder --base-url in the freeze-only preflight & 9 & 1 & 1 \\
\texttt{81c71002} & fix(gpt-oss,qwen): placeholder --base-url in the freeze-only preflights & 18 & 1 & 2 \\
\texttt{bb0a377b} & docs(era): note classic officials are carried forward v0.2.4<->v0.3.0 & 18 & 0 & 2 \\
\texttt{339748af} & docs(era): record original execution eras (v0.2.2 / v0.2.3), not the pinned eras & 13 & 4 & 2 \\
\texttt{5000b8ac} & docs(era): correct origin -- pre-v0.2.2 (original HELM), not v0.2.2 & 19 & 12 & 2 \\
\texttt{b445a7fd} & docs(planning): LiteLLM route-registry plan for the shared-gateway 400 bug & 578 & 0 & 2 \\
\texttt{d5c70c94} & fix(freeze): resolve subset run-entries to their exact dir, not AMBIGUOUS & 76 & 5 & 2 \\
\texttt{f64c20bd} & docs(planning): mark litellm-route-registry plan implemented & 6 & 2 & 1 \\
\texttt{1fbf4cbf} & docs(journal): route-registry implementation session & 70 & 0 & 1 \\
\texttt{9ff068c8} & feat(classic-together): run all targets in parallel, lease arbitrates GPUs & 168 & 167 & 4 \\
\texttt{27e4ee40} & docs(era): document classic-corpus class\_name migration + true origin window & 163 & 12 & 3 \\
\texttt{c44a6467} & chore(submodule): bump infer\_stack to route-registry merge (22d94317) & 1 & 1 & 1 \\
\texttt{c3e68bed} & feat(era-shim): canonicalize relocated class paths (G13 fix) & 214 & 3 & 3 \\
\texttt{3d904b92} & fix(gpt-oss): null-safe chat clients normalize content:null -> "" (HELM untouched) & 250 & 18 & 6 \\
\texttt{32657a62} & fix(docker): bump runner image Python 3.11 -> 3.12 to match eval\_audit floor & 9 & 5 & 3 \\
\texttt{c9f0efac} & fix(docker): prune selenium platform binaries / dangling symlinks before venv COPY & 13 & 0 & 1 \\
\bottomrule\end{longtable}
\section*{2026-07-14 (15 commits)}
\begin{longtable}{@{}p{0.105\textwidth}p{0.665\textwidth}r r r@{}}
\toprule Hash & Subject & + & - & Files \\ \midrule\endhead
\texttt{8d0d2c4b} & feat(reports): instance-level fallback for aggregate score-drift plot & 351 & 42 & 8 \\
\texttt{f34d9f96} & docs(journal): instance-level fallback for aggregate score-drift plot & 52 & 0 & 1 \\
\texttt{c041abb0} & fix(heatmap): headline metric resolves over metric family + main split & 178 & 10 & 2 \\
\texttt{b83c03e9} & docs(journal): headline metric family/main-split resolver fix & 43 & 0 & 1 \\
\texttt{33864be3} & Make eval audit run kwdagger with other\_session\_handler=kill & 1 & 0 & 1 \\
\texttt{ca885180} & fix(reports): align coverage-matrix axis order with the score-drift plot & 86 & 11 & 3 \\
\texttt{cdb31422} & feat(reproduce): add 50\_rsync\_from\_aiq\_gpu.sh to olmo/gpt-oss/qwen runbooks & 360 & 0 & 6 \\
\texttt{d4bf29d7} & feat(coverage-matrix): show match stats per column + align model order & 159 & 3 & 2 \\
\texttt{01351b73} & docs(journal): coverage-matrix annotations + axis alignment & 43 & 0 & 1 \\
\texttt{7a9f1d5f} & feat(coverage-matrix): print agreement proportion in each cell & 120 & 118 & 2 \\
\texttt{93849b09} & fix(heatmap): dedupe stale local attempts in aggregate score-drift & 175 & 1 & 2 \\
\texttt{bec1f57c} & docs(journal): aggregate score-drift stale-local dedupe & 57 & 0 & 1 \\
\texttt{e1d10217} & test(e2e): route-registry survival acceptance check & 201 & 0 & 2 \\
\texttt{a36b787c} & fix(heatmap): keep headline drift cells square for small model counts & 73 & 1 & 2 \\
\texttt{b8b85869} & fix(qwen): serve 8192 window on the instruct-turbo endpoints for from-spec ifeval/mmlu\_pro & 73 & 9 & 3 \\
\bottomrule\end{longtable}
\normalsize}{}
\IfFileExists{appendices/submodule_commit_ledger.tex}{\chapter{Submodule contribution ledger}\label{app:submodule-ledger}
\small
This appendix records Edward-authored commits in the included submodule histories during the internship window. Merge commits can have zero numstat even when they integrated substantial branch work.
\section*{\path{aiq-magnet} (3 commits)}
\begin{longtable}{@{}p{0.09\textwidth}p{0.10\textwidth}p{0.57\textwidth}r r@{}}
\toprule Date & Hash & Subject & + & - \\ \midrule\endhead
2026-05-29 & \texttt{762ca6b5} & Make run matching more robust & 129 & 13 \\
2026-06-25 & \texttt{4b10e1ba} & helm: add materialize\_helm\_run\_from\_spec (faithful run\_spec.json replay) & 931 & 0 \\
2026-06-25 & \texttt{50489f0a} & from-spec: optional --model-deployment rewrite (record the local endpoint) & 208 & 25 \\
\bottomrule\end{longtable}
\section*{\path{cmd_queue} (1 commits)}
\begin{longtable}{@{}p{0.09\textwidth}p{0.10\textwidth}p{0.57\textwidth}r r@{}}
\toprule Date & Hash & Subject & + & - \\ \midrule\endhead
2026-06-23 & \texttt{0cb1ac21} & Add first-class job setup/teardown lifecycle & 107 & 4 \\
\bottomrule\end{longtable}
\section*{\path{infer_stack} (30 commits)}
\begin{longtable}{@{}p{0.09\textwidth}p{0.10\textwidth}p{0.57\textwidth}r r@{}}
\toprule Date & Hash & Subject & + & - \\ \midrule\endhead
2026-05-20 & \texttt{328cd1c4} & Add plan for lora implementation & 151 & 0 \\
2026-05-20 & \texttt{6af064a0} & Merge pull request \#2 from AIQ-Kitware/lora\_branch & 151 & 0 \\
2026-05-27 & \texttt{538e4c5c} & Expose ports in the .env file & 32 & 5 \\
2026-06-23 & \texttt{8150e331} & Fix TUI resize bug and collapse serve and acquire + other stuff & 209 & 175 \\
2026-06-23 & \texttt{79902140} & Add admission queue to acquire (wait\_for\_placement / --queue) & 223 & 3 \\
2026-06-23 & \texttt{55c39a78} & Add `infer-stack gc` to reclaim leaked leases and free GPUs & 145 & 1 \\
2026-06-23 & \texttt{156f6bb7} & Render LiteLLM with a static superset route table (no-blip gateway) & 230 & 24 \\
2026-06-23 & \texttt{af4d1d6e} & Merge feature/litellm-no-blip into the leasing pipeline lifecycle & 230 & 24 \\
2026-06-23 & \texttt{bfb7dbe7} & Merge feature/leasing-and-no-blip into dev/leasing-controller & 597 & 27 \\
2026-06-23 & \texttt{cfddfac0} & Merge dev/leasing-controller: pipeline leasing lifecycle + no-blip gateway & 22700 & 12364 \\
2026-06-24 & \texttt{4381731e} & leasing: tolerate configs without a 'catalog' field in \_load\_catalog & 6 & 2 \\
2026-06-24 & \texttt{995f45e1} & Merge dev/lease-polish: leasing/e2e polish + compose-backend catalog fallback & 524 & 34 \\
2026-06-25 & \texttt{0cde11d5} & Merge dev/lease-polish: protocol-aware generation-gated readiness + --catalog routing for release/gc/evict & 226 & 49 \\
2026-06-25 & \texttt{933c3f2c} & gitignore: ignore setuptools build/ output dir & 1 & 0 \\
2026-06-29 & \texttt{045536cb} & Merge dev/lease-polish: host-port churn fix, reconcile locking, dynamic routing & 2708 & 162 \\
2026-06-30 & \texttt{e2f51fdb} & Merge remote-tracking branch 'origin/dev/lease-polish' into infer-stack-cli-api-migration & 375 & 71 \\
2026-06-30 & \texttt{104d4254} & leasing: release lease when a not-ready acquire times out & 178 & 10 \\
2026-06-30 & \texttt{beb1cc66} & fix(leasing): refuse to mutate without a real cross-process lock; group-write new lock files & 300 & 16 \\
2026-06-30 & \texttt{7fa342dd} & fix(tui): preserve scroll offset when a table rebuilds on poll & 88 & 5 \\
2026-07-01 & \texttt{caa67503} & fix(tui): require a double-click to acquire an endpoint & 180 & 16 \\
2026-07-07 & \texttt{6c586568} & feat(leasing): attention\_backend endpoint option (VLLM\_ATTENTION\_BACKEND) & 101 & 2 \\
2026-07-07 & \texttt{d685573c} & feat(tui): API tab respects endpoint completions vs chat protocol & 122 & 16 \\
2026-07-07 & \texttt{7f5fafed} & fix(tui): handle string-form reclaim spec in endpoint edit screen & 5 & 1 \\
2026-07-08 & \texttt{e5fba7b5} & feat(leasing): reserve-only GPU lease (acquire --reserve-gpus N) & 339 & 14 \\
2026-07-08 & \texttt{1109b0f0} & test(leasing): e2e probe for the reserved-GPU index-frame assumption & 150 & 0 \\
2026-07-08 & \texttt{c119a62e} & docs(changelog): note the reserve-only GPU lease (--reserve-gpus) & 17 & 0 \\
2026-07-13 & \texttt{2d832b61} & feat(leasing): append-only LiteLLM route registry (static-superset) & 783 & 37 \\
2026-07-13 & \texttt{869080bc} & feat(cli): `infer-stack routes` group + route-registry docs & 522 & 18 \\
2026-07-13 & \texttt{c005a887} & fix(tui): surface HTTP error bodies in the API tab & 63 & 2 \\
2026-07-13 & \texttt{22d94317} & Merge feat/litellm-route-registry into infer-stack-cli-api-migration & 1368 & 57 \\
\bottomrule\end{longtable}
\section*{\path{kwdagger} (1 commits)}
\begin{longtable}{@{}p{0.09\textwidth}p{0.10\textwidth}p{0.57\textwidth}r r@{}}
\toprule Date & Hash & Subject & + & - \\ \midrule\endhead
2026-06-23 & \texttt{ce246bae} & Add ProcessNode setup/teardown resource lifecycle & 150 & 0 \\
\bottomrule\end{longtable}
\normalsize}{}

\end{document}
